\documentclass{article} 
\usepackage{amsmath} 
\usepackage{amssymb} 
\usepackage[utf8]{inputenc}
\usepackage{geometry} \geometry{a4paper, margin=1in}
\begin{document} 
\title{\textbf{Recovered Question Bank}}
\date{\today}
\maketitle
\section{...}
\subsection{Easy}
\begin{enumerate}
\item Q. The area of the triangle whose vertices are represented by the complex number 0, z, $ze^{i\alpha}$, ($0 < \alpha < \pi$) equals \hfill \textbf{(B)}
\item Q. $\int \frac{\operatorname{cosec} x}{\cos^2 \left(1+\log \tan \frac{x}{2}\right)}dx$ is equal to \hfill \textbf{(B)}
\end{enumerate}
\subsection{Medium}
\begin{enumerate}
\item Q. If $|z| = \max\{|z - 2|, |z + 2|\}$, then locus of z is \hfill \textbf{(D)}
\item Q. Let f be a positive function. Let $I_1 = \int_{1-k}^k x f \{x(1-x)\}dx$, $I_2 = \int_{1-k}^k f\{x(1-x)\}dx$, where $2k - 1 > 0$. Then, $\frac{I_1}{I_2}$ is \hfill \textbf{(D)}
\item Q. Equation of the line of the shortest distance between the lines $\frac{x}{2} = \frac{y}{-3} = \frac{z}{1}$ & $\frac{x-2}{3} = \frac{y-1}{-5} = \frac{z+2}{2}$ is \hfill \textbf{(A)}
\item Q. If $\sin^{-1}\left(\sin \frac{2x^2+4}{1+x^2}\right) < \pi-3$, then x belongs to \hfill \textbf{(B)}
\item Q. ABC is a triangle where A = (2,3,5), B = (-1,3,2) and C = ($\lambda$, 5, $\mu$). If the median through A is equally inclined to the axes, then the values of $\lambda, \mu$ and the area of the $\triangle ABC$ is \hfill \textbf{(A)}
\item Q. The total number of integers between 1&3600 (both excluded), that are relatively prime to 3600 \hfill \textbf{(A)}
\item Q. The value of $\lim_{n\to\infty} \left[\frac{1}{n^2} \sec^2 \frac{1}{n^2} + \frac{2}{n^2} \sec^2 \frac{4}{n^2} + \dots + \frac{1}{n} \sec^2 1\right]$ is equal to \hfill \textbf{(C)}
\item Q. In a certain test, there are n questions. In the test $2^{n-1}$ students gave wrong answers to at least i questions, where i = 1, 2, . . . .n. If the total number of wrong answers given is 2047, then n is equal to \hfill \textbf{(B)}
\end{enumerate}
\subsection{Hard}
\begin{enumerate}
\item Q. If $C_1 : x^2 + y^2 = (3+2\sqrt{2})^2$ be a circle and PA and PB are pair of tangents on $C_1$, where P is any point on the director circle of $C_1$, then the radius of smallest circle which touch $C_1$ externally and also the two tangents PA and PB is \hfill \textbf{(D)}
\item Q. The coefficient of $x^{1274}$ in the expansion of $(x+1)(x-2)^2(x+3)^3(x-4)^4 \dots (x+49)^{49}(x-50)^{50}$ is \hfill \textbf{(B)}
\item Q. If $f'(x) > 0$ and $f\prime(1) = 0$ such that $g(x) = f (\cot^2 x + 2 \cot x + 2)$, where $0 < x < \pi$, then the interval in which $g(x)$ is decreasing is \hfill \textbf{(D)}
\item Q. If the line $y = \sqrt{3}x$ cuts the curve $x^3 + y^3 + 3xy + 5x^2 + 3y^2 + 4x + 5y \ndash 1 = 0$ at the points A, B, C and O be the origin, then product OA$\cdot$ОB$\cdot$OC is \hfill \textbf{(A)}
\item Q. Consider all the six-digit numbers that can be formed using the digits 1, 2, 3, 4, 5 and 6, each digit being used exactly once. Each of such six-digit numbers have the property that for each digit, not more than two digits, smaller than that digit, appear to the right of that digit. Find the number of such six-digit numbers having the desired property. \hfill \textbf{(D)}
\end{enumerate}
\section{3D Geometry}
\subsection{Medium}
\begin{enumerate}
\item Q. The plane $x - 2y + 3z = 17$ divides the line segment joining $P_1(-2, 4, 7)$ and $P_2(3, -5, 8)$ in the ratio $k:1$. Find $k:1$.
(A) 3:10
(B) 10:3
(C) 7:3
(D) 3:7 \hfill \textbf{((A))}
\item Q. Find the coordinates of the points on the line $\frac{x-1}{2} = \frac{y+2}{3} = \frac{z-3}{6}$ that are at a distance of 3 units from the point $(1, -2, 3)$.
(A) $(\frac{13}{7}, -\frac{5}{7}, \frac{39}{7})$ and $(\frac{1}{7}, -\frac{23}{7}, \frac{3}{7})$
(B) $(4,1,9)$ and $(-2, -5, -3)$
(C) $(\frac{13}{7}, -\frac{5}{7}, \frac{3}{7})$ and $(\frac{1}{7}, -\frac{23}{7}, -\frac{39}{7})$
(D) $(-\frac{11}{7}, \frac{11}{7}, \frac{15}{7})$ and $(\frac{15}{7}, -\frac{21}{7}, -\frac{31}{7})$ \hfill \textbf{((A))}
\item Q. The foot of the perpendicular from the point A(1,1,1) onto the line joining B(5, 4, 4) and C(1, 4, 6) has coordinates
(A) (3, 4, 5)
(B) (2, 3, 4)
(C) (5, 4, 4)
(D) (1, 4, 6) \hfill \textbf{((A))}
\item Q. Find the coordinates of the image P' of the point P(2, −3, 1) in the line $\frac{x+1}{2} = \frac{y-3}{3} = \frac{z+2}{-1}$. That is, reflect P across this line. The options are:
(A) $(-\frac{58}{7}, \frac{18}{7}, -\frac{20}{7})$
(B) $(-\frac{44}{7}, -\frac{6}{7}, -\frac{20}{7})$
(C) $(-\frac{58}{7}, -\frac{18}{7}, \frac{20}{7})$
(D) $(\frac{58}{7}, \frac{18}{7}, -\frac{20}{7})$ \hfill \textbf{((A))}
\item Q. Consider the two lines $L_1 : \vec{r} = (1, -1, -10) + \lambda (2, -3, 8)$ and $L_2 : \vec{r} = (4, -3, -1) + \mu (1, -4, 7)$. Which of the following statements is true?
(A) $L_1$ and $L_2$ are coplanar and intersecting.
(B) $L_1$ and $L_2$ are coplanar but parallel (no intersection).
(C) $L_1$ and $L_2$ are skew (non-coplanar).
(D) $L_1$ and $L_2$ are coincident (identical). \hfill \textbf{((A))}
\item Q. The length of the perpendicular from the point $(3, -1, 11)$ to the line $\frac{x}{2} = \frac{y-2}{3} = \frac{z-3}{4}$ is:
(A) $\sqrt{53}$
(B) $\sqrt{33}$
(C) $\sqrt{29}$
(D) $\sqrt{66}$ \hfill \textbf{((A))}
\item Q. A plane passes through $(1, 2, 2)$ and is perpendicular to two planes $2x - 2y + z = 0$ and $x - y + 2z = 4$. Square of the distance of the plane from the point $(52, 53, 57)$ is
(A) 5202
(B) 5201
(C) 5203
(D) 5000 \hfill \textbf{((A))}
\end{enumerate}
\subsection{Hard}
\begin{enumerate}
\item Q. If the lines whose direction cosines $(l, m, n)$ satisfy $al + bm + cn =0$, $fmn+gnl+hlm=0$ are perpendicular, then the value of $\frac{f}{a} + \frac{g}{b} + \frac{h}{c}$ is
(A) 0
(B) -1
(C) 1
(D) none of these \hfill \textbf{((A))}
\item Q. A line makes angles $\alpha, \beta, \gamma, \delta$ with the four body-diagonals of a cube. Then find the correct value of $\cos^2 \alpha + \cos^2 \beta + \cos^2 \gamma + \cos^2 \delta$.
(A) 3
(B) 4
(C) $\frac{4}{3}$
(D) $\frac{3}{4}$ \hfill \textbf{((C))}
\end{enumerate}
\section{Algebra}
\subsection{Medium}
\begin{enumerate}
\item Q. The number of integral solutions of the equation
$2x - 1 = |x + 7|$
is
(A) 0
(B) 1
(C) 2
(D) infinitely many \hfill \textbf{(B)}
\item Q. The total number of pairs of consecutive odd natural numbers, both greater than 10,
such that their sum is less than 40, is
(A) 3
(B) 4
(C) 5
(D) 6 \hfill \textbf{(B)}
\item Q. If
$(\sqrt{2})^x + (\sqrt{3})^x = (\sqrt{13})^{x/2}$,
then the number of real values of x is
(A) 0
(B) 1
(C) 2
(D) infinitely many \hfill \textbf{(B)}
\item Q. Solve for x:
$\log_7(\log_5(\sqrt{x + 5} + \sqrt{x})) = 0$.
The number of value of x is
(A) 0
(B) 1
(C) 2
(D) 3 \hfill \textbf{(B)}
\item Q. Solve for x:
$\frac{2x}{2x^2 + 5x + 2} > \frac{1}{x + 1}$
(A) $x \in (-2, -1) \cup (-\frac{2}{3}, -\frac{1}{2})$
(B) $x \in (-2, -\frac{2}{3}) \cup (-1, -\frac{1}{2})$
(C) $x \in (-2, -1) \cup (\frac{1}{2}, \infty)$
(D) $x \in (-\infty, -2) \cup (-1, -\frac{2}{3})$ \hfill \textbf{(A)}
\item Q. The number of real solutions of the equation
$4^x - 3^{x - \frac{1}{2}} = 3^{x + \frac{1}{2}} - 2^{2x - 1}$,
$x \in \mathbb{R}$,
is
(A) 0
(B) 1
(C) 2
(D) infinitely many \hfill \textbf{(B)}
\item Q. The total number of integral values of x satisfying
$x^{\log_3 (x^2) + (\log_3 x)^2 - 10} = \frac{1}{x^2}$
is
(A) 0
(B) 1
(C) 2
(D) infinitely many \hfill \textbf{(C)}
\item Q. The number of integer solutions of the equation
$|x| + |\frac{4 - x^2}{x}| = |x + \frac{4 - x^2}{x}|$
is
(A) 0
(B) 2
(C) 4
(D) infinitely many \hfill \textbf{(C)}
\item Q. The number of integral values of x (with $x \neq 0$) that do not satisfy
$|x + \frac{1}{x}| > 2$
is
(A) 0
(B) 1
(C) 2
(D) infinitely many \hfill \textbf{(C)}
\item Q. If a, b, c are distinct positive real numbers satisfying
$a^2 + b^2 + c^2 = 1$,
then ab + bc + ca is necessarily less than
(A) $\frac{1}{2}$
(B) $\frac{2}{3}$
(C) 1
(D) none of these \hfill \textbf{(C)}
\item Q. If
$5\{x\} = x + [x]$, $[x] - \{x\} = \frac{1}{2}$,
where $[x]$ and $\{x\}$ denote the integer and fractional parts of x, then x is
(A) $\frac{1}{2}$
(B) $\frac{3}{2}$
(C) $\frac{5}{2}$
(D) 2 \hfill \textbf{(B)}
\item Q. The number of integral solutions of the inequality
$x + \sqrt{3 - x} \geq \sqrt{3 - x} + 3$
is
(A) 0
(B) 1
(C) 2
(D) infinitely many \hfill \textbf{(B)}
\item Q. Solve for x:
$\frac{(x - 1)(x + 5)}{x(x + 2)} \geq 0$.
(A) $x \in (-\infty, -5] \cup [-2, 0] \cup [1, \infty)$
(B) $x \in (-\infty, -5] \cup (-2, 0) \cup [1, \infty)$
(C) $x \in (-\infty, -5] \cup [1, \infty)$
(D) $x \in (-\infty, -5] \cup (-2, 0) \cup (1, \infty)$ \hfill \textbf{(B)}
\item Q. Solve for x:
$(x + 2) (x - 1)^2 (x - 5) \geq 0$.
(A) $x \in (-\infty, -2] \cup (5, \infty)$
(B) $x \in (-\infty, -2] \cup (5, \infty) \cup \{2\}$
(C) $x \in [-2, -5]$
(D) none of these \hfill \textbf{(D)}
\item Q. Solve for x:
$\frac{(x - 2) x}{(x - 1)^2 (x^2 - 5x + 6)} \geq 0$.
(A) $x \in (-\infty, 0] \cup (3, \infty)$
(B) $x \in (-\infty, 0) \cup (3, \infty)$
(C) $x \in (-\infty, 0] \cup [3, \infty)$
(D) $x \in (-\infty, 0] \cup (2, \infty)$ \hfill \textbf{(A)}
\item Q. Let
$f(x) = x^6 - 2x^5 + x^3 + x^2 - x - 1$, $g(x) = x^4 - x^3 - x^2 - 1$,
and let $a, b, c, d$ be the four roots of $g(x) = 0$. What is
$\sum f(a) + f(b) + f(c) + f(d)$ ? \hfill \textbf{((B))}
\item Q. Let $a_1, a_2, \dots, a_n$ be real numbers such that $\sqrt{a_1} + \sqrt{a_2-1} + \sqrt{a_3 - 2} + \dots + \sqrt{a_n - (n - 1)} = \frac{1}{2}(a_1 + a_2 + \dots + a_n) - \frac{n(n-3)}{4}$. Compute the value of $\sum_{i=1}^{100} a_i$. \hfill \textbf{((D))}
\item Q. The sum of values of $x$ satisfying the equation $(31 + 8\sqrt{15})^{x^2-3} + 1 = (32 + 8\sqrt{15})^{x^2-3}$ is
(A) 3
(B) 0
(C) 2
(D) none of these \hfill \textbf{((B))}
\item Q. If $\alpha, \beta$ are the roots of $x^2 + px + q = 0$ and $x^{2n} + p^n x^n + q^n = 0$ and if $\left(\frac{\alpha}{\beta}\right), \left(\frac{\beta}{\alpha}\right)$ are the roots of $x^n + 1 + (x + 1)^n = 0$, then $n(\in \mathbb{N})$
(A) must be an odd integer.
(B) may be any integer.
(C) must be an even integer.
(D) can not say anything. \hfill \textbf{((C))}
\item Q. Number of values of $x$ for which $\frac{8^x+27^x}{12^x+18^x} = \frac{7}{6}$
(A) 2
(B) 3
(C) 1
(D) no value of x \hfill \textbf{((A))}
\item Q. If $\alpha, \beta, \gamma$ are such that $\alpha+ \beta + \gamma = 2$, $\alpha^2 + \beta^2 + \gamma^2 = 6$, $\alpha^3 + \beta^3 + \gamma^3 = 8$, then $\alpha^4 + \beta^4 + \gamma^4$ is
(A) 18
(B) 15
(C) 10
(D) 36 \hfill \textbf{((A))}
\item Number of ordered pair(s) of $(x, y)$ satisfying the system of simultaneous equations $|x^2 – 2x| + y = 1$ and $x^2 + |y| = 1$ is ($x, y \in R$)\\
(A) 1\\
(B) 2\\
(C) 4\\
(D) infinitely many \hfill \textbf{((C))}
\item If $c < 1$ and the system of equations $x + y – 1 = 0, 2x - y - c = 0$ and $-bx + 3by – c = 0$ is consistent, then the possible real values of $b$ are\\
(A) $b\in (-3,\frac{1}{4})$\\
(B) $b\in (-3,4)$\\
(C) $b\in (-\frac{3}{4},3)$\\
(D) none of these \hfill \textbf{((C))}
\item Q. If $x^2 + y^2 = t + \frac{1}{t}$, $x^4 + y^4 = t^2 + \frac{1}{t^2}$, then - \hfill \textbf{((D))}
\item Q. Let $f(x) = x^3 + x^2 + 100x + 7 \sin x$, then equation $\frac{1}{y-f(1)} + \frac{2}{y-f(2)} + \frac{3}{y-f(3)} = 0$ has \hfill \textbf{((C))}
\item Q. If a, b, c > 0, $a^2 = bc$ and $a + b + c = abc$, then the least value of $a^4 + a^2 + 7$ must be equal to
(A) 16
(B) 18
(C) 21
(D) 19 \hfill \textbf{((D))}
\item Q. Let S be a set of n real numbers with mean m. It is known that mean of S \cup \{15\} = m + 2, mean of S\cup \{15, 1\} = m + 1. How many elements does S have? \hfill \textbf{((A))}
\item Q. Let p(x) be a polynomial satisfying
p(x) – p'(x) = x^n,
where n is a positive integer. Then p(0) equals
(A) n!
(B) (n - 1)!
(C) \frac{1}{n}
(D) \frac{1}{(n-1)!} \hfill \textbf{((A))}
\end{enumerate}
\subsection{Hard}
\begin{enumerate}
\item Q. If $\alpha, \beta, \gamma$ are the roots of the equation $x^3 - 7x + 7 = 0$, then $\frac{1}{\alpha^4} + \frac{1}{\beta^4} + \frac{1}{\gamma^4}$ is
(A) $\frac{7}{3}$
(B) $\frac{3}{7}$
(C) $\frac{4}{7}$
(D) $\frac{7}{4}$ \hfill \textbf{((B))}
\end{enumerate}
\section{Algebra (Progressions)}
\subsection{Medium}
\begin{enumerate}
\item Q. If $a, b, c$ are three real numbers which are in A.P. and $a^2, b^2, c^2$ are in G.P.. If $a < b < c$ and $a + b + c = \frac{3}{2}$, then the value of $a$ is
(A) $\frac{1}{2\sqrt{2}}$
(B) $\frac{1}{2\sqrt{3}}$
(C) $\frac{1}{2} - \frac{1}{\sqrt{3}}$
(D) $\frac{1}{2} - \frac{1}{\sqrt{2}}$ \hfill \textbf{((D))}
\end{enumerate}
\section{Algebra (Quadratic Equations, Progressions)}
\subsection{Medium}
\begin{enumerate}
\item Q. If $a, b, c, d$ are four consecutive terms of an increasing A. P., then the roots of the equation $(x - a)(x - c) + 2(x - b)(x - d) = 0$ are
(A) on-real complex
(B) real and equal
(C) integers
(D) real and distinct \hfill \textbf{((D))}
\item Q. If positive numbers $a^{-1}, b^{-1}$ and $c^{-1}$ are in A. P., then the product of roots of the equation $x^2 - kx + 2b^{101} - a^{101} - c^{101} = 0$, ($k \in R$)
(A) $> 0$
(B) $= 0$
(C) $< 0$
(D) Can't be determined \hfill \textbf{((C))}
\end{enumerate}
\section{Algebraic Expressions And Logarithms}
\subsection{Medium}
\begin{enumerate}
\item The expression $\frac{\frac{x^2 + 3x + 2}{x+2} + 3x - \frac{x(x^3 + 1)}{(x + 1)(x^2 - x + 1)}}{(x - 1) (\log_2 3) (\log_3 4) (\log_4 5) (\log_5 2)} \log_2 8$ reduces to \hfill \textbf{(A)}
\end{enumerate}
\section{Arithmetic Progressions}
\subsection{Medium}
\begin{enumerate}
\item Q. The sums of the first n terms of two arithmetic progressions are in the ratio $\frac{S_{1,n}}{S_{2,n}} = \frac{7n+1}{4n+27}$. Find the ratio of their 11th terms $t_{1,11} : t_{2,11}$.
(A) $\frac{4}{3}$
(B) $\frac{148}{111}$
(C) $\frac{111}{148}$
(D) $\frac{7}{4}$ \hfill \textbf{(A)}
\item Q. Four distinct integers form an increasing arithmetic progression. One of these four numbers equals the sum of the squares of the other three. The four numbers are
(A) -2, -1, 0, 1
(B) -1, 0, 1, 2
(C) 0, 1, 2, 3
(D) -2, 0, 2, 4 \hfill \textbf{(B)}
\item Q. If $a_1, a_2, ..., a_n$ are in arithmetic progression with $a_i > 0$ for all $i$, then the sum
$\sum_{k=1}^{n-1} \frac{1}{\sqrt{a_k} + \sqrt{a_{k+1}}}$ equals
(A) $\frac{n-1}{\sqrt{a_1}+\sqrt{a_n}}$
(B) $\frac{n}{\sqrt{a_1}+\sqrt{a_n}}$
(C) $\frac{n-1}{\sqrt{a_n} - \sqrt{a_1}}$
(D) $\frac{n+1}{\sqrt{a_1} + \sqrt{a_n}}$ \hfill \textbf{(A)}
\end{enumerate}
\section{Binomial Theorem}
\subsection{Medium}
\begin{enumerate}
\item Q. Let $\binom{n}{k}$ be the binomial coefficients. Evaluate the sum $2\binom{n}{1} + 2^3\binom{n}{3} + 2^5\binom{n}{5} + \dots$.
(A) $\frac{3^n - (-1)^n}{2}$
(B) $\frac{3^n + (-1)^n}{2}$
(C) $\frac{3^n + 1}{2}$
(D) $\frac{3^n - 1}{2}$ \hfill \textbf{((A))}
\item Q. If the magnitude of the coefficient of $x^7$ in the expansion of $(ax^2 + \frac{1}{bx})^8$, where a and b
are positive numbers, is equal to the magnitude of the coefficient of $x^{-7}$ in the expansion of
$(ax - \frac{1}{bx^2})^8$, then a and b are connected by the relation
(A) $ab = 1$
(B) $ab = 2$
(C) $a^2b = 1$
(D) $ab^2 = 2$ \hfill \textbf{((A))}
\item What is the the value of $\frac{1}{15} {}^{14}C_1+\frac{5}{2} {}^{14}C_1+\frac{5^2}{3} {}^{14}C_2+\frac{5^3}{4} {}^{14}C_3+\dots.+\frac{5^{14}}{15} {}^{14}C_{14}$ : (where $"C_r$ is the coefficient of $x^r$ in $(1 + x)^n$)\\
(A) $\frac{5^{16}-1}{16}$\\
(B) $\frac{5^{15}-1}{15}$\\
(C) $\frac{6^{15}-1}{15}$\\
(D) $\frac{6^{15}}{15}$ \hfill \textbf{((C))}
\item The first 3 terms in the expansion of $(1 + ax)^n(n \neq 0)$ are $1,6x$, and $16x^2$. Then, the value of $a$ and $n$ are respectively\\
(A) 2 and 9\\
(B) 3 and 2\\
(C) $\frac{2}{3}$ and 9\\
(D) $\frac{2}{3}$ and 6 \hfill \textbf{((C))}
\item Q. If the coefficient of r th, (r + 1) th and (r + 2) th terms in the binomial expansion of $(1 + y)^m$ are in AP, then m and r satisfy the equation
(A) $m^2 - m(4r - 1) + 4r^2 + 2 =0$
(B) $m^2 - m(4r + 1) + 4r^2 - 2 =0$
(C) $m^2 - m(4r + 1) + 4r^2 + 2 = 0$
(D) $m^2 - n(4r - 1) + 4r^2 - 2 = 0$ \hfill \textbf{((B))}
\item Q. Let m, n be natural numbers such that $(1 - y)^m (1 + y)^n = 1 + a_1y + a_2y^2 + \cdots$ and it is given that $a_1 = a_2 = 10$. Determine (m, n). \hfill \textbf{((D))}
\end{enumerate}
\section{Calculus}
\subsection{Medium}
\begin{enumerate}
\item Q. $\lim_{x\to 2} \frac{2^x+2^{3-x}-6}{2^{-x/2}-2^{1-x}}$ is equal to \hfill \textbf{((A))}
\item Q. Evaluate $\lim_{x\to 1} \frac{x^{P+1}-(P+1)x+P}{(x-1)^2}$ where P is a constant. The value of the limit is \hfill \textbf{((A))}
\item Q. Evaluate $\lim_{x\to 1} \frac{\sqrt{x^2+8}-\sqrt{10-x^2}}{\sqrt{x^2+3}-\sqrt{5-x^2}}$. This is of the indeterminate form $\frac{0}{0}$. The value of the limit is \hfill \textbf{((C))}
\item Q. If $\lim_{x\to\infty} (\frac{x^3+1}{x^2+1}-(ax+b)) = 2$, then \hfill \textbf{((C))}
\item Q. Evaluate the limit $\lim_{n\to\infty} \frac{[x]+[2x]+[3x]+\cdots+[nx]}{n^2}$ where [.] denotes the greatest-integer function. The value of the limit is \hfill \textbf{((A))}
\item Q. Evaluate $\lim_{n\to\infty} \frac{\sin(\frac{a}{n})}{\tan(\frac{b}{n+1})}$ where a, b are constants. The value of the limit is \hfill \textbf{((A))}
\item Q. Evaluate $\lim_{x\to\frac{\pi}{2}} (\sec x - \tan x)$. This limit is of the $\infty-\infty$ form. The value of the limit is \hfill \textbf{((A))}
\item Q. Evaluate $\lim_{x\to 0} (\frac{a^x+b^x+c^x}{3})^{1/x}$, where a, b, c > 0. The value of this limit is \hfill \textbf{((C))}
\item Q. Let $f(x) = \begin{cases} \frac{1}{e^{1/x}+1}, & x \ne 0, \\ 0, & x = 0. \end{cases}$ Which of the following is true? \hfill \textbf{((C))}
\item Q. Let $f(x) = \begin{cases} a \sin^{2n} x, & x \ge 0, n \to \infty, \\ b \cos^{2m} x - 1, & x < 0, m \to \infty. \end{cases}$ Then which statement is true? \hfill \textbf{((A))}
\item Q. f is a continuous function on R. Given that $x^2 + (f(x) – 2) x - \sqrt{3} f(x) + 2\sqrt{3}-3 = 0 \forall x$, then the value of f($\sqrt{3}$) is \hfill \textbf{((B))}
\item Q. Let
$J = \int_{0}^{1} \frac{x}{1 + x^8} dx$.
Consider the assertions:
I. $J > \frac{1}{4}$, II. $J < \frac{\pi}{8}$.
Which of the following is true? \hfill \textbf{((A))}
\item Q. Let $f : (-1, 1) \to \mathbb{R}$ be a differentiable function satisfying
$(f'(x))^4 = 16(f(x))^2$ for all $x \in (-1,1)$, $f(0) = 0$.
How many such functions $f$ exist? \hfill \textbf{((D))}
\item Q. For $x \in \mathbb{R}$, define
$f(x) = |\sin x|$, $g(x) = \int_{0}^{x} f(t) dt$, $p(x) = g(x) - \frac{2}{\pi} x$.
Which of the following is true? \hfill \textbf{((A))}
\item Q. Let $f : [0, 1] \to [0, 1]$ be a continuous function such that
$x^2 + (f(x))^2 \le 1$ for all $x \in [0, 1]$,
and
$\int_{0}^{1} f(x)dx = \frac{\pi}{4}$.
Then
$\int_{\frac{1}{2}}^{\frac{\sqrt{2}}{2}} \frac{f(x)}{\sqrt{1-x^2}} dx$
equals \hfill \textbf{((A))}
\item Q. Define
$f(x) = \begin{cases} \frac{\sin(x^2)}{x}, & x \neq 0, \\ 0, & x = 0. \end{cases}$
At $x = 0$, $f$ is \hfill \textbf{((D))}
\item Q. Let $f: \mathbb{R} \rightarrow \mathbb{R}$ be defined by $f(x) = x^3 + ax^2 + bx + c$, where $a, b, c \in \mathbb{R}$. For which conditions on $a, b, c$ is $f$ bijective?
(A) $b^2 \geq 3a, c \in \mathbb{R}$
(B) $a^2 \leq 3b, c \in \mathbb{R}$
(C) $b^2 \geq 3a, c \in \mathbb{R}$
(D) $a^2 \geq 3b, c \in \mathbb{R}$ \hfill \textbf{((B))}
\item Q. Evaluate the limit $L = \lim_{x\to\frac{\pi}{2}} \left( \frac{1 - \tan(\frac{x}{2})}{1 + \tan(\frac{x}{2})} \right) \frac{1 - \sin x}{(\pi - 2x)^3}$
(A) $\frac{1}{8}$
(B) 0
(C) $\frac{1}{32}$
(D) does not exist \hfill \textbf{((C))}
\item Q. Let $f: \mathbb{R}\to \mathbb{R}$ be defined by $f(x) = \frac{1}{3}x^3 - 2x^2 + mx + 1$, where $m \in \mathbb{R}$. For which $m$ is $f$ strictly monotonic on $\mathbb{R}$?
(A) $(4,\infty)$
(B) $[4,\infty)$
(C) $(-\infty,4]$
(D) $(-\infty,4)$ \hfill \textbf{((B))}
\item Q. Evaluate $\int_0^{\frac{\pi}{2}} \frac{2^{\sin x}}{2^{\sin x} + 2^{\cos x}} dx$.
(A) $2\pi$
(B) $\frac{\pi}{4}$
(C) $\pi$
(D) $\frac{\pi}{2}$ \hfill \textbf{((B))}
\item Q. If $I_1 = \int_0^{\pi} \frac{x \sin x}{1+\cos^2 x} dx$, $I_2 = \int_0^{\pi} x \sin^4 x dx$, then $I_1 : I_2$ is equal to
(A) 3:4
(B) 1:2
(C) 4:3
(D) 2:3 \hfill \textbf{((C))}
\item Q. If $\int \frac{3 \sin x + 2 \cos x}{3 \cos x + 2 \sin x} dx = ax + b \ln |2 \sin x + 3 \cos x| + C$, then
(A) $a = \frac{12}{13}, b = \frac{15}{30}$
(B) $a = -\frac{7}{13}, b = \frac{6}{13}$
(C) $a = \frac{12}{13}, b = -\frac{15}{39}$
(D) $a = -\frac{7}{13}, b = -\frac{6}{13}$ \hfill \textbf{((C))}
\item Q. Water runs into an inverted conical tent at the rate of $20\text{ft}^3/\text{min}$ and leaks out at the rate of $5\text{ft}^3/\text{min}$. The height of the water is three times the radius of the water's surface. The rate at which the radius of the water surface is increasing when the radius is 5 ft, is
(A) $\frac{1}{5\pi} \text{ft/min}$
(B) $\frac{1}{10\pi} \text{ft/min}$
(C) $\frac{15}{5\pi} \text{ft/min}$
(D) None of these \hfill \textbf{((A))}
\item Q. If the function $f$ defined as $f(x) = \frac{1}{x} - \frac{k}{e^{2x}-1}$, $x \neq 0$ is continuous at $x = 0$, then ordered pair $(k, f(0))$ is equal to
(A) (2,1)
(B) (3,1)
(C) (3,2)
(D) $\left(\frac{1}{2},2\right)$ \hfill \textbf{((B))}
\item Q. The value of the integral $e^{x}\left[\frac{2\tan x}{1+\tan x}+\cot^{2}\left(x+\frac{\pi}{4}\right)\right]dx$ is equal to
(A) $e^{x}\tan\left(\frac{\pi}{4}-x\right)+C$
(B) $e^{x}\tan\left(x-\frac{\pi}{4}\right)+C$
(C) $e^{x}\tan\left(\frac{3\pi}{4}-x\right)+C$
(D) None of these \hfill \textbf{((B))}
\item Q. min $(x_1-x_2)^2 + (3 + \sqrt{1 - x_1^2} - \sqrt{4x_2})^2$, $\forall x_1, x_2 \in R$ is
(A) $4\sqrt{5}+1$
(B) $3 - 2\sqrt{2}$
(C) $\sqrt{5}+1$
(D) $\sqrt{5}-1$ \hfill \textbf{((B))}
\item Q. If $f(x) = x(\sqrt{x} + \sqrt{x + 1})$, then
(A) $f(x)$ is continuous but not differentiable at $x = 0$
(B) $f(x)$ is differentiable at $x = 0$
(C) $f(x)$ is not differentiable at $x = 0$
(D) None of the above \hfill \textbf{((C))}
\item Q. The difference between the greatest and least values of the function $f(x) = -x + \sin 2x$
on $[-\frac{\pi}{2}, \frac{\pi}{2}]$ is
(A) $\sqrt{3}+\sqrt{2}$
(B) $\sqrt{3}+\sqrt{2}+\frac{\pi}{6}$
(C) $\frac{\pi}{3}+\frac{\sqrt{3}}{2}$
(D) $\frac{\pi}{3}-\frac{\sqrt{3}}{2}$ \hfill \textbf{((C))}
\item $\int \frac{x}{(x^2-a^2)(x^2-b^2)} dx$ is equal to\\
(A) $\frac{1}{2(a^2-b^2)} \log\left|\frac{x^2-a^2}{x^2-b^2}\right|+c$\\
(B) $\frac{1}{2(b^2-a^2)} \log\left|\frac{x^2-a^2}{x^2-b^2}\right|+c$\\
(C) $\frac{1}{2(a^2-b^2)} \log\left|\frac{x^2-b^2}{x^2-a^2}\right|+c$\\
(D) $\frac{1}{2(a^2-b^2)} \log\left|\frac{x^2-b^2}{x^2-a^2}\right|+c$ \hfill \textbf{((C))}
\item The function $f(x) = 1 + x(\sin x) [\cos x]$, $0 < x \le \pi/2$, where $[\cdot]$ represents the greatest integer function,\\
(A) is continuous on $(0, \pi/2)$\\
(B) is strictly decreasing in $(0, \pi/2)$\\
(C) is strictly increasing in $(0, \pi/2)$\\
(D) has global maximum value 2 \hfill \textbf{((A))}
\item $\int_{-\log 3}^{\log 3} \cot^{-1} (\frac{e^x-1}{e^x+1}) dx =$\\
(A) $\frac{\pi}{2} \log 3$\\
(B) $\pi\log 3$\\
(C) $0$\\
(D) $\pi\log 9$ \hfill \textbf{((B))}
\item Q. If $\int_{0}^{f(x)} t^2dt = x \cos \pi x$, then $f'(9)$ is \hfill \textbf{((A))}
\item Q. If $f(x) = p|\sin x| + qe^{|x|} + r|x|^3$ and if $f(x)$ is differentiable at $x = 0$, then \hfill \textbf{((B))}
\item Q. If $I = \int_{-\pi/6}^{\pi/6} \frac{\pi+4x^5}{1-\sin(|x|+\frac{\pi}{6})} dx$, then $I$ is equal to \hfill \textbf{((D))}
\item Q. The function $f(x) = \lim_{n\to\infty}(\sin x)^n$ is discontinuous at
(A) $n\pi + (-1)^n\frac{\pi}{2}, n \in I$
(B) $n\pi + \frac{\pi}{2}, n \in I$
(C) $2n\pi\pm\frac{\pi}{2}, n \in I$
(D) None of these \hfill \textbf{((B))}
\item Q. At present, a firm is manufacturing 2000 items. It is estimated that the rate of change of production P w.r.t. additional number of workers x is given by $\frac{dP}{dx} = 100 - 12\sqrt{x}$. If the firm employs 25 more workers, then the new level of production of items is
(A) 3500
(B) 4500
(C) 2500
(D) 3000 \hfill \textbf{((A))}
\item Q. Let a > 0 be a real number. Then the limit
L = \lim_{x\to 2} \frac{a^x + a^{3-x} – (a^2 + a)}{a^{3-x} - a^{x/2}}
is equal to
(A) 2lna
(B) -\frac{1}{2}a
(C) \frac{a^2+a}{2}
(D) \frac{2}{3}(1-a) \hfill \textbf{((D))}
\item Q. Let
f(x) = ax^2 − 2 + \frac{1}{x},
where a is a real constant. The smallest value of a for which f(x) \ge 0 for all x > 0 is
(A) \frac{22}{3^3}
(B) \frac{23}{3^3}
(C) \frac{24}{3^3}
(D) \frac{25}{3^3} \hfill \textbf{((D))}
\item Q. Let f: R \to R be a continuous function satisfying
f(x) + \int_0^x t f(t) dt + x^2 = 0 for all x\in R.
Which of the following statements is true?
(A) \lim_{x\to\infty} f(x) = 2.
(B) \lim_{x\to\infty} f(x) = −2.
(C) f(x) has more than one point in common with the x-axis.
(D) f(x) is an odd function. \hfill \textbf{((B))}
\item Q. The curve y = 2x - 4x^3 intersects the horizontal line y = c at points A(a, c) and B(b, c)
(a < b), enclosing regions I above the line and II below as shown.
If the areas of I and II are equal, then a + b equals
(A) \frac{2}{\sqrt{7}}
(B) \frac{3}{\sqrt{7}}
(C) \frac{4}{\sqrt{7}}
(D) \frac{5}{\sqrt{7}} \hfill \textbf{((A))}
\item Q. Define f: R \to R by
f(x) = max{\|x\|, \|x − 1\|, . . ., \|x − 2n\|},
where n is a fixed natural number. Then
\int_0^{2n} f(x) dx
equals
(A) n
(B) n^2
(C) 3n
(D) 3n^2 \hfill \textbf{((D))}
\item Q. If p(x) is a cubic polynomial satisfying
p(1) = 3, p(0) = 2, p(-1) = 4,
then
\int_{-1}^1 p(x) dx =?
(A) 2
(B) 3
(C) 4
(D) 5 \hfill \textbf{((D))}
\item Q. Let x > 0 be fixed. Evaluate the integral
I = \int_0^{\infty} e^{-t} \|x - t\| dt.
(A) x + 2e^{-x} − 1
(B) x - 2e^{-x} + 1
(C) x + 2e^{-x} + 1
(D) -x - 2e^{-x} + 1 \hfill \textbf{((A))}
\item Q. Consider the function
f(x) = \begin{cases}
\frac{x+5}{x-2}, & x\ne 2, \\
1, & x = 2.
\end{cases}
Then f(f(x)) is discontinuous
(A) at all real numbers
(B) at exactly two values of x
(C) at exactly one value of x
(D) at exactly three values of x \hfill \textbf{((B))}
\item Q. Let f : [0, \pi] \to R be defined by
f(x) = \begin{cases}
\sin x, & x \text{ irrational,} \\
\tan^2x, & x \text{ rational.}
\end{cases}
The number of points in [0, \pi] at which f is continuous is
(A) 6
(B) 4
(C) 2
(D) 0 \hfill \textbf{((B))}
\item Q. A continuous function f : R \to R satisfies
f(x) = x + \int_0^x f(t) dt.
Which of the following identities holds for all x, y \in R?
(A) f(x + y) = f(x) + f(y)
(B) f(x + y) = f(x) f(y)
(C) f(x + y) = f(x) + f(y) + f(x) f(y)
(D) f(x + y) = f(xy) \hfill \textbf{((C))}
\item Q. Let f be a differentiable function satisfying
$\int_0^x (x-p+1) f(p) dp = x^4 + x^2 \forall x \in [-2,\infty)$.
Then the value of
$\int_{-2}^2 (f'(x) + f(x)) dx$
is
(A) 0
(B) 36
(C) 72
(D) 18 \hfill \textbf{((C))}
\item Q. Let
$P = \lim_{x\to 0} \frac{1 - (\cos 2x \cos 4x \cos 6x \cos 8x \cos 10x)^3}{\tan^2 x}$
Then the value of $P/10$ is
(A) 66
(B) 33
(C) 22
(D) 11 \hfill \textbf{((B))}
\item Q. If
$A = \lim_{n\to\infty} \sum_{r=1}^n \sin(\frac{r}{n})$, $B = \lim_{n\to\infty} \frac{1}{n} \sum_{r=1}^n \sin(\frac{\pi r}{n})$,
then the value of $2A + B$ is
(A) $1+\frac{2}{\pi}$
(B) $\frac{2}{\pi}$
(C) $\frac{\pi-1}{2}$
(D) $\frac{\pi}{2}-1$ \hfill \textbf{((A))}
\item Q. Let
$f(x) = \frac{1}{\sin^2 x - (x-a)^2}$, $x \in (0,1)$.
Which of the following statements is incorrect?
(A) $f(x)$ is always increasing in $(0,1)$.
(B) $f(x)$ is always decreasing in $(0,1)$.
(C) $f(x)$ is increasing if $a \in (2,3)$.
(D) $f(x)$ is decreasing if $a \in (2,3)$. \hfill \textbf{((D))}
\item Q. Let
$f(t) = \int_0^t e^{x^2} \frac{x^5}{(x^4+2x^2+2)^2} dx$.
Then $f(1) + f'(1)$ is equal to
(A) $\frac{3e}{10} - \frac{1}{4}$
(B) $\frac{7e}{50} - \frac{1}{4}$
(C) $\frac{7e}{50} - \frac{1}{2}$
(D) $\frac{2e}{5} - \frac{1}{2}$ \hfill \textbf{((B))}
\item Q. If the area bounded between the curves $y = \tan x$, $y = \sin x$, $x = 0$ and $x = \frac{\pi}{4}$ is A, then A is equal to
(A) $\frac{1}{2} \ln 2 + \frac{1}{\sqrt{2}}$
(B) $\frac{1}{2} \ln 2$
(C) $\frac{1}{\sqrt{2}}$
(D) none of these \hfill \textbf{((D))}
\item Q. Let $f(x) = x + \sin x$ and define
$I_1 = \int_0^\pi f(x)dx$, $I_2 = \int_0^\pi f^{-1}(x)dx$, $I_3 = \int_0^1 e^{-x^2}dx$.
Which of the following is true?
(A) $I_1 > I_2 > I_3$
(B) $I_3 > I_2 > I_1$
(C) $I_1 > I_3 > I_2$
(D) $I_3 > I_1 > I_2$ \hfill \textbf{((A))}
\item Q. Let
$f(x) = \int_0^x (\ln t)^2 dt$, $g(x) = \int_0^x \ln t dt$.
Find the area bounded between the curves $y = f'(x)$ and $y = g'(x)$ for $x$ from 1 to $e$.
(A) $e-3$
(B) 3
(C) $e$
(D) $3-e$ \hfill \textbf{((D))}
\item Q. Let
$L_1 = \lim_{x\to -\infty} \sec^{-1} x$, $L_2 = \lim_{x\to \infty} \sec^{-1} x$,
and evaluate
$\int_{L_1}^{L_2} \frac{x^3 + 1}{x^2 + 1} dx$.
(A) $\pi$
(B) $\frac{\pi}{2}$
(C) $\frac{\pi}{4}$
(D) 0 \hfill \textbf{((D))}
\item Q. Water is being filled at the rate of 100L/min in a right-circular conical vessel (vertex downwards) with axis vertical. Its semi-vertical angle satisfies $\tan\theta = \frac{5}{12}$. At the instant when the water depth is 3 m, the rate (in $m^3/min$) at which the wet curved surface area is increasing is:
(A) $\frac{13}{75}$
(B) $\frac{13}{25\pi}$
(C) $\frac{1}{5\pi}$
(D) $\frac{1}{11\pi}$ \hfill \textbf{((A))}
\item Q. If $y = y(x)$ is defined implicitly by
$\ln(x + y) = 6xy$,
then $\frac{d^2 y}{dx^2}$ at $x = 0$ is equal to
(A) 84
(B) 96
(C) 108
(D) 120 \hfill \textbf{((B))}
\item Q. Evaluate
$\lim_{x\to 0} (3 - 2 \cos x \sqrt{\cos 2x})^{\frac{x+3}{x^2}} = e^a$.
Then $a$ is equal to
(A) 6
(B) 9
(C) 12
(D) 15 \hfill \textbf{((B))}
\end{enumerate}
\subsection{Hard}
\begin{enumerate}
\item Q. Evaluate $\lim_{n\to\infty} n^{-n^2} ((n+1)(n+\frac{1}{2})(n+\frac{1}{2^2})\cdots(n+\frac{1}{2^{n-1}}))^n$. \hfill \textbf{((A))}
\item Q. Let $f(x) = \begin{cases} \frac{a(1-x \sin x) + b \cos x + 5}{x^2}, & x < 0, \\ 3, & x = 0, \\ (1+\frac{cx + dx^3}{x^2})^{1/x}, & x > 0. \end{cases}$ If f is continuous at x = 0, then (a, b, c, d) is \hfill \textbf{((A))}
\item Q. Let $f(x) = \begin{cases} \cos^{-1}\{ \cot x \}, & x < \frac{\pi}{2}, \\ [\pi |x|] - 1, & x \ge \frac{\pi}{2}, \end{cases}$ where {.} denotes the fractional-part function and [.] the greatest-integer function. The jump of discontinuity at $x = \frac{\pi}{2}$ is \hfill \textbf{((B))}
\item Q. Consider $f(x) = \begin{cases} x [x]^2 \log_{1+x} 2, & -1 < x < 0, \\ \frac{\ln(e^{x^2}+2\sqrt{\{x\}})}{\tan \sqrt{x}}, & 0 < x < 1, \end{cases}$ where [.] is the greatest-integer function and {.} the fractional part. Which of the following is true about f at x = 0? \hfill \textbf{((D))}
\end{enumerate}
\section{Calculus (Definite Integrals)}
\subsection{Easy}
\begin{enumerate}
\item Q. If $\int_{-1}^1 f(x)dx = 4$ and $\int_2^4 (3 - f(x))dx = 7$, then $\int_{-1}^2 f(x)dx$ is
(A) $1$
(B) $2$
(C) $4$
(D) $5$ \hfill \textbf{((D))}
\end{enumerate}
\section{Calculus (Limits)}
\subsection{Medium}
\begin{enumerate}
\item Q. $\lim_{x\to 0} \frac{\sqrt{1-\cos x^2}}{(1-\cos x)}$ equals-
(A) $\sqrt{2}$
(B) $\frac{1}{\sqrt{2}}$
(C) $1$
(D) none of these \hfill \textbf{((A))}
\end{enumerate}
\section{Calculus - Applications Of Derivatives}
\subsection{Medium}
\begin{enumerate}
\item Q. The sum of values of x that represent points of relative optima of the function $f(x) = 2 \cos x + x$, $0 \le x \le \pi$, is:
(A) $\frac{5\pi}{6}$
(B) $2\pi$
(C) $\pi$
(D) $\frac{\pi}{6}$ \hfill \textbf{((B))}
\end{enumerate}
\section{Calculus - Area Under Curve}
\subsection{Medium}
\begin{enumerate}
\item Q. Let the area enclosed by the lines $x + y = 2$, $y = 0$, $x = 0$ and the curve $f(x) = \min\{x^2 + \frac{3}{4}, 1 + |x|\}$ be A. Then the value of 12A is:
(A) 15
(B) 16
(C) 17
(D) 18 \hfill \textbf{((C))}
\end{enumerate}
\section{Calculus - Continuity And Differentiability}
\subsection{Medium}
\begin{enumerate}
\item Q. If the function
$f(x) = \begin{cases} ae^x + be^{-x}, & -1 \le x < 1, \\ cx^2, & 1 \le x \le 3, \\ ax^2 + 2cx, & 3 < x \le 4 \end{cases}$
is continuous for some real constants a, b, c and satisfies $f'(0) + f'(2) = e$, then the value of a is:
(A) $\frac{e}{e^2 + 3e + 13}$
(B) $\frac{e}{e^2 - 3e + 13}$
(C) $\frac{1}{e^2 - 3e + 13}$
(D) $\frac{e}{e^2 - 3e - 13}$ \hfill \textbf{((B))}
\end{enumerate}
\section{Calculus - Integrals And Functions}
\subsection{Medium}
\begin{enumerate}
\item Q. If $a = \lim_{n \to \infty} \sum_{k=1}^n \frac{2n}{n^2 + k^2}$ and $f(x) = \sqrt{\frac{1 - \cos x}{1 + \cos x}}$, $x \in (0,1)$, then which of the following holds?
(A) $\sqrt{2} f(\frac{\pi}{8}) = f'(\frac{\pi}{8})$
(B) $f(\frac{\pi}{8}) f'(\frac{\pi}{8}) = \sqrt{2}$
(C) $f(\frac{\pi}{8}) = \sqrt{2} f'(\frac{\pi}{8})$
(D) $2\sqrt{2} f(\frac{\pi}{8}) = f'(\frac{\pi}{8})$ \hfill \textbf{((A))}
\end{enumerate}
\section{Circles}
\subsection{Medium}
\begin{enumerate}
\item Q. Find the equation of the circle which cuts each of the circles $x^2 + y^2 = 4$, $x^2 + y^2 – 6x – 8y + 10 = 0$ and $x^2 + y^2 + 2x – 4y − 2 = 0$ at the extremities of a diameter.
(A) $x^2 + y^2 - 4x - 6y - 4 = 0$
(B) $x^2 + y^2 + 4x - 6y - 4 = 0$
(C) $x^2 + y^2 - 4x + 6y − 4 = 0$
(D) $x^2 + y^2 + 4x + 6y − 4 = 0$ \hfill \textbf{((A))}
\item Q. Let $S_1 : x^2 + y^2 + 4x - 6y - 12 = 0$ and $S_2 : x^2 + y^2 - 3x - 2y + 1 = 0$ be two distinct circles with centres $C_1$ and $C_2$ respectively. If A and B lie on their common chord as well as on a pair of common tangents, and $C_1A^2 - C_2B^2 = \frac{a}{b}$, where a, b $\in \mathbb{N}$ are coprime, then the value of $a+b$ is \hfill \textbf{((C))}
\end{enumerate}
\section{Combinatorics}
\subsection{Easy}
\begin{enumerate}
\item Q. In a football championship, each pair of teams plays exactly one match. A total of 153 matches were played. How many teams participated?
(A) 21
(B) 18
(C) 17
(D) 19 \hfill \textbf{((B))}
\end{enumerate}
\subsection{Medium}
\begin{enumerate}
\item Q. Let $E$ be the set of 26 English letters, $V = \{a, e, i, o, u\}$ the vowels, and $C = E \setminus V$ the consonants. How many four-letter “words” (strings, allowing repetition) contain at least one vowel and at least one consonant? \hfill \textbf{((A))}
\item Q. What is the HCF of $^{2n}C_1, ^{2n}C_3, ^{2n}C_5, \dots, ^{2n}C_{2n-1}$ (where $n = 2^m.q$, $q$ being an odd integer). \hfill \textbf{((A))}
\item Q. The sum $S = \frac{1}{9!} + \frac{1}{3!7!} + \frac{1}{5!5!} + \frac{1}{7!3!} + \frac{1}{9!}$ equals \hfill \textbf{((A))}
\item Q. The number of integral solutions of $x + y + z = 0$ with $x \ge -5, y \ge -5, z \ge -5$ is \hfill \textbf{((B))}
\item Q. Let $a = \ ^{x+2}P_{x+2}$, $b = \ ^{x}P_{11}$, $c = \ ^{x-11}P_{x-11}$, and suppose $a = 182 bc$. Then x is \hfill \textbf{((B))}
\item Q. In how many ways can 8 distinct books be distributed among 3 students if each student receives at least 2 books? \hfill \textbf{((C))}
\item Q. Find the number of permutations of the letters a, b, c, d, e, f, g taken all at a time if neither the pattern "beg” nor the pattern “cad” appears as consecutive letters. \hfill \textbf{((A))}
\item Q. If all distinct permutations of the letters of the word “PROPER" are listed in lexicographic (dictionary) order, what is the rank of “PROPER" in that list? \hfill \textbf{((C))}
\item Q. How many 4-letter “words" (arrangements) can be formed using the letters of the word “PARALLELOPIPED"? (The multiset of available letters is {P, P, P, L, L, L, A, A, E, E, R, O, I, D}.) \hfill \textbf{((B))}
\item Q. How many 6-digit numbers can be formed using digits from the set {1, 2, 3, 4, 5} if every digit that appears in the number appears at least twice? \hfill \textbf{((B))}
\item Q. A student may select at most n books from a shelf of $2n + 1$ distinct books. If the total number of ways in which he can select books (selecting 1 book, or 2 books, ..., up to n books) is 63, then n equals \hfill \textbf{((B))}
\item Q. In how many ways can the number 18900 be factored as a product of two coprime positive integers? \hfill \textbf{((C))}
\item Q. In how many ways can 12 identical apples be distributed among four children if each child receives at least 1 apple and at most 4 apples? \hfill \textbf{((C))}
\item Q. Find the number of positive integer solutions $(x, y, z)$ to $x+y+z \ge 150$, subject to $1 \le x, y, z \le 60$. \hfill \textbf{((B))}
\item Q. Let $N_1$ be the number of positive-integer solutions $(x, y, z) \in N^3$ of $xyz = 360$, and let $N_2$ be the total number of integer solutions $(x, y, z) \in Z^3$ of the same equation. Then $(N_1, N_2) =?$ \hfill \textbf{((B))}
\item Q. A person writes letters to six friends and addresses the corresponding envelopes. In how many ways can the letters be placed in the envelopes so that (i) all the letters are in the wrong envelopes (a derangement of 6), and (ii) at least two letters are in the wrong envelopes? \hfill \textbf{((A))}
\item Q. Find the sum of all 4-digit numbers formed using the digits 2, 4, 6, 8 exactly once each. \hfill \textbf{((C))}
\item Q. How many ways can three numbers be selected from the set $\{5^1, 5^2, ..., 5^{11}\}$ so that they form a geometric progression? \hfill \textbf{((B))}
\item Q. A regular polygon has 15 vertices. In how many ways can one choose 3 vertices to form a triangle such that none of its sides is a side of the polygon? \hfill \textbf{((C))}
\end{enumerate}
\section{Complex Numbers}
\subsection{Medium}
\begin{enumerate}
\item Q. Let $z = x+iy$ and $w = u+iv$ be complex numbers on the unit circle such that $z^2+w^2 = 1$. Then the number of ordered pairs $(z, w)$ is \hfill \textbf{((C))}
\item Q. Let $z = x + iy$ be a complex number. If $\frac{z-i}{z-1}$ is always purely imaginary, the locus of $z$ is
(A) the circle with center $(\frac{1}{2}, \frac{1}{2})$ and radius $\frac{1}{\sqrt{2}}$
(B) the circle with center $(-\frac{1}{2}, -\frac{1}{2})$ and radius $\frac{1}{\sqrt{2}}$
(C) the circle with center $(2, 2)$ and radius $\frac{1}{2}$
(D) the circle with center $(-2, -2)$ and radius $\frac{1}{2}$ \hfill \textbf{((A))}
\item Q. If $|z_1| = |z_2| = |z_3| = 1$ and $z_1 + z_2 + z_3 = \sqrt{2} + i$, then $z_1\bar{z_2} + z_2\bar{z_3} + z_3\bar{z_1}$ is
(A) purely imaginary
(B) purely real
(C) a positive real number
(D) a negative real number \hfill \textbf{((A))}
\item Q. A triangle with vertices represented by complex numbers $z_0, z_1, z_2$ has their opposite side lengths in the ratio $2 : \sqrt{6} : \sqrt{3} - 1$ respectively. Then
(A) $(z_2-z_0)^4 = -9(7+4\sqrt{3}) (z_1 - z_0)^4$
(B) $(z_2-z_0)^4 = 9(7+4\sqrt{3}) (z_1 - z_0)^4$
(C) $(z_2-z_0)^4 = (7+4\sqrt{3}) (z_1 - z_0)^4$
(D) None of these \hfill \textbf{((A))}
\item Q. If $\omega$ is an imaginary cube root of unity, then expression $(2-\omega) (2 - \omega^2)+2(3-\omega) (3 - \omega^2)+\dots+ (n - 1)(n - \omega) (n - \omega^2)$ is equal to
(A) $\frac{1}{4}n^2(n + 1)^2 + n$
(B) $\frac{1}{4}n^2(n + 1)^2 - n$
(C) $\frac{1}{4}n^2(n + 1) - n$
(D) $\frac{1}{4}n(n+1)^2 - n$ \hfill \textbf{((B))}
\item Q. If $|z - 2| = \min\{|z - 1|, |z - 5|\}$, where z is a complex number, then-
(A) $\text{Re}(z) = \frac{3}{2}$
(B) $\text{Re}(z) = \frac{7}{2}$
(C) $\text{Re}(z) \in \left\{\frac{3}{2}, \frac{7}{2}\right\}$
(D) None of these \hfill \textbf{((C))}
\item Q. If $\cos \alpha + \cos \beta + \cos \gamma = 0$ and $\sin \alpha + \sin \beta + \sin \gamma = 0$, then $\cos 3\alpha + \cos3\beta + \cos 3\gamma$ equals
(A) 0
(B) $\cos(\alpha + \beta + \gamma)$
(C) $3\cos(\alpha + \beta + \gamma)$
(D) $3\sin(\alpha + \beta + \gamma)$ \hfill \textbf{((C))}
\item Let $z_1$ and $z_2$ be $n^{th}$ roots of unity which subtend a right angle at the origin, then $n$ must be of the form (where $k$ is an integer)\\
(A) $4k +1$\\
(B) $4k +2$\\
(C) $4k+3$\\
(D) $4k$ \hfill \textbf{((D))}
\item Q. The value of $i^{57} + \frac{1}{i^{125}}$ is: (A) 0 (B) -2i (C) 2i (D) 2 \hfill \textbf{((A))}
\item Q. If $x = -5 + 2\sqrt{-4}$, then the value of $x^4 + 9x^3 + 35x^2 - x + 4$ is (A) -160 (B) 160 (C) -4 (D) 0 \hfill \textbf{((A))}
\item Q. Among the complex numbers z satisfying $|z - 25i| \le 15$, the one with the least positive argument is (A) 12 + 16i (B) -12+16i (C) 10i (D) 40i \hfill \textbf{((A))}
\item Q. Let $x_n = \cos(\frac{\pi}{2^n}) + i \sin(\frac{\pi}{2^n})$. Evaluate the infinite product $\prod_{n=1}^{\infty} x_n$. (A) -1 (B) 1 (C) 0 (D) $\infty$ \hfill \textbf{((A))}
\item Q. If $|\frac{z-i}{z+i}| = 1$, then the locus of z in the complex plane is (A) the x-axis (B) the y-axis (C) the line x = 1 (D) the line y = 1 \hfill \textbf{((A))}
\item Q. If $z_1$ and $z_2$ are complex numbers such that $|\frac{z_1 - 2z_2}{2 - z_1z_2}| = 1$ and $|z_2| \ne 1$, then $|z_1|$ equals (A) 1 (B) 2 (C) $\sqrt{2}$ (D) 4 \hfill \textbf{((B))}
\item Q. The locus of the complex number z in the Argand plane satisfying $\log_{1/2}(\frac{|z-1|+4}{3|z-1|-2}) > 1$, $|z-1| \ne \frac{2}{3}$, is (A) a circle (B) the interior of a circle (C) the exterior of a circle (D) none of these \hfill \textbf{((C))}
\item Q. Let $z_1, z_2, z_3$ be the complex affixes of the vertices A, B, C of an isosceles right triangle with right angle at C. Which of the following identities holds? (A) $(z_1-z_2)^2 = 2(z_1-z_3) (z_3-z_2)$ (B) $(z_1-z_2)^2 = -2(z_1-z_3) (z_3-z_2)$ (C) $(z_1-z_3)^2 = 2(z_1-z_2) (z_2-z_3)$ (D) $(z_2-z_3)^2 = 2 (z_2-z_1) (z_1-z_3)$ \hfill \textbf{((A))}
\item Q. Let $\alpha$ and $\beta$ be the non-real cube roots of unity. Then for any integer n, $\alpha^n + \beta^n = ?$ (A) $2\cos(\frac{2n\pi}{3})$ (B) $\cos(\frac{2n\pi}{3})$ (C) $2i\sin(\frac{2n\pi}{3})$ (D) $i\sin(\frac{2n\pi}{3})$ \hfill \textbf{((A))}
\item Q. Let $\alpha, \beta, \gamma$ be the roots of $x^3 - 3x^2 + 3x + 7 = 0$, and let $\omega$ be a non-real cube root of unity. Then the value of $\frac{\alpha-1}{\beta-1} + \frac{\beta-1}{\gamma-1} + \frac{\gamma-1}{\alpha-1}$ is (A) 3 (B) $3\omega$ (C) $3\omega^2$ (D) 0 \hfill \textbf{((C))}
\item Q. Let z be a nonzero complex number satisfying $(z + 1)\left(z + \frac{1}{z} + 1\right) = 1$. Compute the value of $(3z^{100} + 2z^{-100} + 1)(z^{100} + 2z^{-100} + 3)$. \hfill \textbf{((A))}
\item Q. If the distance of the point $2i$ from the line $(1 + i)z - (1 - i)\bar{z} + 2i = 0$ (where $i = \sqrt{-1}$) is d, then the value of $\sqrt{2}d$ is \hfill \textbf{((C))}
\end{enumerate}
\subsection{Hard}
\begin{enumerate}
\item Q. Let $z_1, z_2, z_3$ be the complex affixes of the vertices $A, B, C$ of an isosceles right triangle with right angle at $C$. Which of the following identities holds?
(A) $(z_1-z_2)^2 = 2(z_1-z_3) (z_3-z_2)$
(B) $(z_1-z_2)^2 = -2(z_1-z_3) (z_3-z_2)$
(C) $(z_1-z_3)^2 = 2(z_1-z_2) (z_2-z_3)$
(D) $(z_2-z_3)^2 = 2 (z_2-z_1) (z_1-z_3)$ \hfill \textbf{(A)}
\end{enumerate}
\section{Conic Sections}
\subsection{Medium}
\begin{enumerate}
\item Q. The exhaustive set of values of $a^2$ such that there exists a tangent to the ellipse $x^2 + a^2y^2 = a^2$ such that the portion of the tangent intercepted by the hyperbola $a^2x^2 – y^2 = 1$ subtends a right angle at the centre of the curves is \hfill \textbf{((A))}
\item Q. The sides $AC$ and $AB$ of $\triangle ABC$ touch the conjugate hyperbola of the hyperbola $\frac{x^2}{a^2} - \frac{y^2}{b^2} = 1$ . If the vertex $A$ lies on the ellipse $\frac{x^2}{a^2} + \frac{y^2}{b^2} = 1$, then the side $BC$ must touch \hfill \textbf{((D))}
\item Q. If P is any point on an ellipse with foci $S_1, S_2$ and eccentricity $e = \frac{1}{2}$, and $\angle PS_1S_2 = \alpha, \angle PS_2S_1 = \beta, \angle S_1 PS_2 = \gamma$, then $\lambda = \frac{\cot(\alpha/2) + \cot(\beta/2)}{\cot(\gamma/2)}$, is equal to \hfill \textbf{((B))}
\end{enumerate}
\section{Conic Sections - Hyperbola}
\subsection{Medium}
\begin{enumerate}
\item Q. The equation of the hyperbola whose foci are at $(\pm 8,0)$ and whose latus rectum has length 24 is:
(A) $x^2 - 3y^2 = 48$
(B) $x^2 - 4y^2 = 48$
(C) $4x^2 - y^2 = 48$
(D) $3x^2 - y^2 = 48$ \hfill \textbf{((D))}
\end{enumerate}
\section{Conic Sections - Pair Of Straight Lines}
\subsection{Medium}
\begin{enumerate}
\item Q. Let a and b be non-zero real numbers. Then the equation $(ax^2 + by^2 + c) (x^2 - 5xy + 6y^2) = 0$ represents a circle and two straight lines. What are the values of $(a,b)$?
(A) $(1, \frac{7}{3})$
(B) $(\frac{2}{3}, 1)$
(C) $(\frac{1}{3}, 2)$
(D) $(\frac{1}{3}, \frac{5}{3})$ \hfill \textbf{((C))}
\end{enumerate}
\section{Conics}
\subsection{Medium}
\begin{enumerate}
\item Q. If S and S' are the foci of the ellipse $\frac{x^2}{25} + \frac{y^2}{16} = 1$, and P is any point on it, then range of values of SP$\cdot$S'P is
(A) $9 \le f(\theta) \le 16$
(B) $9 \le f(\theta) \le 25$
(C) $16 \le f(\theta) \le 25$
(D) $1 \le f(\theta) \le 16$ \hfill \textbf{((C))}
\end{enumerate}
\section{Continuity}
\subsection{Medium}
\begin{enumerate}
\item Q. Let
$f(x) = \begin{cases} \frac{1}{e^{1/x} + 1}, & x \ne 0, \\ 0, & x = 0. \end{cases}$
Which of the following is true?
(A) $\lim_{x \to 0^+} f(x) = 1$
(B) $\lim_{x \to 0^-} f(x) = 0$
(C) f(x) is discontinuous at x = 0
(D) f(x) is continuous at x = 0 \hfill \textbf{((C))}
\item Q. Consider
$f(x) = \begin{cases} x [x]^2 \log_{1+x} 2, & -1 < x < 0, \\ \frac{\ln(e^{x^2}+2\sqrt{\{x\}})}{\tan \sqrt{x}}, & 0 < x < 1, \end{cases}$
where $[.]$ is the greatest-integer function and \{\cdot\} the fractional part. Which of the following is true about f at x = 0?
(A) If f(0) = ln 2, then f is continuous at x = 0.
(B) If f(0) = 2, then f is continuous at x = 0.
(C) If f(0) = $e^2$, then f is continuous at x = 0.
(D) f has an irremovable discontinuity at x = 0. \hfill \textbf{((D))}
\item Q. Let
$f(x) = \begin{cases} a \sin^{2n} x, & x \ge 0, n \to \infty, \\ b \cos^{2m} x - 1, & x < 0, m \to \infty. \end{cases}$
Then which statement is true?
(A) $f(0^-) \ne f(0^+)$
(B) $f(0^+) \ne f(0)$
(C) $f(0^-) = f(0)$
(D) f is continuous at x = 0 \hfill \textbf{((A))}
\item Q. f is a continuous function on R. Given that
$x^2 + (f(x) - 2)x - \sqrt{3} f(x) + 2\sqrt{3} - 3 = 0 \quad \forall x$,
then the value of $f(\sqrt{3})$ is
(A) $\frac{2(\sqrt{3}-2)}{\sqrt{3}}$
(B) $2(1 - \sqrt{3})$
(C) 0
(D) cannot be determined \hfill \textbf{((B))}
\end{enumerate}
\subsection{Hard}
\begin{enumerate}
\item Q. Let
$f(x) = \begin{cases} \frac{a(1-x \sin x) + b \cos x + 5}{x^2}, & x < 0, \\ 3, & x = 0, \\ (1+\frac{cx + dx^3}{x^2})^{1/x}, & x > 0. \end{cases}$
If f is continuous at x = 0, then (a, b, c, d) is
(A) (-1, -4, 0, ln 3)
(B) (-1, -4, 1, ln 3)
(C) (1, -4, 0, ln 3)
(D) (-1, -4, 0, 3) \hfill \textbf{((A))}
\end{enumerate}
\section{Coordinate Geometry}
\subsection{Medium}
\begin{enumerate}
\item Q. Two circles of equal radius intersect at $(0,1)$ and $(0, -1)$. The tangent at $(0,1)$ to one circle passes through the center of the other. The distance between their centers is
(A) 2
(B) $2\sqrt{2}$
(C) 1
(D) $\sqrt{2}$ \hfill \textbf{((A))}
\item Q. The parabola $y = (x - 2)^2 - 1$ meets the line $x - y = 3$ in two points. The tangents to the parabola at those points intersect at
(A) $(\frac{5}{2},-1)$
(B) $(\frac{3}{2},1)$
(C) $(-\frac{5}{2},1)$
(D) $(-\frac{3}{2},-1)$ \hfill \textbf{((A))}
\item Q. A line is drawn from $A(-2,0)$ to intersect the curve $y^2 = 4x$ in $P$ and $Q$ in the first quadrant such that $\frac{1}{AP}+\frac{1}{AQ} < \frac{1}{4}$, then slope of the line is always
(A) $> \sqrt{3}$
(B) $< \frac{1}{\sqrt{3}}$
(C) $> \sqrt{2}$
(D) $> \frac{1}{\sqrt{3}}$ \hfill \textbf{((A))}
\item Q. From a point R(5,8), two tangents RP and RQ are drawn to a given circle S = 0 whose radius is 5. If circumcentre of the triangle PQR is (2,3), then the equation of circle S=0 is
(A) $x^2 + y^2 + 2x + 4y - 20 = 0$
(B) $x^2 + y^2 + x + 2y - 20 = 0$
(C) $x^2 + y^2 - x - 2y - 20 = 0$
(D) $x^2 + y^2 - 4x - 6y - 12 = 0$ \hfill \textbf{((A))}
\item Q. A triangle has two vertices A = (0, 1), B = (2, 3), and its centroid is at G = (3,3). Find the third vertex C, the circumcenter P, the circumradius R, and the orthocenter H.\n(A) C = (3,5), P = (-\frac{9}{2}, \frac{15}{2}), R = \frac{1}{2}\sqrt{10}, H = (14, −6)\n(B) C = (3,5), P = (-\frac{9}{2}, \frac{1}{2}), R = \frac{1}{2}\sqrt{5}, H = (14, 6)\n(C) C = (5,3), P = (-\frac{9}{2}, \frac{15}{2}), R = \frac{1}{2}\sqrt{10}, H = (−14, -6)\n(D) C = (3,5), P = (-\frac{9}{2}, \frac{15}{2}), R = \frac{3}{2}\sqrt{10}, H = (14, −6) \hfill \textbf{(A)}
\item Q. A rigid rod of length $l$ has its ends always sliding on the x- and y-axes, respectively. Let P(x, y) be the point on the rod that divides it in the ratio $m_1 : m_2$ (measured from the x-axis end towards the y-axis end). The locus of P is\n(A) $m_1 x^2 + m_2 y^2 = \frac{l^2}{(m_1 + m_2)^2}$\n(B) $(m_2 x)^2 + (m_1 y)^2 = (\frac{m_1 m_2 l}{m_1+m_2})^2$\n(C) $(m_1 x)^2 + (m_2 y)^2 = (\frac{m_1 m_2 l}{m_1+m_2})^2$\n(D) none of these \hfill \textbf{(C)}
\item Q. Let A = (a,0) and B = (-a, 0) be fixed points. A variable point C = (x, y) moves so that cot ∠A + cot ∠B = λ, where λ is a constant. The locus of C is\n(A) λy = 2a\n(B) y = λα\n(C) ay = 2x\n(D) none of these \hfill \textbf{(A)}
\item Q. Two points A and B move along the positive x- and y-axes respectively so that OA+OB = K (a constant). Let M = (x, y) be the foot of the perpendicular from the origin O onto the line AB. The locus of M is\n(A) $(x - y) (x^2 + y^2) = K xy$\n(B) $(x + y) (x^2 – y^2) = K xy$\n(C) $(x + y) (x^2 + y^2) = K xy$\n(D) $(x + y) (x^2 + y^2) = \frac{1}{2} K xy$ \hfill \textbf{(C)}
\item Q. A variable line through P(-2, -3) meets the pair of lines $x^2 + 3y^2 + 4xy – 8x – 6y – 9 =0$ at points Q and R. If PQ · PR = 20, then the equations of the line are\n(A) $x − y − 1 = 0$ and $3x - y + 3 = 0$\n(B) $x + y − 1 = 0$ and $3x + y + 3 = 0$\n(C) $y − x − 1 = 0$ and $y - 3x-3=0$\n(D) none of these \hfill \textbf{(A)}
\item Q. Find the equations of the two legs of an isosceles right-angled triangle whose hypotenuse is $3x + 4y = 4$ and whose opposite vertex is C = (2, 2).\n(A) $7y - x - 12 = 0$ and $7x + y - 16 = 0$\n(B) $7y + x - 12 = 0$ and $7x - y + 16 = 0$\n(C) $y - 7x - 12 = 0$ and $x + 7y - 16 = 0$\n(D) $x - 7y+ 12 = 0$ and $7x + y - 16 = 0$ \hfill \textbf{(A)}
\item Q. A circle of radius 3 has its center on the line $y = x - 1$ and passes through the point (7,3). Which of the following gives the equations of all such circles?\n(A) $x^2 + y^2 - 8x - 6y + 16 = 0$ and $x^2 + y^2 – 14x – 12y + 76 = 0$\n(B) $x^2 + y^2 - 8x - 6y – 16 = 0$ and $x^2 + y^2 – 14x – 12y – 76 = 0$\n(C) $x^2 + y^2-6x - 8y + 16 = 0$ and $x^2 + y^2 — 12x – 14y + 76 = 0$\n(D) $x^2 + y^2 - 8x - 6y + 16 = 0$ and $x^2 + y^2 – 14x – 12y + 28 = 0$ \hfill \textbf{(A)}
\item Q. A variable line $y = 2x + a$ lies entirely between (and does not meet or touch) the two circles $C_1: x^2 + y^2 - 2x - 2y + 1 = 0$ and $C_2: x^2 + y^2 – 16x – 2y + 61 = 0$. Find the range of the parameter a.\n(A) $-15 < a < -1$\n(B) $2\sqrt{5}-15 < a < -1-\sqrt{5}$\n(C) $-1-\sqrt{5} < a < 2\sqrt{5}-15$\n(D) none of these \hfill \textbf{(B)}
\item Q. Find the equation of the circle whose center is (3, -1) and which cuts off a chord of length 6 on the line $2x – 5y + 18 = 0$.\n(A) $(x - 3)^2 + (y + 1)^2 = 38$\n(B) $(x + 3)^2 + (y − 1)^2 = 38$\n(C) $(x - 3)^2 + (y + 1)^2 = \sqrt{38}$\n(D) none of these \hfill \textbf{(A)}
\item Q. The straight line $ax + by = 2$, $a, b \neq 0$, is tangent to the circle $x^2 + y^2 - 2x = 3$ and is normal to the circle $x^2 + y^2 - 4y = 6$. Find (a, b).\n(A) (1, -1)\n(B) (1, 2)\n(C) (-\frac{4}{3}, 1)\n(D) (2, 1) \hfill \textbf{(C)}
\item Q. From the point (14,7), three normals are drawn to the parabola $y^2 - 16x - 8y = 0$. The feet of these normals are\n(A) (0,0), (8,16), (3, -4)\n(B) (0,0), (4,8), (3, -4)\n(C) (1,1), (8,16), (3, -4)\n(D) (0,0), (8,16), (4, -3) \hfill \textbf{(A)}
\item Q. Let $y^2 = 4ax$ be a parabola. A chord of this parabola subtends a right angle at the vertex. If (h, k) is the midpoint of such a chord, then its locus is\n(A) $k^2 = 2a (h - 4a)$\n(B) $k^2 = 4a (h - 4a)$\n(C) $k^2 = 2a (h + 4a)$\n(D) $k^2 = a (h - 4a)$ \hfill \textbf{(A)}
\item Q. The parabolas $y^2 = 4ax$ and $y^2 = 4c (x - b)$, with $a \neq c$, admit a common normal (other than the x-axis). Which of the following is the correct necessary and sufficient condition?\n(A) $\frac{b}{a-c} > 2$\n(B) $\frac{b}{a-c} < 2$\n(C) $\frac{b}{c-a} > 2$\n(D) none of these \hfill \textbf{(A)}
\item Q. A line $lx + my = n$ is a normal to the ellipse $\frac{x^2}{a^2} + \frac{y^2}{b^2} = 1$ if and only if its coefficients satisfy which of the following conditions?\n(A) $n^2(\frac{l^2}{a^2} + \frac{m^2}{b^2}) = (a^2-b^2)^2$\n(B) $n^2(\frac{l^2}{a^2} + \frac{m^2}{m^2}) = (a^2-b^2)^2$\n(C) $n^2(\frac{l}{a^2} + \frac{b}{m}) = a^2-b^2$\n(D) $n^2(l^2 + b^2) = a^2 + b^2$ \hfill \textbf{(A)}
\item Q. A normal to the hyperbola $\frac{x^2}{a^2} - \frac{y^2}{b^2} = 1$ meets the x- and y-axes at M and N, respectively. From M draw MP $||$ x-axis, and from N draw NP $||$ y-axis, meeting at P(x, y). The locus of P is\n(A) $a^2x^2 - b^2y^2 = (a^2 + b^2)^2$\n(B) $a^2x^2 - b^2y^2 = a^2 + b^2$\n(C) $a^2x^2 – b^2y^2 = (a^2 – b^2)^2$\n(D) none of these \hfill \textbf{(A)}
\item Two straight lines are perpendicular to each other. One of them touches the parabola $y^2 = 4a(x + a)$ and the other touches $y^2 = 4b(x + b)$. Their point of intersection lies on the line\\
(A) $x - a + b = 0$\\
(B) $x - a - b = 0$\\
(C) $x + a + b = 0$\\
(D) $x - a - b = 0$ \hfill \textbf{((C))}
\item If $d$ is the distance between parallel tangents with positive slope to $y^2 = 4x$ and $x^2 + y^2 – 2x + 4y – 11 = 0$, then\\
(A) $10 < d < 2$\\
(B) $4 < d < 6$\\
(C) $d < 4$\\
(D) none of these \hfill \textbf{((C))}
\item Find the condition for the quadrilateral to be concyclic whose consecutive sides are $a_rx + b_ry + c_r = 0$; $r = 1, 2, 3, 4.$\\
(A) $(a_1b_2-a_2b_1) (a_3a_4 + b_3b_4) \pm (a_3b_4 - a_4b_3) (a_1a_2 + b_1b_2) = 0$\\
(B) $(a_1b_2 + a_2b_1) (a_3a_4 + b_3b_4) \pm (a_3b_4 + a_4b_3) (a_1a_2 + b_1b_2) = 0$\\
(C) $(a_1b_2 + a_2b_1) (a_3a_4 + b_3b_4) \pm (a_3b_4 - a_4b_3) (a_1a_2 + b_1b_2) = 0$\\
(D) $(a_1b_2 - a_2b_1) (a_3a_4 + b_3b_4) \pm (a_3b_4 + a_4b_3) (a_1a_2 + b_1b_2) = 0$ \hfill \textbf{((A))}
\item Q. If the tangent at (1, 2) to the circle $C_1 : x^2+y^2 = 5$ intersects the circle $C_2 : x^2+y^2 = 9$ at A and B and tangents at A and B to the second circle meet at point C, then the coordinates of C are given by
(A) (4,5)
(B) $(\frac{9}{5}, \frac{18}{5})$
(C) $(4, -5)$
(D) $(\frac{9}{5}, \frac{18}{5})$ \hfill \textbf{((D))}
\item Q. If the points $\left(\frac{a^3}{a-1}, \frac{a^2-3}{a-1}\right), \left(\frac{b^3}{b-1}, \frac{b^2-3}{b-1}\right), \left(\frac{c^3}{c-1}, \frac{c^2-3}{c-1}\right)$ are collinear, and $\alpha abc + \beta (a + b + c) = \gamma (ab + bc + ca)$, where $\alpha, \beta, \gamma \in \mathbb{N}$ are coprime, then the value of $[\alpha + \beta + \gamma]$ is \hfill \textbf{((C))}
\end{enumerate}
\subsection{Hard}
\begin{enumerate}
\item Q. The lines $L_1 : y - \sqrt{3}x + \sqrt{3} = 0$, $L_2 : y + \sqrt{3}x = 6 + \sqrt{3}$, $L_3 : y = 0$ enclose a triangle. Its incircle radius is
(A) 4
(B) 3
(C) 2
(D) 1 \hfill \textbf{((D))}
\item Q. A variable point P divides every chord (having slope 2 ) of the parabola $y^2 = 4x$ into two parts such that the length of one part is twice the other. If locus of P is also a parabola, then find the coordinates of focus of the locus of point P.
(A) $\left(\frac{1}{3}, \frac{8}{9}\right)$
(B) $\left(\frac{1}{9}, \frac{8}{3}\right)$
(C) $\left(\frac{2}{9}, \frac{8}{3}\right)$
(D) $\left(\frac{1}{3}, \frac{8}{9}\right)$ \hfill \textbf{((D))}
\item Q. If a circle of radius R passes through the origin O and intersects the coordinate axes at A and B, then the locus of the foot of perpendicular from O on AB is :
(A) $(x^2 + y^2) (x + y) = R^2xy$
(B) $(x^2 + y^2)^3 = 4R^2x^2y^2$
(C) $(x^2 + y^2)^2 = 4R^2x^2y^2$
(D) $(x^2 + y^2)^2 = 4Rx^2y^2$ \hfill \textbf{((B))}
\end{enumerate}
\section{Coordinate Geometry (Circles)}
\subsection{Medium}
\begin{enumerate}
\item Q. The distinct points $A(0, 0)$, $B(0, 1)$, $C(1, 0)$ and $D(2a, 3a)$ are concyclic, then
(A) $a$ can attain only rational values
(B) $a$ is irrational
(C) cannot be concyclic for any $a$
(D) none of the above \hfill \textbf{((A))}
\end{enumerate}
\section{Coordinate Geometry (Ellipse)}
\subsection{Medium}
\begin{enumerate}
\item Q. Let $Q = (3, \sqrt{5})$ and $R = (7,3\sqrt{5})$. A point $P$ in the XY-plane varies in such a way that perimeter of $\triangle PQR$ is $16$. Then the maximum area of $\triangle PQR$ is
(A) 6 sq. units
(B) 12 sq. units
(C) 18 sq. units
(D) 9 sq. units \hfill \textbf{((B))}
\item Q. If $PQR$ is an equilateral triangle inscribed in the auxiliary circle of the ellipse $\frac{x^2}{a^2} + \frac{y^2}{b^2} = 1(a > b)$ and $P'Q'R'$ is the corresponding triangle inscribed within the ellipse, then centroid of the triangle $P'Q'R'$ lies
(A) at the centre of ellipse.
(B) at the focus of ellipse.
(C) between focus and centre on major axis.
(D) none of these. \hfill \textbf{((A))}
\end{enumerate}
\section{Coordinate Geometry (Hyperbola)}
\subsection{Hard}
\begin{enumerate}
\item Q. Let $LL'$ be the latus rectum through the focus $S$ of a hyperbola and $A$ be the farther vertex of the conic. If $\triangle ALL'$ is equilateral then its eccentricity is
(A) $\sqrt{3}$
(B) $\sqrt{3}+1$
(C) $\frac{(\sqrt{3}+1)}{\sqrt{2}}$
(D) $1 + \frac{1}{\sqrt{3}}$ \hfill \textbf{((D))}
\end{enumerate}
\section{Coordinate Geometry (Straight Lines, Area)}
\subsection{Medium}
\begin{enumerate}
\item Q. A system of line is given as $y = m_i x + c_i$, where $m_i$ can take any value out of $0, 1, -1$ and when $m_i$ is positive then $c_i$ can be $1$ or $-1$ when $m_i$ equal $0, c_i$ can be $0$ or $1$ and when $m_i$ equals $-1, c_i$ can take $0$ or $2$. Then the area enclosed by all these straight line is
(A) $\frac{3}{2}(\sqrt{2}-1)$ sq. units
(B) $\frac{3}{\sqrt{2}}$ sq. units
(C) $\frac{3}{2}$ sq. units
(D) None of these \hfill \textbf{((C))}
\end{enumerate}
\section{Determinants}
\subsection{Medium}
\begin{enumerate}
\item Q. If $\Delta(x) = \begin{vmatrix} 3 & 3x & 3x^2 + 2a^2 \\ 3x & 3x^2 + 2a^2 & 3x^3 + 6a^2x \\ 3x^2 + 2a^2 & 3x^3 + 6a^2x & 3x^4 + 12a^2x^2 + 2a^4 \end{vmatrix}$, then which of the following
options are always true for all $a \in R$ ?
(A) $\Delta'(x) = 0$
(B) $\Delta(x)$ is not independent of $x$
(C) $\int_{0}^{1} \Delta(x)dx = 16a^6$
(D) $y=\Delta(x)$ is the equation of a parabola. \hfill \textbf{((A))}
\end{enumerate}
\section{Differential Calculus}
\subsection{Medium}
\begin{enumerate}
\item Q. Tangent drawn at any point on the curve $(\frac{x}{a})^{2/3}+(\frac{y}{b})^{2/3} = 1$ has p&q as its intercepts on the x and y axis respectively. Then, $\frac{p^2}{a^2}+\frac{q^2}{b^2} =$
(A) 2
(B) 1
(C) 3
(D) 0 \hfill \textbf{((B))}
\end{enumerate}
\section{Differential Equations}
\subsection{Easy}
\begin{enumerate}
\item Determine which of the following functions are homogeneous:
(i) $f(x, y) = x^2 - xy$, (ii) $f(x, y) = \frac{xy}{x + y^2}$, (iii) $f(x, y) = \sin(xy)$.

(A) (i) is homogeneous; (ii) and (iii) are not.
(B) (ii) is homogeneous; (i) and (iii) are not.
(C) (iii) is homogeneous; (i) and (ii) are not.
(D) All three are homogeneous. \hfill \textbf{((A))}
\end{enumerate}
\subsection{Medium}
\begin{enumerate}
\item Find the order and degree of the following differential equations:
(i) $\frac{d^{2}y}{dx^{2}} = \sqrt[3]{\frac{dy}{dx} + 3}$, (ii) $\frac{d^{2}y}{dx^{2}} = \sin(\frac{dy}{dx})$, (iii) $\frac{dy}{dx} = \sqrt{3x + 5}$.

(A) (i) order 2, degree 3; (ii) order 2, degree not defined; (iii) order 1, degree 1.
(B) (i) order 2, degree 2; (ii) order 2, degree 1; (iii) order 1, degree 2.
(C) (i) order 1, degree 3; (ii) order 2, degree not defined; (iii) order 1, degree 1.
(D) (i) order 2, degree 3; (ii) order 1, degree not defined; (iii) order 1, degree 1. \hfill \textbf{((A))}
\item Find the differential equation of all parabolas whose axis is parallel to the $x$-axis and having latus rectum $a$.

(A) $y'' + 2(y')^3 = 0$
(B) $ay'' + 2(y')^3 = 0$
(C) $ay'' + 2(y')^2 = 0$
(D) $y'' + 2(y')^2 = 0$ \hfill \textbf{((B))}
\item Find the differential equation whose solution represents the one-parameter family $c (y + c)^2 = x^3$.

(A) $(\frac{2x}{3}y' - y) (\frac{2x}{3}y')^2 = x^3$
(B) $(2xy' - y) (\frac{2x}{3}y')^2 = x^3$
(C) $(\frac{2x}{3}y' - y) (2xy')^2 = x^3$
(D) $(\frac{2x}{3}y' + y) (\frac{2x}{3}y')^2 = x^3$ \hfill \textbf{((A))}
\item Solve the differential equation $xy \frac{dy}{dx} = \frac{1 + y^2}{1 + x^2} (1 + x + x^2)$.
The general solution can be written in which of the following forms?

(A) $\sqrt{1 + y^2} = Cxe^{-\tan^{-1}x}$
(B) $1 + y^2 = Cx^2 e^{2 \tan^{-1} x}$
(C) $\sqrt{1 + y^2} = Cxe^{\tan^{-1} x}$
(D) $\ln(1 + y^2) = 2\ln x + 2\tan^{-1}x + \ln C$ \hfill \textbf{((C))}
\item Solve the differential equation $\frac{dy}{dx} = \cos(x + y) - \sin(x + y)$.
The general solution can be expressed as:

(A) $\ln|\tan \frac{x + y}{2} - 1| = x + C$
(B) $\ln|1 - \tan \frac{x + y}{2}| = -x + C$
(C) $\ln|\tan \frac{x + y}{2} + 1| = x + C$
(D) $\ln|1 + \tan \frac{x + y}{2}| = -x + C$ \hfill \textbf{((B))}
\item Solve the differential equation $\frac{x+y \frac{dy}{dx}}{y-x \frac{dy}{dx}} = x^2 + 2y^2 + \frac{y^4}{x^2}$.
The general solution can be written implicitly as

(A) $\frac{1}{2(x^2+y^2)} = \frac{y}{x} + C$
(B) $\frac{1}{2(x^2+y^2)} = \frac{y}{x} + C$
(C) $\frac{1}{x^2+y^2} = \frac{y}{x} + C$
(D) $\frac{1}{2(x^2+y^2)} = -\frac{y}{x} + C$ \hfill \textbf{((A))}
\item The solution of the differential equation $\frac{dy}{dx} = \frac{\sin y + x}{\sin 2y - x \cos y}$ is:

(A) $\sin^2 y = x \sin y + \frac{x^2}{2} + C$
(B) $\sin^2 y = x \sin y - \frac{x^2}{2} + C$
(C) $\sin^2 y = x + \sin y + \frac{x^2}{2} + C$
(D) $\sin^2 y = x - \sin y + \frac{x^2}{2} + C$ \hfill \textbf{((A))}
\item Solve the differential equation $(1+ 2e^{x/y}) dx + 2e^{x/y} (1-\frac{x}{y}) dy = 0$.
The general solution can be written in which of the following forms?

(A) $x + 2y e^{x/y} = C$
(B) $y + 2x e^{x/y} = C$
(C) $\ln y + \ln(\frac{x}{y} + 2e^{x/y}) = \ln C$
(D) $x + 2y e^{-x/y} = C$ \hfill \textbf{((A))}
\item The solution of the differential equation $(x^2 - 1) \frac{dy}{dx} + 2xy = \frac{1}{x^2-1}$ is:

(A) $y (x^2 - 1) = \frac{1}{2} \ln|\frac{x-1}{x+1}| + C$
(B) $y (x^2 + 1) = \frac{1}{2} \ln|\frac{x-1}{x+1}| - C$
(C) $y (x^2 - 1) = \frac{5}{2} \ln|\frac{x-1}{x+1}| + C$
(D) none of these \hfill \textbf{((A))}
\item Solve the differential equation $x \frac{dy}{dx} + y = x^3 y^6$.
The general solution can be written as

(A) $y^{-5} = \frac{5}{2} x^3 + C x^5$
(B) $y^{-5} = - \frac{5}{2} x^3 + C x^5$
(C) $y^{-5} = \frac{5}{2} x^{-3} + C x^5$
(D) $y^{-5} = \frac{5}{2} x^3 + C x^{-5}$ \hfill \textbf{((A))}
\item Find the value of $k$ such that the family of parabolas $y = cx^2 + k$ is the orthogonal trajectories of the family of ellipses $x^2 + 2y^2 - y = c$.

(A) $k = 0$
(B) $k = \frac{1}{4}$
(C) $k = \frac{1}{2}$
(D) $k = -\frac{1}{4}$ \hfill \textbf{((B))}
\item Solve the differential equation $\frac{y+\sin x \cos^2(xy)}{\cos^2(xy)} dx + (\frac{x}{\cos^2(xy)} + \sin y) dy = 0$.
Its general solution can be written as

(A) $\tan(xy) - \cos x - \cos y = C$
(B) $\tan(xy) + \cos x + \cos y = C$
(C) $\tan(xy) - \sin x - \sin y = C$
(D) $\tan(xy) + \sin x + \sin y = C$ \hfill \textbf{((A))}
\item A tank initially contains 5000 L of water in which 20 kg of salt is dissolved. Brine containing 0.03 kg of salt per liter enters at 25 L/min, and the well-mixed solution drains at the same rate. How much salt remains in the tank after 30 minutes?

(A) 38.1 kg
(B) 45.0 kg
(C) 33.6 kg
(D) 50.0 kg \hfill \textbf{((A))}
\item Q. Solve the initial-value problem $\frac{dy}{dx} + \frac{3}{\cos^2 x} y = \frac{1}{\cos^2 x}$, $x \in (-\frac{\pi}{2}, \frac{\pi}{2})$, with $y(\frac{\pi}{4}) = \frac{4}{3}$. Then $y(-\frac{\pi}{4}) =$
(A) $\frac{4}{3}$
(B) $\frac{1}{3} + e^6$
(C) $\frac{1}{3} + e^3$
(D) $-\frac{1}{3}$ \hfill \textbf{((B))}
\item Q. The general solution of the differential equation $\frac{dy}{dx} = \sqrt{1 - x^2 - y^2 + x^2y^2}$ is:
(A) $\sin^{-1}y = x\sqrt{1-x^2} + c$
(B) $2\sin^{-1} y = x\sqrt{1-x^2} + \sin^{-1}x + c$
(C) $\cos^{-1} y = x\sqrt{1-x^2} + c$
(D) $2\cos^{-1}y = x\sqrt{1-x^2} + \sin^{-1}x + c$ \hfill \textbf{((B))}
\item Q. A tangent to the curve, $y = f(x)$ at P$(x, y)$ meets x-axis at A and y-axis at B. If AP : BP = 1 : 3 and $f(1) = 1$, then the curve also passes through the point
(A) $(3,24)$
(B) $(1,4)$
(C) $(2,\frac{1}{8})$
(D) $(3,\frac{1}{28})$ \hfill \textbf{((C))}
\item Q. The general solution of the differential equation $\frac{dy}{dx} = (x^3 - 2x \tan^{-1} y) (1 + y^2)$ is- \hfill \textbf{((B))}
\item Q. If $f(x)$ satisfies the differential equation $\frac{dy}{dx} = (x - y)^2$ and given that $y(1) = 1$, then
(A) $-\ln \left|\frac{1-x+y}{1+x-y}\right| = 2(x - 1)$
(B) $\ln \left|\frac{2-x}{2-y}\right| = x+y-1$
(C) $\ln \left|\frac{1-x+y}{1+x-y}\right| = 2(x - 1)$
(D) $\ln \left|\frac{1-x+y}{1+x-y}\right| + \ln |x| = 0$ \hfill \textbf{((A))}
\end{enumerate}
\subsection{Hard}
\begin{enumerate}
\item Solve the differential equation $y(xy+1) dx + x(1+xy+x^2y^2) dy = 0$.
The general solution can be written implicitly as

(A) $2x^2y^2 \ln(y) - 2xy - 1 = Cx^2y^2$
(B) $2x^2y^2 \ln(xy) - 2xy - 1 = Cx^2y^2$
(C) $2x^2y^2 \ln((\frac{v}{x})) - 2xy - 1 = Cx^2y^2$
(D) $2x^2y^2 \ln(x^2y^2) - 2xy - 1 = Cx^2y^2$ \hfill \textbf{((A))}
\end{enumerate}
\section{Discontinuity}
\subsection{Medium}
\begin{enumerate}
\item Q. Let
$f(x) = \begin{cases} [\cos^{-1}\{cotx\}], & x < \frac{\pi}{2}, \\ [\pi[x] - 1], & x \ge \frac{\pi}{2}, \end{cases}$
where \{\} denotes the fractional-part function and $[.]$ the greatest-integer function. The jump of discontinuity at x = $\frac{\pi}{2}$ is
(A) $\frac{\pi}{2}$
(B) $\frac{\pi}{2} - 1$
(C) $\pi-1$
(D) $1-\frac{\pi}{2}$ \hfill \textbf{((B))}
\end{enumerate}
\section{Domain Of Composite Functions}
\subsection{Medium}
\begin{enumerate}
\item Q. Let
$\quad f(x) = \begin{cases} [x], & -2 \le x \le -1, \\ |x| + 1, & -1 < x \le 2, \end{cases} \quad g(x) = \begin{cases} [x], & -\pi \le x < 0, \\ \sin x, & 0 \le x \le \pi. \end{cases}$
Define $\ h(x) = g(f(x)).$ The domain of h is
(A) $( -2,2)$
(B) $[-2,2)$
(C) $[-2,2]$
(D) $(-\infty,\infty)$ \hfill \textbf{((C))}
\end{enumerate}
\section{Domain Of Functions}
\subsection{Medium}
\begin{enumerate}
\item Q. Find the domain of the function
$\quad y = \log_{x-4}(x^2 - 11x + 24).$
(A) $(3,4) \cup (8, \infty)$
(B) $(8,\infty)$
(C) $(4,5)$
(D) $(-\infty,3) \cup (8,\infty)$ \hfill \textbf{((B))}
\item Q. Find the domain of the function
$\quad f(x) = \log_2(-\log_{1/2}(1+\frac{1}{\sqrt{x}}) - 1).$
(A) $(0,1)$
(B) $(1,\infty)$
(C) $(0,\infty)$
(D) $(-\infty,0)$ \hfill \textbf{((A))}
\item Q. Find the domain of the function
$\quad f(x) = \frac{1}{\sqrt{[[|x|-5]] - 11}},$
where $[\cdot]$ denotes the greatest integer function.
(A) $(-\infty,-11] \cup [11, \infty)$
(B) $(-\infty, -16] \cup [16, \infty)$
(C) $(-\infty, -17] \cup [17, \infty)$
(D) $(-\infty, -6] \cup [6, \infty)$ \hfill \textbf{((C))}
\end{enumerate}
\section{Equations With Floor And Fractional Part Functions}
\subsection{Medium}
\begin{enumerate}
\item Q. Solve for x in the equation
$\quad |2x-1| = 3[x] + 2\{x\},$
where $[x]$ is the greatest integer function and $ \{x\} = x-[x]$ the fractional part. The solution
is
(A) $\quad x=\frac{1}{4}$
(B) $\quad x=\frac{1}{2}$
(C) $\quad x=1$
(D) no real solution \hfill \textbf{((A))}
\end{enumerate}
\section{Exponential And Logarithmic Equations}
\subsection{Medium}
\begin{enumerate}
\item The number of real solutions of the equation $\frac{2009^x}{2010} = 2009^{\log_x (2010)}$ can be expressed in lowest terms as $\frac{m}{n}$. What is $m - n$? \hfill \textbf{(B)}
\end{enumerate}
\section{Functional Equations}
\subsection{Medium}
\begin{enumerate}
\item Q. Let $f$ be a function on the natural numbers satisfying
$f(x + y) = f(x) f(y)$ for all $x, y \in N, f(1) = 2$.
Find the natural number $a$ such that for all $n \geq 1$,
$\sum_{k=1}^{n} f(a+k) = 16(2^n – 1)$.
(A) 1
(B) 2
(C) 3
(D) 4 \hfill \textbf{(C)}
\end{enumerate}
\section{Functions}
\subsection{Medium}
\begin{enumerate}
\item Q. f is an even function and g is an odd function. If $f(x) + g(x) = 2^{x+1}$, then
(A) $f(x) < 0$ for all x
(B) f is a constant function
(C) $0 < f(x) < 1$ for all x
(D) $f(x) \ge 2$ for all x \hfill \textbf{((D))}
\item Q. The number of real solutions of $2 \sin (e^x) = 5^x + 5^{-x}$ in $[0, 1]$ is /are
(A) 0
(B) 1
(C) 2
(D) 4 \hfill \textbf{((A))}
\item Q. Equation $|x - 4| + |x + 4| = ax + 8$ has \hfill \textbf{((C))}
\item Q. If $f(x) = \frac{1}{1-x}$, then the point of discontinuity of the function $f[f(f(x))]$ is/are
(A) $\{0,-1\}$
(B) $\{0,1\}$
(C) $\{1,-1\}$
(D) none of these \hfill \textbf{((B))}
\item Q. If $\alpha$ and $\beta$ are real numbers and $f(x) = \alpha \sin x + \beta \sqrt{x} + 4, x \in \mathbb{R}$, and if $f(\log_{10}(\log_3 10)) = 5$, then the value of $f(\log_{10}(\log_{10} 3))$ is \hfill \textbf{((A))}
\end{enumerate}
\section{Functions (Inverse Functions)}
\subsection{Medium}
\begin{enumerate}
\item Q. If $f : [1,\infty) \to [2, \infty)$ is given by $f(x) = x + \frac{1}{x}$, then $f^{-1}(x)$ equals
(A) $\frac{x+\sqrt{x^2-4}}{2}$
(B) $\frac{1+x^2}{x}$
(C) $\frac{x-\sqrt{x^2-4}}{2}$
(D) $1 + \sqrt{x^2-4}$ \hfill \textbf{((A))}
\end{enumerate}
\section{Geometric Progressions}
\subsection{Medium}
\begin{enumerate}
\item Q. If $a, b, c$ are in geometric progression, then the quadratic equations
$ax^2 + 2bx + c = 0$ and $dx^2 + 2ex + f = 0$
have a common root precisely when the numbers $\frac{d}{a}, \frac{e}{b}, \frac{f}{c}$ are in
(A) arithmetic progression (A.P.)
(B) geometric progression (G.P.)
(C) harmonic progression (H.P.)
(D) none of these \hfill \textbf{(A)}
\item Q. A three-digit number has digits in geometric progression. The sum of its right-hand and left-hand digits exceeds twice the middle digit by 1, and the sum of the left-hand and middle digits is two-thirds of the sum of the middle and right-hand digits. The sum of the digits of the number is
(A) 18
(B) 19
(C) 21
(D) 17 \hfill \textbf{(B)}
\end{enumerate}
\section{Geometry}
\subsection{Medium}
\begin{enumerate}
\item Q. In $\triangle ABC$, let $D$ be the midpoint of $BC$, and suppose
$\frac{\cot \angle CAD}{\cot \angle BAD} = \frac{2}{1}$
Let $G$ be the centroid of the triangle. What is the measure of $\angle BGA$? \hfill \textbf{((A))}
\item Q. On a semicircle with diameter $AB$, points $C$ and $D$ satisfy $AC = 1$, $CD = 2$, and $DB = 3$. If the length of $AB$ is $x$, then $x$ lies in which interval? \hfill \textbf{((B))}
\item Q. Among all the parallelograms whose diagonals are 10 and 4, the one having maximum area has its perimeter lying in the interval \hfill \textbf{((C))}
\end{enumerate}
\section{Inequalities}
\subsection{Medium}
\begin{enumerate}
\item Q. If $a, b, x, y$ are positive integers satisfying
$\frac{1}{x} + \frac{1}{y} = 1$,
then which of the following inequalities always holds?
(A) $\frac{a^x}{x} + \frac{b^y}{y} \geq ab$
(B) $\frac{a^x}{x} + \frac{b^y}{y} \leq ab$
(C) $\frac{a^x}{x} + \frac{b^y}{y} \geq a+b$
(D) $\frac{a^x}{x} + \frac{b^y}{y} \leq a+b$ \hfill \textbf{(A)}
\end{enumerate}
\subsection{Hard}
\begin{enumerate}
\item Q. The number of solutions of the following inequality
$(2^{1/\sin^2x_2}) \cdot (3^{1/\sin^2x_3})\cdot (4^{1/\sin^2x_4}) ...... (n^{1/\sin^2x_n}) \le n!$
where $x_i \in (0, 2\pi)$ for $i = 2, 3, . . . . . . n$ is
(A) $1$
(B) $2^{n-1}$
(C) $n^n$
(D) infinite number of solution \hfill \textbf{((B))}
\item Q. Suppose the cubic equation
$2x^3 + ax^2 + bx + 4 = 0$,
with $a, b > 0$, has three real roots. Then which of the following inequalities holds?
(A) $a+b \geq 6(2^{1/3} + 4^{1/3})$
(B) $a+b \geq 3(2^{1/3} + 4^{1/3})$
(C) $a+b \geq 6(2^{2/3} + 4^{2/3})$
(D) $a+b \geq 6(2^{1/3} + 4^{2/3})$ \hfill \textbf{(A)}
\end{enumerate}
\section{Inequalities And Logarithms}
\subsection{Medium}
\begin{enumerate}
\item The solution set of the inequality $\frac{(3^x - 4^x) \ln(x + 2)}{x^2 - 3x - 4} \leq 0$ is \hfill \textbf{(D)}
\end{enumerate}
\section{Inequalities And Square Roots}
\subsection{Medium}
\begin{enumerate}
\item Sum of all the real solutions of the inequality $\frac{(x^2 + 2) \sqrt{x^2 - 16}}{(x^2 + 2) (x^2 - 9)} \leq 0$ is \hfill \textbf{(D)}
\end{enumerate}
\section{Integration}
\subsection{Medium}
\begin{enumerate}
\item Q. $\int \frac{2x^{12}+5x^9}{(1+x^3+x^5)^3} dx$ equals
(A) $\frac{x^{10}}{2(x^5+x^3+1)^2} + C$
(B) $\frac{x^2+2x}{(x^5+x^3+1)} + C$
(C) $\log(x^5+x^3+1 + \sqrt{2x^{12}+5x^9}) + C$
(D) None of the above \hfill \textbf{((A))}
\end{enumerate}
\section{Inverse Functions}
\subsection{Medium}
\begin{enumerate}
\item Q. Let
$\quad f(x) = \begin{cases} x, & x < 1, \\ x^2, & 1 \le x \le 4, \\ 8\sqrt{x}, & x > 4. \end{cases}$
Find $\ f^{-1}(x).$
(A) $\quad f^{-1}(x) = \begin{cases} x, & x < 1, \\ \sqrt{x}, & 1 \le x \le 16, \\ \frac{x^2}{64}, & x > 16. \end{cases}$
(B) $\quad f^{-1}(x) = \begin{cases} x, & x < 1, \\ \sqrt{x}, & 1 \le x \le 4, \\ \frac{x^2}{64}, & x > 16. \end{cases}$
(C) $\quad f^{-1}(x) = \begin{cases} x, & x < 1, \\ \sqrt{x}, & 1 \le x \le 16, \\ \frac{x^2}{8}, & x > 16. \end{cases}$
(D) $\quad f^{-1}(x) = \begin{cases} x, & x < 1, \\ x^2, & 1 \le x \le 16, \\ \frac{x^2}{64}, & x > 16. \end{cases}$ \hfill \textbf{((A))}
\end{enumerate}
\section{Limits}
\subsection{Medium}
\begin{enumerate}
\item Q. If $\lim_{x\to\infty} (\sqrt{x^4 + ax^3 + 3x^2 + bx + 2} - \sqrt{x^4 + 2x^3 – cx^2 + 3x – d})$ is finite, then the value
of $a$ is
(A) 3
(B) 5
(C) 2
(D) any real number \hfill \textbf{((C))}
\item Q. $\lim_{x \to 2} \frac{2^x + 2^{3-x} - 6}{2^{x/2} - 2^{1-x}}$ is equal to
(A) 8
(B) 4
(C) 2
(D) none of these \hfill \textbf{((A))}
\item Q. Evaluate
$\lim_{x \to 1} \frac{x^{P+1} - (P+1)x + P}{(x-1)^2}$
where P is a constant. The value of the limit is
(A) $\frac{P(P+1)}{2}$
(B) P
(C) $\frac{(P-1)P}{2}$
(D) $\frac{(P+1)^2}{2}$ \hfill \textbf{((A))}
\item Q. Evaluate
$\lim_{x \to 1} \frac{\sqrt{x^2+8} - \sqrt{10-x^2}}{\sqrt{x^2+3} - \sqrt{5-x^2}}$
This is of the indeterminate form $\frac{0}{0}$. The value of the limit is
(A) $\frac{1}{2}$
(B) $\frac{1}{3}$
(C) $\frac{2}{3}$
(D) $\frac{3}{2}$ \hfill \textbf{((C))}
\item Q. If
$\lim_{x \to \infty} (\frac{x^3+1}{x^2+1} - (ax+b)) = 2$,
then
(A) a = 1, b = 1
(B) a = 1, b = 2
(C) a = 1, b = -2
(D) none of these \hfill \textbf{((C))}
\item Q. Evaluate the limit
$\lim_{n \to \infty} \frac{[x] + [2x] + [3x] + \cdots + [nx]}{n^2}$
where $[.]$ denotes the greatest-integer function. The value of the limit is
(A) $\frac{x}{2}$
(B) x
(C) $\frac{x+1}{2}$
(D) $\frac{x-1}{2}$ \hfill \textbf{((A))}
\item Q. Evaluate
$\lim_{n \to \infty} \frac{\sin(\frac{a}{n})}{\tan(\frac{b}{n+1})}$
where a, b are constants. The value of the limit is
(A) $\frac{a}{b}$
(B) $\frac{b}{a}$
(C) 1
(D) 0 \hfill \textbf{((A))}
\item Q. Evaluate
$\lim_{x \to \frac{\pi}{2}} (\sec x - \tan x)$.
This limit is of the $\infty - \infty$ form. The value of the limit is
(A) 0
(B) 1
(C) -1
(D) does not exist \hfill \textbf{((A))}
\item Q. Evaluate
$\lim_{x \to 0} (\frac{a^x + b^x + c^x}{3})^{1/x}$
where a, b, c > 0. The value of this limit is
(A) abc
(B) $\sqrt{abc}$
(C) $(abc)^{1/3}$
(D) none of these \hfill \textbf{((C))}
\item Q. Evaluate
$\lim_{n \to \infty} n^{-n^2} ((n+1)(n+\frac{1}{2})(n+\frac{1}{2^2}) \cdots (n+\frac{1}{2^{n-1}}))^n$.
(A) $e^2$
(B) e
(C) $e^3$
(D) none of these \hfill \textbf{((A))}
\end{enumerate}
\section{Linear Algebra}
\subsection{Easy}
\begin{enumerate}
\item Q. Let
$A_{\alpha} = \begin{pmatrix} \cos \alpha & \sin \alpha \\ - \sin \alpha & \cos \alpha \end{pmatrix}$.
Compute $A_{\alpha} A_{\beta}$. Which of the following is correct?
(A) $A_{\alpha+\beta}$
(B) $A_{\alpha\beta}$
(C) $A_{\alpha-\beta}$
(D) none of these \hfill \textbf{((A))}
\item Q. If
$A=\begin{pmatrix} 1 & \tan x \\ -\tan x & 1 \end{pmatrix}$,
and the function $f(x) = \det(A^T A^{-1})$, (assuming the question intends $f(x) = \det(A^T A^{-1})$ as per solution)
then which of the following can be the value of
$f(f(f(\dots f(x))))$ ($n \ge 2$)?

(A) $f^n(x)$
(B) 1
(C) $f^{n-1}(x)$
(D) $n f(x)$ \hfill \textbf{((B))}
\end{enumerate}
\subsection{Medium}
\begin{enumerate}
\item Q. For the determinant
$D=\begin{vmatrix} 2 & -3 & 1 \\ 4 & 0 & 5 \\ -1 & 6 & 7 \end{vmatrix}$,
find the minors and cofactors of the elements -3, 5, -1, and 7. Which of the following sets is correct?
(A) Minor(-3) = 33, Cofactor(-3) = -33; Minor(5) = 9, Cofactor(5) = -9; Minor(-1) = -15, Cofactor(-1) = 15; Minor(7) = 12, Cofactor(7) = 12
(B) Minor(-3) = 33, Cofactor(-3) = -33; Minor(5) = 9, Cofactor(5) = -9; Minor(-1) = -15, Cofactor(-1) = -15; Minor(7) = 12, Cofactor(7) = 12
(C) Minor(-3) = $-33$, Cofactor(-3) = 33; Minor(5) = 9, Cofactor(5) = -9; Minor(-1) = -15, Cofactor(-1) = $-15$; Minor(7) = 12, Cofactor(7) = 12
(D) Minor(-3) = 33, Cofactor(-3) = -33; Minor(5) = 9, Cofactor(5) = $-9$; Minor(-1) = -15, Cofactor(-1) = -15; Minor(7) = 12, Cofactor(7) = $-12$ \hfill \textbf{((B))}
\item Q. Let
$D_r = \begin{vmatrix} r & r^3 & 2 \\ n & n^3 & 2n \\ \frac{n(n+1)}{2} & (\frac{n(n+1)}{2})^2 & 2(n+1) \end{vmatrix}$.
Find $\sum_{r=0}^{n} D_r$.
(A) 0
(B) $\frac{n(n+1)}{2}$
(C) $n(n + 1)$
(D) $(\frac{n(n+1)}{2})^2$ \hfill \textbf{((A))}
\item Q. Evaluate the determinant
$D = \begin{vmatrix} a+b+nc & (n-1)a & (n-1)b \\ (n-1)c & b+c+na & (n-1)b \\ (n-1)c & (n-1)a & c+a+nb \end{vmatrix}$.
The value of D is
(A) $(a + b + c)^3$
(B) $n (a + b + c)^3$
(C) $(n - 1) (a + b + c)^3$
(D) none of these \hfill \textbf{((B))}
\item Q. Evaluate the cyclic determinant
$D = \begin{vmatrix} 1 & \alpha & \alpha^2 \\ \alpha & \alpha^2 & 1 \\ \alpha^2 & 1 & \alpha \end{vmatrix}$.
Which of the following is equal to D?
(A) $-(1- \alpha^3)^2$
(B) $(1 - \alpha^3)^2$
(C) $-(1 + \alpha^3)^2$
(D) $(\alpha - 1)^2 (1 + \alpha + \alpha^2)^2$ \hfill \textbf{((A))}
\item Q. If $x, y, z$ (not all zero) satisfy the system
$\begin{cases} \sin(3\theta) x - y + z = 0, \\ \cos(2\theta) x + 4y + 3z = 0, \\ 2x + 7y + 7z = 0, \end{cases}$
then the possible values of $\theta$ in $[0, 2\pi]$ are
(A) $0, \frac{\pi}{6}, \frac{5\pi}{6}, \pi, 2\pi$
(B) $\frac{\pi}{6}, \frac{5\pi}{6}$
(C) $0, \pi, 2\pi$
(D) $\frac{\pi}{2}, \frac{3\pi}{2}$ \hfill \textbf{((A))}
\item Q. For which values of $\lambda$ is the following system of equations consistent (i.e., has a nontrivial solution)?
$\begin{cases} x+y-3= 0, \\ (1 + \lambda) x + (2 + \lambda) y - 8 = 0, \\ x - (1 + \lambda) y + (2 + \lambda) = 0? \end{cases}$
(A) $\lambda = 1$
(B) $\lambda = - \frac{5}{3}$
(C) $\lambda = 1, - \frac{5}{3}$
(D) $\lambda = -1, \frac{5}{3}$ \hfill \textbf{((C))}
\item Q. If the product of $n$ matrices
$\begin{pmatrix} 1 & 1 \\ 0 & 1 \end{pmatrix} \begin{pmatrix} 1 & 2 \\ 0 & 1 \end{pmatrix} \begin{pmatrix} 1 & 3 \\ 0 & 1 \end{pmatrix} \dots \begin{pmatrix} 1 & n \\ 0 & 1 \end{pmatrix}$
is equal to
$\begin{pmatrix} 1 & 378 \\ 0 & 1 \end{pmatrix}$,
then the value of $n$ is
(A) 26
(B) 27
(C) 377
(D) 378 \hfill \textbf{((B))}
\item Q. If
$A = \begin{pmatrix} 0 & 1 \\ -1 & 0 \end{pmatrix}$
and $(aI_2+bA)^2 = A$,
then
(A) $a = b = \sqrt{2}$
(B) $a = b = \frac{1}{\sqrt{2}}$
(C) $a = b = \sqrt{3}$
(D) $a = b = \frac{1}{\sqrt{3}}$ \hfill \textbf{((B))}
\item Q. Let
$A=\begin{pmatrix} 2 & 1 \\ 4 & 1 \end{pmatrix}, B=\begin{pmatrix} 3 & 4 \\ 2 & 3 \end{pmatrix}, C=\begin{pmatrix} 3 & -4 \\ -2 & 3 \end{pmatrix}$.
Evaluate the infinite sum
$\operatorname{tr}(A) + \operatorname{tr}(\frac{A(BC)}{2}) + \operatorname{tr}(\frac{A(BC)^2}{4}) + \operatorname{tr}(\frac{A(BC)^3}{8}) + \dots$
(A) 6
(B) 9
(C) 12
(D) none of these \hfill \textbf{((A))}
\item Q. Let A and B be distinct $n \times n$ matrices satisfying
$A^3 = B^3$ and $A^2B = B^2A$.
Determine which of the following statements is correct:
(A) $\det(A^2 + B^2) = 0$
(B) $\det(A - B) = 0$
(C) Both $\det(A^2 + B^2)$ and $\det(A - B)$ must be zero
(D) At least one of $\det(A^2 + B^2)$ or $\det(A - B)$ must be zero \hfill \textbf{((D))}
\item Q. Let
$A = \begin{pmatrix} x & 3 & 2 \\ 1 & y & 4 \\ 2 & 2 & z \end{pmatrix}$,
with $xyz = 60$ and $8x + 4y + 3z = 20$. Then $A (\operatorname{adj} A)$ equals
(A) $\begin{pmatrix} 64 & 0 & 0 \\ 0 & 64 & 0 \\ 0 & 0 & 64 \end{pmatrix}$
(B) $\begin{pmatrix} 88 & 0 & 0 \\ 0 & 88 & 0 \\ 0 & 0 & 88 \end{pmatrix}$
(C) $\begin{pmatrix} 68 & 0 & 0 \\ 0 & 68 & 0 \\ 0 & 0 & 68 \end{pmatrix}$
(D) $\begin{pmatrix} 34 & 0 & 0 \\ 0 & 34 & 0 \\ 0 & 0 & 34 \end{pmatrix}$ \hfill \textbf{((C))}
\item Q. Let
$A = \left\{\mathbf{a} = (a_1, a_2, a_3) \in \mathbb{R}^3 : \left(\frac{a_1}{2} + \frac{a_2}{4} + \frac{a_3}{8}\right)^2 = \frac{a_1^2}{4} + \frac{a_2^2}{16} + \frac{a_3^2}{8}\right\}$.
Then how many elements does $A$ have? \hfill \textbf{((B))}
\end{enumerate}
\subsection{Hard}
\begin{enumerate}
\item Q. Evaluate the determinant
$D = \begin{vmatrix} (x - a)^2 & (x - b)^2 & (x - c)^2 \\ (y - a)^2 & (y - b)^2 & (y - c)^2 \\ (z - a)^2 & (z - b)^2 & (z - c)^2 \end{vmatrix}$.
Which of the following is equal to D?
(A) $(x - y)(y - z)(z - x)(a - b)(b - c)(c - a)$
(B) $2 (x - y)(y - z)(z - x)(a - b)(b - c)(c - a)$
(C) $\frac{1}{2} (x - y)(y - z)(z - x)(a - b)(b - c)(c - a)$
(D) $(x - y)^2(y - z)^2(z - x)^2$ \hfill \textbf{((B))}
\item Q. Let $\alpha$ and $\beta$ be the roots of the equation $ax^2 + bx + c = 0$, and define $S_n = \alpha^n + \beta^n$ for $n \ge 1$. Evaluate the determinant
$D = \begin{vmatrix} 3 & 1+S_1 & 1+S_2 \\ 1+S_1 & 1+S_2 & 1+S_3 \\ 1+S_2 & 1+S_3 & 1+S_4 \end{vmatrix}$.
Which of the following is equal to D?
(A) $\frac{(b^2 - 4ac) (a + b + c)^2}{a^4}$
(B) $\frac{b^2-4ac}{a^2}$
(C) $\frac{(b^2 - 4ac) (a + b + c)}{a^3}$
(D) $\frac{(b^2 - 4ac) (a + b + c)^2}{a^3}$ \hfill \textbf{((A))}
\end{enumerate}
\section{Linear Algebra - System Of Equations}
\subsection{Medium}
\begin{enumerate}
\item Q. If the system of linear equations
$\begin{cases} 2x + 2y + 3z = a, \\ 3x - y + 5z = b, \\ x - 3y + 2z = c, \end{cases}$
where a, b, c are non-zero real numbers, has more than one solution, then:
(A) $a + b + c = 0$
(B) $b - c - a = 0$
(C) $b + c - a = 0$
(D) $b - c + a = 0$ \hfill \textbf{((B))}
\end{enumerate}
\section{Logarithm Properties}
\subsection{Medium}
\begin{enumerate}
\item Evaluate the expression $\frac{1}{\log_a(abc)} + \frac{1}{\log_b(abc)} + \frac{1}{\log_c(abc)}$ \hfill \textbf{(B)}
\end{enumerate}
\section{Logarithmic Equations}
\subsection{Medium}
\begin{enumerate}
\item If $\log_{(2x+3)} (6x^2 + 23x +21) = 4 - \log_{(3x+7)} (4x^2 + 12x + 9)$, then the value of $x$ is equal to \hfill \textbf{(C)}
\end{enumerate}
\section{Logarithmic Inequalities}
\subsection{Medium}
\begin{enumerate}
\item The complete solution set of the inequality $\frac{1}{\log_4(\frac{x+1}{x+2})} < \frac{1}{\log_4(x + 3)}$ can be written in the form $(-a, \infty)$. Determine $a$. \hfill \textbf{(A)}
\end{enumerate}
\section{Logarithms And Exponents}
\subsection{Medium}
\begin{enumerate}
\item If $a, b, c$ are positive real numbers such that $a^{\log_3 7} = 27$, $b^{\log_7 11} = 49$, $c^{\log_{11} 25} = \sqrt{11}$, then the value of $a^{(\log_3 7)^2} + b^{(\log_7 11)^2} + c^{(\log_{11} 25)^2}$ equals \hfill \textbf{(B)}
\end{enumerate}
\section{Mathematics}
\subsection{Medium}
\begin{enumerate}
\item Q. The sums of the first n terms of two arithmetic progressions are in the ratio
$\frac{S_{1,n}}{S_{2,n}} = \frac{7n+1}{4n+27}$
Find the ratio of their $11^{th}$ terms $t_{1,11} : t_{2,11}$.
(A) $\frac{4}{3}$
(B) $\frac{148}{111}$
(C) $\frac{111}{148}$
(D) $\frac{7}{4}$ \hfill \textbf{((A))}
\item Q. Four distinct integers form an increasing arithmetic progression. One of these four num-bers equals the sum of the squares of the other three. The four numbers are
(A) -2, -1, 0, 1
(B) −1, 0, 1, 2
(C) 0, 1, 2, 3
(D) -2, 0, 2, 4 \hfill \textbf{((B))}
\item Q. Let a, b, c, d, p be distinct real numbers satisfying
$(a^2 + b^2 + c^2)p^2 - 2p (ab+bc+cd) + (b^2+c^2 + d^2) \le 0$.
Then a, b, c, d must be in
(A) Arithmetic progression (A.P.)
(B) Geometric progression (G.P.)
(C) Harmonic progression (H.P.)
(D) None of these \hfill \textbf{((B))}
\item Q. If $a_1, a_2, ..., a_n$ are in arithmetic progression with $a_i > 0$ for all i, then the sum
$\sum_{k=1}^{n-1} \frac{1}{\sqrt{a_k} + \sqrt{a_{k+1}}}$
equals
(A) $\frac{n-1}{\sqrt{a_1}+\sqrt{a_n}}$
(B) $\frac{n}{\sqrt{a_1}+\sqrt{a_n}}$
(C) $\frac{n-1}{\sqrt{a_n} - \sqrt{a_1}}$
(D) $\frac{n+1}{\sqrt{a_1} + \sqrt{a_n}}$ \hfill \textbf{((A))}
\item Q. If a, b, c are in geometric progression, then the quadratic equations
$ax^2 + 2bx + c = 0$ and $dx^2 + 2ex + f = 0$
have a common root precisely when the numbers $\frac{d}{a}, \frac{e}{b}, \frac{f}{c}$ are in
(A) arithmetic progression (A.P.)
(B) geometric progression (G.P.)
(C) harmonic progression (H.P.)
(D) none of these \hfill \textbf{((A))}
\item Q. A three-digit number has digits in geometric progression. The sum of its right-hand and left-hand digits exceeds twice the middle digit by 1, and the sum of the left-hand and middle digits is two-thirds of the sum of the middle and right-hand digits. The sum of the digits of the number is
(A) 18
(B) 19
(C) 21
(D) 17 \hfill \textbf{((B))}
\item Q. Let a be a fixed real number and x, y, z satisfy
$\frac{a-x}{px} = \frac{a-y}{qy} = \frac{a-z}{rz}$
If p, q, r are in arithmetic progression, then x, y, z are in
(A) arithmetic progression (A.P.)
(B) geometric progression (G.P.)
(C) harmonic progression (H.P.)
(D) none of these \hfill \textbf{((C))}
\item Q. Let $z_1, z_2, z_3$ be the complex affixes of the vertices A, B, C of an isosceles right triangle with right angle at C. Which of the following identities holds?
(A) $(z_1-z_2)^2 = 2(z_1-z_3) (z_3-z_2)$
(B) $(z_1-z_2)^2 = -2(z_1-z_3) (z_3-z_2)$
(C) $(z_1-z_3)^2 = 2(z_1-z_2) (z_2-z_3)$
(D) $(z_2-z_3)^2 = 2 (z_2-z_1) (z_1-z_3)$ \hfill \textbf{((A))}
\item Q. Suppose the cubic equation
$2x^3 + ax^2 + bx + 4 = 0$,
with a, b > 0, has three real roots. Then which of the following inequalities holds?
(A) $a+b \ge 6(2^{1/3} + 4^{1/3})$
(B) $a+b \ge 3(2^{1/3} + 4^{1/3})$
(C) $a+b \ge 6(2^{2/3} + 4^{2/3})$
(D) $a+b \ge 6(2^{1/3} + 4^{2/3})$ \hfill \textbf{((A))}
\item Q. If a, b, x, y are positive integers satisfying
$\frac{1}{x} + \frac{1}{y} = 1$,
then which of the following inequalities always holds?
(A) $\frac{a^x}{x} + \frac{b^y}{y} \ge ab$
(B) $\frac{a^x}{x} + \frac{b^y}{y} < ab$
(C) $\frac{a^x}{x} + \frac{b^y}{y} \ge a+b$
(D) $\frac{a^x}{x} + \frac{b^y}{y} < a+b$ \hfill \textbf{((A))}
\item Q. Find the sum of the infinite series
$4 - 9x + 16x^2 - 25x^3 + 36x^4 - 49x^5 + ...$
for $|x| < 1$. The sum S equals
(A) $\frac{4 + 3x + x^2}{(1-x)^3}$
(B) $\frac{4 + 3x + x^2}{(1+x)^3}$
(C) $\frac{4- 3x + x^2}{(1+x)^3}$
(D) $\frac{4+ 3x - x^2}{(1+x)^3}$ \hfill \textbf{((B))}
\item Q. Find the sum of the first 16 terms of the series
$\frac{1^3}{1} + \frac{1^3 + 2^3}{1+3} + \frac{1^3 + 2^3 + 3^3}{1+3+5} + ...$
The value is
(A) 450
(B) 456
(C) 446
(D) none of these \hfill \textbf{((C))}
\item Q. Let $T_r$ be a sequence whose partial sums satisfy
$\sum_{r=1}^{n} T_r = \frac{n(n+1)(n+2)(n+3)}{8}$.
Compute the infinite series
$\sum_{r=1}^{\infty} \frac{1}{T_r}$.
(A) $\frac{1}{2}$
(B) $\frac{1}{3}$
(C) $\frac{2}{3}$
(D) $\frac{3}{4}$ \hfill \textbf{((A))}
\item Q. The positive integers are partitioned into consecutive “groups” as follows:
(1) , (2,3,4), (5,6,7,8,9), ...
1st group 2nd group 3rd group
where the $n^{th}$ group contains $2n - 1$ numbers. What is the sum of the numbers in the $n^{th}$ group?
(A) $n^3 + (n - 1)^3$
(B) $n^3 + (n + 1)^3$
(C) $2n^3 - n$
(D) $(2n - 1) n^2$ \hfill \textbf{((A))}
\item Q. Let f be a function on the natural numbers satisfying
$f(x + y) = f(x) f(y)$ for all $x, y \in \mathbb{N}$, $f(1) = 2$.
Find the natural number a such that for all $n \ge 1$,
$\sum_{k=1}^{n} f(a+k) = 16(2^n – 1)$.
(A) 1
(B) 2
(C) 3
(D) 4 \hfill \textbf{((C))}
\item A triangle has base AB = 10 cm and base angles $\angle A = 50^\circ$, $\angle B = 70^\circ$. If its perimeter can be written as Perimeter = $x + y \cos z^\circ$, where $z \in (0,90^\circ)$, then the value of $x + y + z$ is
(A) 60
(B) 55
(C) 50
(D) 40 \hfill \textbf{((D))}
\item Over the reals one can factorize $z^5 - 32 = (z - 2) (z^2 - pz + 4) (z^2 - qz + 4)$. If this holds, then the quantity $p^2 + 2q$ equals
(A) 2
(B) 3
(C) 4
(D) 8 \hfill \textbf{((C))}
\item Let $g: (0,\infty) \to R$ be twice-differentiable and satisfy $g''(x) - 3g'(x) > 3$, $x \ge 0$, and $g'(0) = -1$. Then for $x > 0$, the function $g(x) + x$ is
(A) an increasing function
(B) a decreasing function
(C) a constant function
(D) none of these \hfill \textbf{((A))}
\item How many integers between 1 and 10000 (inclusive) have at least one digit equal to 8 and at least one digit equal to 9?
(A) 972
(B) 974
(C) 970
(D) 980 \hfill \textbf{((B))}
\item The line $2x - y + 1 = 0$ is tangent to a circle at the point $A(2,5)$. The centre $O(h, k)$ of the circle lies on the line $x - 2y = 4$. The radius of the circle is
(A) $3\sqrt{5}$
(B) $5\sqrt{3}$
(C) $2\sqrt{5}$
(D) $5\sqrt{2}$ \hfill \textbf{((A))}
\item Let $X = \{x \in R : \cos(\sin x) = \sin(\cos x)\}$. How many elements does X contain?
(A) 0
(B) 2
(C) 4
(D) Not finite \hfill \textbf{((A))}
\item There are six boxes labeled $B_1, B_2, ..., B_6$. In each of $n$ trials, two fair dice $D_1, D_2$ are thrown. If $D_1$ shows $j$ and $D_2$ shows $k$, then $j$ balls are placed into box $B_k$. What is the probability that box $B_1$ contains at most one ball after $n$ trials?
(A) $(\frac{5}{6})^n + n (\frac{5}{6})^{n-1} \frac{1}{6}$
(B) $(\frac{5}{6})^n + n (\frac{5}{6})^{n-1} \frac{1}{6}$
(C) $(\frac{5}{6})^n + n (\frac{5}{6})^{n-1} \frac{1}{6^1}$
(D) $(\frac{5}{6})^n + n (\frac{5}{6})^{n-1} \frac{1}{6^2}$ \hfill \textbf{((D))}
\item The minimum distance between a point on the curve $y = e^x$ and a point on the curve $y = \ln x$ is
(A) $\frac{1}{\sqrt{2}}$
(B) $\sqrt{2}$
(C) $\sqrt{3}$
(D) $2\sqrt{2}$ \hfill \textbf{((B))}
\item Let $[x]$ and $\{x\}$ denote the integer and fractional parts of a real number $x$, respectively. Evaluate the integral $\int_{0}^{5} [x] \{x\}dx$.
(A) $\frac{5}{2}$
(B) 5
(C) $\frac{69}{2}$
(D) $\frac{71}{2}$ \hfill \textbf{((B))}
\item There are three distinct parallel lines in a plane. On each line $p$ points are chosen (so in total $3p$ points). What is the maximum number of triangles whose vertices are chosen from these points?
(A) $3p^2(p-1) + 1$
(B) $3 p^2(p-1)$
(C) $p^2(4p - 3)$
(D) $p^2(2p - 3)$ \hfill \textbf{((C))}
\item Let A, B, C be $n \times n$ matrices with $\det(A) = 2$, $\det(B) = 3$, and $\det(C) = 5$. Compute $\det (A^2 BC^{-1})$.
(A) $\frac{6}{5}$
(B) $\frac{12}{5}$
(C) $\frac{18}{5}$
(D) $\frac{24}{5}$ \hfill \textbf{((B))}
\item Let $f(x) = x^3 + \alpha x^2 + \beta x + \gamma$, where $\alpha, \beta, \gamma$ are rational numbers. Suppose two of the roots of $f(x) = 0$ are the eccentricities of a parabola and of a rectangular hyperbola. Then $\alpha + \beta + \gamma$ equals
(A) -1
(B) 0
(C) 1
(D) 2 \hfill \textbf{((A))}
\item For the hyperbola $\frac{x^2}{a^2} - \frac{y^2}{b^2} = 1$, the chords of contact of tangents drawn from the points (-4, 2) and (2, 1) are perpendicular. Then the eccentricity $e$ of the hyperbola is
(A) $\frac{\sqrt{7}}{2}$
(B) $\sqrt{\frac{5}{2}}$
(C) $\sqrt{\frac{3}{2}}$
(D) $\sqrt{2}$ \hfill \textbf{((C))}
\item If $\log_{245} 175 = a$ and $\log_{1715} 875 = b$, then the value of $\frac{1-ab}{a-b}$ is
(A) 2
(B) 3
(C) 5
(D) 7 \hfill \textbf{((C))}
\item Let Z be a nonzero complex number satisfying $Z^3 + Z^{-3} \le 2$. Then the maximum possible value of $| Z + Z^{-1}|$ is
(A) 2
(B) $3\sqrt{2}$
(C) $2\sqrt{2}$
(D) 1 \hfill \textbf{((A))}
\item Which of the following does not represent a hyperbola?
(A) $(x - 1)(y - 3) = 3$
(B) $xy = 1$
(C) $x^2 - y^2 = 0$
(D) $x^2 - y^2 = 5$ \hfill \textbf{((C))}
\item In a sports club of 50 members, 22 play cricket, 20 play tennis, and 18 play badminton. Every member plays at least one game, and no one plays all three. How many members play exactly two of these games?
(A) 18
(B) 32
(C) 10
(D) 8 \hfill \textbf{((C))}
\item Let $\alpha = a + i\beta$ with $a, \beta \in \mathbb{R}$. Consider the linear system
$$ \begin{cases} 4ix + (1+i) y = 0, \\ 8(\cos\frac{2\pi}{3}+i\sin\frac{2\pi}{3}) x + \alpha y = 0. \end{cases} $$
If this system has a nontrivial solution, then $\frac{a}{\beta} =$
(A) $2 - \sqrt{3}$
(B) $-2 + \sqrt{3}$
(C) $2 + \sqrt{3}$
(D) $-2 - \sqrt{3}$ \hfill \textbf{((A))}
\item Let $f$ and $g$ be twice-differentiable even functions on $(-2, 2)$ satisfying
$f(\frac{1}{2}) = 0$, $f(-\frac{1}{2}) = 0$, $f(1) = 1$, $g(\frac{1}{2}) = 0$, $g(1) = 2$.
Consider the equation
$f(x)g''(x) + f'(x) g'(x) = 0$,
for $x \in (-2,2)$. The minimum number of solutions of this equation in $(-2, 2)$ is
(A) 2
(B) 3
(C) 4
(D) 5 \hfill \textbf{((C))}
\item The direction ratios of two lines $L_1$ and $L_2$ are $(4, -1,3)$ and $(2, -1,2)$, respectively. A vector $\vec{V}$ is perpendicular to both $L_1$ and $L_2$ and satisfies $|\vec{V}| = 15$. If
$\vec{V} = x_1 \hat{i} + x_2 \hat{j} + x_3 \hat{k}$,
find $x_1 + x_2 + x_3$.
(A) 10
(B) 15
(C) 20
(D) 5 \hfill \textbf{((B))}
\item If the area bounded by the curve $2y^2 = 3x$, lines $x + y = 3$, $y = 0$ and outside the circle $(x - 3)^2 + y^2 = 2$ is $A$, then $4(\pi+ 4 A)$ is equal to
(A) 36
(B) 42
(C) 48
(D) 56 \hfill \textbf{((B))}
\item Different arithmetic progressions are formed with first term 100 and last term 199, each having an integral common difference. Among those progressions with between 3 and 33 terms (inclusive), the sum of all possible common differences is:
(A) 45
(B) 53
(C) 66
(D) 99 \hfill \textbf{((B))}
\item Let
$$ X=\begin{pmatrix} 0 & 1 & 0 \\ 0 & 0 & 1 \\ 0 & 0 & 0 \end{pmatrix}, \quad Y = \alpha I + \beta X + \gamma X^2, $$
with $\alpha, \beta, \gamma \in \mathbb{R}$. If
$$ Y^{-1} = \begin{pmatrix} 1 & -\frac{1}{5} & \frac{1}{5} \\ 0 & \frac{1}{5} & -\frac{2}{5} \\ 0 & 0 & \frac{1}{5} \end{pmatrix}, $$
then the value of $(\alpha - \beta + \gamma)^2$ is
(A) 64
(B) 100
(C) 144
(D) 36 \hfill \textbf{((B))}
\item A real-valued function $f$ satisfies the functional equation
$f(x - y) = f(x)f(y) - f(a - x) f(a + y)$,
where $a$ is a constant and $f(0) = 1$. Then $f(2a - x) =$
(A) $-f(x)$
(B) $f(x)$
(C) $f(a) + f(a - x)$
(D) $f(-x)$ \hfill \textbf{((A))}
\item What is the minimum distance between the circles
$x^2 + y^2 = 144$ and $x^2 + y^2 - 6x - 8y = 0$?
(A) 17
(B) 2
(C) 7
(D) 0 \hfill \textbf{((B))}
\item If $a_r$ is the coefficient of $x^{10-r}$ in the expansion of $(1 + x)^{10}$, evaluate
$$ \sum_{r=1}^{10} r^3 \left(\frac{a_r}{a_{r-1}}\right)^2. $$
(A) 5445
(B) 1210
(C) 3025
(D) 4895 \hfill \textbf{((B))}
\item What is the distance from the point $(1, 1)$ to the line
$2x - 3y = 4$
measured along the direction of the line
$x + y = 1$ ?
(A) $\sqrt{2}$
(B) $\frac{1}{\sqrt{2}}$
(C) $5\sqrt{2}$
(D) $2\sqrt{2}$ \hfill \textbf{((A))}
\item If
$u + iv = \frac{3i}{x + iy+2}$,
express $y$ in terms of $u$ and $v$.
(A) $\frac{5u}{u^2 + v^2}$
(B) $\frac{3u}{u^2 + v^2}$
(C) $\frac{4u}{u^2 + v^2}$
(D) $\frac{2u}{u^2 + v^2}$ \hfill \textbf{((B))}
\item Let $f(x)$ be a polynomial of degree four which has extreme values at $x = 1$ and $x = 2$. If
$$ \lim_{x\to 0}\left[1+\frac{f(x)}{x^2}\right] = 3, $$
then $f(2)$ equals
(A) 4
(B) -4
(C) -8
(D) 0 \hfill \textbf{((D))}
\item Find the particular solution of the differential equation
$\frac{dy}{dx} + \frac{y}{2} \sec x = \frac{\tan x}{2 y}$,
$0 < x < \frac{\pi}{2}$,
satisfying $y(0) = 1$. The answer may be written implicitly in terms of $y^2$. Which of the following is correct?
(A) $y^2 = 1+\frac{x}{\sec x + \tan x}$
(B) $y = 1-\frac{x}{\sec x + \tan x}$
(C) $y = 1+\frac{x}{\sec x + \tan x}$
(D) $y^2 = 1-\frac{x}{\sec x + \tan x}$ \hfill \textbf{((D))}
\item The set of all real numbers $x$ for which $x^2 - |x + 2| + x > 0$, is \hfill \textbf{((B))}
\item Number of positive integral solutions of $15 < x_1 + x_2 + x_3 \leq 20$ is \hfill \textbf{((D))}
\item If $x = (\sqrt{3} + 1)^n$, then $[x]$ is (where $[x]$ denotes the greatest integer less than or equal to $x$) \hfill \textbf{((A))}
\item Given that, $x + y + z = 15$, where $a, x, y, z$ and $b$ are in AP and $\frac{1}{x} + \frac{1}{y} + \frac{1}{z} = \frac{5}{3}$, where $a, x, y, z$ and $b$ are in HP. Then, \hfill \textbf{((A))}
\item A five digit number is chosen at random. The probability that all the digits are distinct and digits at odd places are odd and digits at even places are even, is \hfill \textbf{((D))}
\item A square matrix $A$ of order 3 satisfies $A^2 = I - 2A$, where $I$ is an identity matrix of order 3. If $A^n = 29A - 12I$, then the value of $n$ is equal to \hfill \textbf{((C))}
\item The roots of the equation $x^4 - 2x^2 + 4 = 0$ are the vertices of a \hfill \textbf{((D))}
\item The value of $\int e^{\sec x} \cdot \sec^3 x (\sin^2 x + \cos x + \sin x + \sin x \cos x) dx$ is: \hfill \textbf{((C))}
\item If $x$ and $y$ are positive real numbers then minimum value of $x^{100} + \frac{100}{x} + y^{50} + \frac{50}{y}$ is \hfill \textbf{((D))}
\item The function $f(x) = x^3 + ax^2 + bx + c$, $a^2 < 3b$ has: \hfill \textbf{((C))}
\item Let $a, b \in N, a \neq b$ and the two quadratic equations $(a - 1)x^2 - (a^2 + 2) x + a^2 + 2a = 0$ and $(b - 1)x^2 - (b^2 + 2) x + b^2 + 2b = 0$ have a common root. The value of $ab$ is \hfill \textbf{((C))}
\item A student is given a true - false exam with 10 questions. If he gets 8 or more correct answers he passes the exam. Given that he guesses at the answers to each question, what is the probability that he passes the exam. \hfill \textbf{((A))}
\item If $\sin^{-1} \sin(\sqrt{1 - \alpha}) = \sqrt{1 - \alpha}$, then \hfill \textbf{((D))}
\item The value of $(^{21}C_1 - ^{10}C_1)+(^{21}C_2 - ^{10}C_2)+(^{21}C_3 - ^{10}C_3)+(^{21}C_4 - ^{10}C_4)+...+(^{21}C_{10} - ^{10}C_{10})$ is \hfill \textbf{((D))}
\item If the distance of any point $(x, y)$ from the origin is defined as $d(x, y) = \max\{|x|, |y|\}$, $d(x, y) = a$, a non zero constant, then the locus is \hfill \textbf{((C))}
\end{enumerate}
\section{Matrices}
\subsection{Medium}
\begin{enumerate}
\item Q. Let $P = \begin{bmatrix} 1 & 0 & 0 \\ 4 & 1 & 0 \\ 16 & 4 & 1 \end{bmatrix}$ and $I$ be the identity matrix of order 3 . If $Q = [q_{ij}]$ is a matrix such that $P^{50} – Q = I$, then the value of $\frac{q_{31}+q_{32}}{q_{21}}$ is equal to \hfill \textbf{((B))}
\item Q. The matrix $A = \begin{pmatrix} 1 & 2 & 2 \\ 2 & 1 & 1 \\ 2 & 2 & 1 \end{pmatrix}$ is a zero divisor of the polynomial $f(x) = x^2 - 4x - 5$. If m is the trace of the matrix $A^3$, then $\left\lfloor\frac{m}{20}\right\rfloor$ is (where $[\cdot]$ denotes the greatest-integer function) \hfill \textbf{((A))}
\end{enumerate}
\section{Matrices And Determinants}
\subsection{Medium}
\begin{enumerate}
\item Q. If $A = \begin{pmatrix} 1 & 1 \\ 1 & 1 \end{pmatrix}$ and $\det (A^n - I) = 1 - \lambda^n, n \in \mathbb{N}$, then the value of $\lambda$ is \hfill \textbf{((B))}
\end{enumerate}
\section{Matrix Algebra}
\subsection{Medium}
\begin{enumerate}
\item Q. Let a matrix A = $\begin{pmatrix} 0 & 2q & r \ p & q & -r \ p & -q & r \end{pmatrix}$. If $AA^T = I_3$, then the value of $|p|$ equals:
(A) $\frac{1}{\sqrt{6}}$
(B) $\frac{1}{\sqrt{2}}$
(C) $\frac{1}{\sqrt{3}}$
(D) $\frac{1}{2}$ \hfill \textbf{((B))}
\end{enumerate}
\section{Matrix And Determinants}
\subsection{Medium}
\begin{enumerate}
\item Q. $\begin{vmatrix} 1 & 0 & 0 \\ 0 & 2 & 0 \\ 0 & 0 & 3 \end{vmatrix} + \begin{vmatrix} 2 & 0 & 0 \\ 0 & 3 & 0 \\ 0 & 0 & 4 \end{vmatrix} + \dots + \begin{vmatrix} n & 0 & 0 \\ 0 & n+1 & 0 \\ 0 & 0 & n+2 \end{vmatrix}$ is equal to
(A) $\frac{n^2-1}{2}$
(B) $\frac{n(n+1)(n+2)(n+3)}{4}$
(C) $\frac{n^2-1}{2}$
(D) None of these \hfill \textbf{((B))}
\item Q. Let A be a 2 $\times$ 2 matrix
Statement-1: adj(adj A) = A
Statement-2: $|\text{adj A}| = |A|$
(A) Statement- 1 is true, Statement-2 is true; Statement-2 is not a correct explanation for Statement-1
(B) Statement- 1 is true, Statement- 2 is false
(C) Statement- 1 is false, Statement- 2 is true
(D) Statement-1 is true, Statement-2 is true; Statement-2 is correct explanation for Statement-1 \hfill \textbf{((A))}
\item Q. Matrix A is such that $A^2 = 2A - I$, where I is the identity matrix. Then for $n \ge 2$, $A^n$ is equal to
(A) $nA - (n - 1)I$
(B) $nA - I$
(C) $2^{n-1}A - (n - 1)I$
(D) $2^{n-1}A - I$ \hfill \textbf{((A))}
\item Q. Consider the set A of all determinants of order 3 with entries 0 or 1 only. Let B be the subset of A consisting of all determinants with value 1. Let C be the subset of the set of all determinants with value -1. Then,
(A) C is empty
(B) B has as many elements as C
(C) A = B$\cup$C
(D) B has twice as many elements as C \hfill \textbf{((B))}
\end{enumerate}
\section{Mensuration}
\subsection{Easy}
\begin{enumerate}
\item Q. A farmer connects a pipe of internal diameter 20 cm from a canal into a cylindrical tank in her field, which is 10 m in diameter and 2 m deep if water flows through the pipe at the rate of 3 km/hr. In how much time will the tank be filled?
(A) 1 hr
(B) 1 hr 20 minutes
(C) 11 hr 30 minutes
(D) 11 hr 40 minutes \hfill \textbf{((D))}
\end{enumerate}
\section{Name}
\subsection{Medium}
\begin{enumerate}
\item Q. For $x \in (0, \frac{\pi}{2})$, define
$f(x) = \sqrt{x}$, $g(x) = \tan x$, $h(x) = \frac{1-x^2}{1+x^2}$
Let $\phi(x) = (h \circ f) (g(x))$. Then
$\phi(\frac{\pi}{3}) =$
(A) $\tan \frac{7\pi}{12}$
(B) $\tan \frac{\pi}{12}$
(C) $\tan \frac{5\pi}{12}$
(D) $\tan \frac{11\pi}{12}$ \hfill \textbf{((D))}
\item Q. Let
$u = \frac{2z + i}{z - ki}$, $z = x + iy$, $k > 0$.
The locus $\Re(u) + \Im(u) = 1$ meets the y-axis at two points P, Q with PQ = 5. Then
$k = ...$
(A) $\frac{3}{2}$
(B) $\frac{1}{2}$
(C) 2
(D) 4 \hfill \textbf{((C))}
\item Q. Q. Let $0 < x < y < z$ be real numbers such that $\frac{1}{x}, \frac{1}{y}, \frac{1}{z}$ are in arithmetic progression
and $x, \sqrt{2}y, z$ are in geometric progression. If
$xy + yz + zx = \frac{3}{\sqrt{2}} xyz$,
then the value of
$3 (x + y + z)^2$
is
(A) 50
(B) 100
(C) 150
(D) 200 \hfill \textbf{((C))}
\item Q. Find the coefficient of $x^2$ in
$(2 – x^2) [(1 + 2x + 3x^2)^6 + (1 – 4x^2)^6]$.
(A) 107
(B) 106
(C) 108
(D) 155 \hfill \textbf{((B))}
\item Q. In a geometric progression, the pth, qth, and rth terms are a, b, c, respectively. Show that
$a^{q-r} b^{r-p} c^{p-q} =?$
(A) 0
(B) pqr
(C) abc
(D) 1 \hfill \textbf{((D))}
\item Q. If
$\lim_{t\to 0} k t \csc t = \lim_{t\to 0} t \csc(kt)$,
then k equals
(A) a number $\neq \pm1$
(B) -1
(C) 1
(D) 1 or -1 \hfill \textbf{((D))}
\item Q. A man of height 1.8m walks away from a lamp-post of height 6 m at 1.5 m/s. The tip of
his shadow moves at speed p m/s. Then [p] (the greatest integer less than or equal to p) is
(A) 0
(B) 2
(C) 1
(D) 3 \hfill \textbf{((B))}
\item Q. Let
$f(x) = \int \frac{\sqrt{x}}{(1+x)^2} dx (x > 0)$.
Then f(3) - f(1) equals
(A) $\frac{\pi}{12} + \frac{1}{2} - \frac{\sqrt{3}}{4}$
(B) $-\frac{\pi}{12} + \frac{1}{2} + \frac{\sqrt{3}}{4}$
(C) $-\frac{\pi}{6} + \frac{1}{2} + \frac{\sqrt{3}}{4}$
(D) $\frac{\pi}{6} + \frac{1}{2} - \frac{\sqrt{3}}{4}$ \hfill \textbf{((A))}
\item Q. A triangle has its side-midpoints at (0,1), (1,1), and (1,0). The x-coordinate of its
incentre is
(A) $2-\sqrt{2}$
(B) $1 + \sqrt{2}$
(C) $2 + \sqrt{2}$
(D) $1 - \sqrt{2}$ \hfill \textbf{((A))}
\item Q. A circle $\omega$ has equation
$x^2 + y^2 + 2gx + 2fy + c = 0$.
Let AB be a chord of $\omega$ that subtends a right angle at the origin O. If $OP \perp AB$, then the
locus of the foot P is
(A) a parabola
(B) a circle
(C) an ellipse
(D) a straight line \hfill \textbf{((B))}
\item Q. The latus rectum of the parabola $y^2 = 5x$ has endpoints $(x_1,y_1)$ and $(x_2,y_2)$. Compute
$4x_1x_2 + y_1y_2$.
(A) 0
(B) $\frac{5}{4}$
(C) 5
(D) 25 \hfill \textbf{((A))}
\item Q. Solve the differential relation
$x^3dy + xydx = x^2 dy + 2y dx$,
with the initial condition $y(2) = e$ and $x > 1$. Find $y(4)$.
(A) $\frac{\sqrt{e}}{2}$
(B) $\frac{4}{3}\sqrt{e}$
(C) $\frac{\sqrt{e}}{3} + \sqrt{e}$
(D) $\frac{\sqrt{e}}{4} + \sqrt{e}$ \hfill \textbf{((B))}
\item Q. A point $P(\alpha, \beta, \gamma)$ lies on the plane
$x + 2y+3z - 7 = 0$.
Find the minimum value of $\alpha^2 + \beta^2 + \gamma^2$.
(A) $\frac{\sqrt{7}}{2}$
(B) 0
(C) 7
(D) $\frac{7}{2}$ \hfill \textbf{((D))}
\item Q16. The numbers 1 to 100 are written on slips, placed in a box. A draws one slip at
random (then replaces it), and then B draws one slip at random. What is the probability
that B's number is larger than A's?
(A) $\frac{99}{20000}$
(B) $\frac{99}{500}$
(C) $\frac{49}{500}$
(D) $\frac{99}{200}$ \hfill \textbf{((D))}
\item Q. In a college of 300 students, every student reads 5 newspapers, and each newspaper is
read by 60 students. The total number of distinct newspapers is
(A) exactly 25
(B) at most 20
(C) at most 15
(D) at least 30 \hfill \textbf{((A))}
\end{enumerate}
\section{Not Specified}
\subsection{Easy}
\begin{enumerate}
\item Q. The equation of the circle passing through the points $(1,0)$ and $(0,1)$ with the smallest possible radius is: \hfill \textbf{((B))}
\item Q. The arithmetic mean of the terms of the sequence $3, 7, 11, 15, \dots, (4n - 1)$ is $31$. Find $n$. \hfill \textbf{((D))}
\end{enumerate}
\subsection{Medium}
\begin{enumerate}
\item Q. If $a^2 + b^2 = 7$ and $a^3 + b^3 = 10$, then the maximum value of $|a + b|$ is
(A) 4
(B) 5
(C) 6
(D) 3 \hfill \textbf{((B))}
\item Q. The area common to the curves $y = \sin^{-1}(\sin x)$ and $y = [\sin^{-1}(\sin x)]$ on the interval $0 \le x \le \pi$ is
(A) $(\pi - 2)^2$
(B) $\frac{(\pi-2)^2}{2}$
(C) $\frac{(\pi-2)^2}{4}$
(D) None of these \hfill \textbf{((C))}
\item Q. If $I_1 = \int_0^1 \frac{e^x}{x + 3}\,dx$ and $I_2 = \int_0^1 \frac{x^3}{e^{x^4}(4 - x^4)}\,dx$, then the value of $\frac{I_1}{I_2}$ is
(A) 4
(B) 4e
(C) 3e
(D) 2e \hfill \textbf{((B))}
\item Q. If $\sin \alpha + \cos \beta = \frac{1}{\sqrt{2}}$, $\cos \alpha + \sin \beta = \sqrt{\frac{2}{3}}$, then the value of $\tan(\frac{\alpha-\beta}{2})$ is
(A) $3\sqrt{3}-6$
(B) $\sqrt{3}-2$
(C) $4\sqrt{3}-7$
(D) $4\sqrt{3}+7$ \hfill \textbf{((C))}
\item Q. A die is so loaded that the probability of throwing face $k$ is proportional to $k^2$. Given that the outcome is not even, the probability that the number 3 appears is
(A) $\frac{9}{70}$
(B) $\frac{35}{91}$
(C) $\frac{9}{91}$
(D) $\frac{9}{35}$ \hfill \textbf{((D))}
\item Q. The digit in the unit place of $15^{981} + 9^{789} – 6^{50}$ is
(A) 8
(B) 6
(C) 5
(D) 9 \hfill \textbf{((A))}
\item Q. The circle $x^2 + y^2 - 8x = 0$ and the hyperbola $\frac{x^2}{9} - \frac{y^2}{4} = 1$ have a common tangent with positive slope. Its equation is
(A) $2x - \sqrt{5}y - 20 = 0$
(B) $2x - \sqrt{5}y + 4 = 0$
(C) $3x – 4y +8=0$
(D) $4x - 3y +4=0$ \hfill \textbf{((B))}
\item Q. Two vertices of a triangle are $(5, -1)$ and $(-2, 3)$. If the orthocenter of the triangle is at the origin, the third vertex is
(A) $(4,7)$
(B) $(3,7)$
(C) $(-4, -7)$
(D) $(4,-7)$ \hfill \textbf{((C))}
\item Q. Which of the following intervals is a possible domain of the function $f(x) = \log_{\{x\}}([x]) + \log_{[x]}(\{x\})$, where $[x]$ is the greatest integer (floor) of $x$ and $\{x\} = x - [x]$ is the fractional part?
(A) $(0,1)$
(B) $(1,2)$
(C) $(2,3)$
(D) $(3,5)$ \hfill \textbf{((C))}
\item Q. If $f: \mathbb{R} \to \mathbb{R}$ is differentiable with $f(1) = 4$, then $\lim_{x \to 1} \frac{1}{x-1} \int_4^{f(x)} 2t\,dt$ equals
(A) $8 f'(1)$
(B) $4 f'(1)$
(C) $2 f'(1)$
(D) $f'(1)$ \hfill \textbf{((A))}
\item Q. If $f(x) = (x + 2019)^n$ where $n$ is a positive integer, then the value of $f(0) + f'(0) + \frac{f''(0)}{2!} + \dots + \frac{f^{(n-1)}(0)}{(n-1)!}$ is
(A) $2019^n$
(B) $2020^n$
(C) $2020^n - 1$
(D) $n 2019^n$ \hfill \textbf{((C))}
\item Q. The shortest distance from the line $3x + 4y = 25$ to the circle $x^2 + y^2 = 6x – 8y$ is
(A) $\frac{7}{5}$
(B) $\frac{9}{5}$
(C) $\frac{11}{5}$
(D) $\frac{32}{5}$ \hfill \textbf{((A))}
\item Q. Let $g(x) = \cos(x^2)$, $f(x) = \sqrt{x}$, and let $\alpha, \beta (\alpha < \beta)$ be the roots of $18x^2 – 9\pi x + \pi^2 = 0$. The area bounded by the curve $y = (g \circ f)(x)$, the lines $x = \alpha, x = \beta$, and the $x$-axis $(y = 0)$ is
(A) $\frac{\sqrt{3}+1}{2}$
(B) $\frac{\sqrt{3}-\sqrt{2}}{2}$
(C) $\frac{\sqrt{2}-1}{2}$
(D) $\frac{\sqrt{3}-1}{2}$ \hfill \textbf{((D))}
\item Q. Let $f(x) = \int_0^x \sin(t^2 - t + x)\,dt$. Then the value of $f''(-\frac{1}{2}) + f(-\frac{1}{2})$ is
(A) 0
(B) 1
(C) $\frac{1}{2}$
(D) $-\frac{1}{2}$ \hfill \textbf{((A))}
\item Q. The range of the function $f(x) = \frac{1}{5 \cos x - 12 \sin x + 15}$ is: \hfill \textbf{((C))}
\item Q. Let $z$ be a complex number with $|z| = 1$ and $\arg z = \theta$. Compute $\arg \left(\frac{1+z}{1+\bar{z}}\right)$. \hfill \textbf{((D))}
\item Q. The term independent of $x$ in the expansion of $\left(\frac{x+1}{x^{2/3} - x^{1/3}+1} - \frac{x-1}{x - \sqrt{x}}\right)^{10}$ is: \hfill \textbf{((C))}
\item Q. A train must stop at exactly 4 out of 10 stations along its route. In how many ways can these 4 stopping-stations be chosen so that at least two of the chosen stops are consecutive? \hfill \textbf{((C))}
\item Q. The sum of the first $n$ terms of an arithmetic progression (A.P.) is $S(n) = 2n + 3n^2$. A second A.P. has the same first term but twice the common difference of the first. The sum of its first $n$ terms is \hfill \textbf{((C))}
\item Q. If $y = \sec(\tan^{-1}x)$, then the derivative $\frac{dy}{dx}$ at $x = 1$ is: \hfill \textbf{((A))}
\item Q. The area enclosed by the parabola $y = x^2 + 1$ and the normal to it whose slope is $-1$ is: \hfill \textbf{((A))}
\item Q. The normal to the parabola $y^2 = 4ax$ at the point with parameter $t_1$ meets the curve again at the point with parameter $t_2$. If $\angle POQ = \frac{\pi}{2}$, where $O$ is the origin, then which relation holds? \hfill \textbf{((B))}
\item Q. A plane P meets the coordinate axes at A, B, and C. The centroid of $\triangle ABC$ is $(1, 1, 2)$. Find the equation of the line through this centroid that is perpendicular to the plane P. \hfill \textbf{((C))}
\item Q. Let $\vec{a}$ and $\vec{b}$ be non-perpendicular vectors, and $\vec{c}$, $\vec{d}$ satisfy $\vec{b} \times \vec{c} = \vec{b} \times \vec{d}$ and $\vec{a} \cdot \vec{d} = 0$. Then $\vec{d}$ must be \hfill \textbf{((C))}
\item Q. Evaluate $\cos\left(\frac{2\pi}{7}\right) + \cos\left(\frac{4\pi}{7}\right) + \cos\left(\frac{6\pi}{7}\right)$. \hfill \textbf{((B))}
\end{enumerate}
\subsection{Hard}
\begin{enumerate}
\item Q. Define $g(x) = \frac{1}{\ln(1+x)} - \frac{1}{x}$, $x > 0$. For which interval does $g(x)$ take its values? \hfill \textbf{((B))}
\item Q. Suppose $y = y(x)$ satisfies the differential equation $\frac{dy}{dx} + y \tan x = 2x + x^2 \tan x$, $x \in (-\frac{\pi}{2}, \frac{\pi}{2})$, with the initial condition $y(0) = 1$. Which of the following is true? \hfill \textbf{((D))}
\end{enumerate}
\section{Number Theory}
\subsection{Medium}
\begin{enumerate}
\item Q. If $(1 + 3 + 5 + \dots + p) + (1 + 3 + 5 + \dots + q) = (1 + 3 + 5 + \dots + r)$, where each set of parentheses contains the sum of consecutive odd integers as shown, the smallest possible value of $p + q + r$. (where $p > 6$) is
(A) 12
(B) 21
(C) 45
(D) 54 \hfill \textbf{((B))}
\item Q. Let S be the set of pairs of natural numbers (x, y) with x \leq y satisfying $\frac{xy}{x + y + \gcd(x, y)}$. How many such pairs (x, y) are there? \hfill \textbf{((B))}
\end{enumerate}
\section{One-To-One Functions And Derivatives}
\subsection{Medium}
\begin{enumerate}
\item Q. Let
$\quad f(x) = \frac{x^2 + 3x + a}{x^2+x+1},\quad f: \mathbb{R}\to\mathbb{R}.$ For which real value(s) of the parameter a is f one-to-one (injective) on R?
(A) a = 2
(B) a = 1
(C) a = 3
(D) No real value of a \hfill \textbf{((D))}
\end{enumerate}
\section{Parabola}
\subsection{Medium}
\begin{enumerate}
\item Q. A ray of light moving parallel to the x-axis is reflected from the parabolic mirror $(y - 2)^2 = 4(x + 1)$. The point on the axis of the parabola through which the reflected ray must pass is $(0,k)$, where $k \in \mathbb{N}$. Then k is \hfill \textbf{((B))}
\end{enumerate}
\section{Parity Of Functions}
\subsection{Medium}
\begin{enumerate}
\item Q. Which of the following correctly describes the parity of the functions below?
(i) $\ f(x) = x^2 \sin x,$
(ii) $\ f(x) = \sqrt{1 + x + x^2} - \sqrt{1 - x + x^2},$
(iii) $\ f(x) = \ln(\frac{1-x}{1+x}),$
(iv) $\ f(x) = \sin x - \cos x,$
(v) $\ f(x) = \frac{e^x + e^{-x}}{2}$
(A) (i, ii, iii) are odd, (iv) is neither, (v) is even.
(B) (i, ii) are odd, (iii, v) are even, (iv) is neither.
(C) (i, ii, iii, v) are odd, (iv) is neither.
(D) (i, ii, v) are odd, (iii, iv) are neither. \hfill \textbf{((A))}
\end{enumerate}
\section{Periodicity Of Functions}
\subsection{Medium}
\begin{enumerate}
\item Q. Determine which of the following functions are periodic, and if so, find their fundamental
period (where $[x]$ is the greatest-integer function):
(i) $\ f(x) = e^{x-[x]} + \sin x,$
(ii) $\ f(x) = \sin(\frac{\pi x}{\sqrt{2}}) + \cos(\frac{\pi x}{\sqrt{3}}),$
(iii) $\ f(x) = \sin(\frac{\pi x}{\sqrt{3}}) + \cos(\frac{2\pi x}{\sqrt{3}}).$
(A) (i) not periodic, (ii) not periodic, (iii) periodic with period $4\sqrt{3}.$
(B) (i) periodic with period $2\pi,$ (ii) not periodic, (iii) periodic with period $2\sqrt{3}.$
(C) (i) not periodic, (ii) periodic with period $2\sqrt{6},$ (iii) periodic with period $4\sqrt{3}.$
(D) (i) periodic with period 1, (ii) periodic with period $2\sqrt{2},$ (iii) not periodic. \hfill \textbf{((A))}
\end{enumerate}
\section{Permutations And Combinations}
\subsection{Easy}
\begin{enumerate}
\item Q. How many number of five digits can be made with at least one repeated digit?
(A) $62784$
(B) $61485$
(C) $65648$
(D) none of these \hfill \textbf{((A))}
\end{enumerate}
\subsection{Medium}
\begin{enumerate}
\item Q. A five letter word is to be formed such that the letters appearing in the odd numbered
positions are taken from the letters which appear without repetition in the word MATH-
EMATICS. Further the letters appearing in the even numbered positions are taken from
the letters which appear with repetitions in the same word MATHEMATICS. In how many
different ways the five letter word can be formed-
(A) 390
(B) 600
(C) 540
(D) 450 \hfill \textbf{((C))}
\item Q. How many numbers of five digits can be made having exactly two identical digits? How many of these will have the repeated digits in consecutive places?
(A) 45360,18144
(B) 45370,18144
(C) 18144,45360
(D) 3456,1845 \hfill \textbf{((A))}
\end{enumerate}
\section{Polynomial Inequalities}
\subsection{Medium}
\begin{enumerate}
\item Number of positive integral values of $x$ satisfying the inequality $\frac{(x - 4)^{2017} (x + 8)^{2016} (x + 1)}{x^{2016} (x - 2)^3 (x + 3)^5 (x - 6) (x + 9)^{2018}} \leq 0$ is \hfill \textbf{(D)}
\end{enumerate}
\section{Probability}
\subsection{Easy}
\begin{enumerate}
\item Q. Three different numbers are selected at random from the set A = {1,2,3, ..., 15}. The probability that the product of two is equal to the third one is
(A) $\frac{7}{65}$
(B) $\frac{1}{65}$
(C) $\frac{1}{13}$
(D) $\frac{5}{13}$ \hfill \textbf{((A))}
\end{enumerate}
\subsection{Medium}
\begin{enumerate}
\item Q. Ravi and Rashmi are each holding 2 red cards and 2 black cards (all four red and all four black cards are identical). Ravi picks a card at random from Rashmi and then Rashmi picks a card at random from Ravi. This process is repeated a second time. Let $p$ be the probability that both have all 4 cards of the same colour. Then, $p$ satisfies \hfill \textbf{((A))}
\item Q. If the integers $m\&n$ are chosen at random from 1 to 100, then the probability that a
number of the form $7^m + 7^n$ is divisible by 5 equals:
(A) $\frac{1}{8}$
(B) $\frac{1}{4}$
(C) $\frac{1}{10}$
(D) $\frac{1}{7}$ \hfill \textbf{((B))}
\item Q. An urn contains $m$ white and $n$ black balls. A ball is drawn at random and is put back into the urn along with $k$ additional balls of the same colour as that of the ball drawn. A ball is again drawn at random. What is the probability that the ball drawn now is white?
(A) $\frac{m}{m+n}$
(B) $\frac{m}{m-n}$
(C) $\frac{m}{n}$
(D) None of these \hfill \textbf{((A))}
\item Let X be a set containing 10 elements and P(X) be its power set. If A and B are picked up at random from P(X), with replacement, then the probability that A and B have equal number of elements is:\\
(A) $\frac{(2^{10}-1)}{2^{20}}$\\
(B) $\frac{{}^{20}C_{10}}{2^{20}}$\\
(C) $\frac{{}^{20}C_{10}}{2^{10}}$\\
(D) $\frac{(2^{10}-1)}{2^{10}}$ \hfill \textbf{((B))}
\item Q. Let X_{10} = \{1,2,...,10\}. A subset A \subseteq X_{10} is chosen uniformly at random subject
to the condition that any two elements of A differ by at least 3. If p = Pr(1 \in A) and
q = Pr(2 \in A), then
(A) p > q and p -q= \frac{1}{6}
(B) p < q and q-p= \frac{1}{6}
(C) p > q and p-q= \frac{1}{10}
(D) p < q and q-p= \frac{1}{10} \hfill \textbf{((C))}
\item Q. An urn contains marbles of four colours: red, white, blue, and green. When four marbles
are drawn without replacement, the following events are equally likely:
1. four red marbles,
2. one white and three red marbles,
3. one white, one blue, and two red marbles,
4. one marble of each colour.
What is the smallest total number of marbles in the urn satisfying these conditions?
(A) 19
(B) 21
(C) 46
(D) 69 \hfill \textbf{((B))}
\item Q. In an examination consisting of four parts A, B, C, D each containing 30 questions, a student randomly attempts each question (with probability 1/2). The probability that he attempts at least 60 questions is
(A) $\frac{1}{2} + \frac{1}{2} \frac{\binom{120}{60}}{2^{120}}$
(B) $\frac{1}{2} - \frac{1}{2} \frac{\binom{120}{60}}{2^{120}}$
(C) $\frac{1}{2} + \frac{\binom{120}{60}}{2^{120}}$
(D) $\frac{1}{2} - \frac{\binom{120}{60}}{2^{120}}$ \hfill \textbf{((A))}
\item Q. If n positive integers are chosen at random (each last digit equally likely), then the probability that the last digit of their product is one of 2, 4, 6, 8 equals
(A) $(\frac{4}{10})^n$
(B) $(\frac{4}{10})^n - (\frac{3}{10})^n$
(C) $(\frac{3}{10})^n$
(D) $(\frac{2}{5})^n$ \hfill \textbf{((D))}
\end{enumerate}
\subsection{Hard}
\begin{enumerate}
\item Q. Let $X \sim \text{Binomial}(n,p)$. Suppose the difference of its mean and variance is 1, and $2 P(X = 2) = 3 P(X = 1)$. Then $n^2 P(X > 1)$ equals
(A) 15
(B) 16
(C) 12
(D) 11 \hfill \textbf{((D))}
\end{enumerate}
\section{Properties Of Logarithms And Domain}
\subsection{Medium}
\begin{enumerate}
\item Q. The functions
$\quad f(x) = \log(x - 1) - \log(x - 2),\quad g(x) = \log(\frac{x-1}{x-2})$
are identical for values of x in which interval?
(A) $[1,2]$
(B) $[2,\infty)$
(C) $(2,\infty)$
(D) $(-\infty,\infty)$ \hfill \textbf{((C))}
\end{enumerate}
\section{Quadratic Equations}
\subsection{Medium}
\begin{enumerate}
\item Q. If $a, b \in R$, then the equation $x^2 - abx - a^2 = 0$ has
(A) one positive and one negative root
(B) both positive roots
(C) both negative roots
(D) non-real roots \hfill \textbf{((A))}
\item Q. If $\alpha$ and $\beta$, $\alpha$ and $\gamma$, $\alpha$ and $\delta$ are the roots of the equations $ax^2+2bx+c = 0$, $2bx^2+cx+a = 0$ and $cx^2 + ax + 2b = 0$, respectively, where $a, b$ and $c$ are positive real numbers, then $\alpha+\alpha^2 =$
(A) $abc$
(B) $a + 2b + c$
(C) -1
(D) 0 \hfill \textbf{((C))}
\item Q. If $\alpha, \beta$ are the roots of the quadratic equation $x^2 - 3x + 5 = 0$, then the equation whose roots are $\alpha^2 - 3\alpha + 7$ and $\beta^2 - 3\beta + 7$ is (A) $x^2 + 4x + 1 = 0$ (B) $x^2 - 4x + 4 = 0$ (C) $x^2 - 4x - 1 = 0$ (D) $x^2 + 2x + 3 = 0$ \hfill \textbf{((B))}
\item Q. Find all integral values of a for which the quadratic equation $(x - a) (x - 10) + 1 = 0$ has integral roots. (A) a = 8 or 12 (B) a = 9 or 11 (C) a = 7 or 13 (D) a = 6 or 14 \hfill \textbf{((A))}
\item Q. If the equation $\frac{x^2 - bx}{ax - c} = \frac{k-1}{k+1}$ has roots equal in magnitude and opposite in sign, then the value of k is (A) $\frac{a+b}{a-b}$ (B) $\frac{a-b}{a+b}$ (C) $\frac{a}{b}+1$ (D) $\frac{a}{b}-1$ \hfill \textbf{((B))}
\item Q. Determine all values of $\lambda$ for which the equation $(\lambda^2 - 5\lambda + 6) x^2 + (\lambda^2 - 3\lambda + 2) x + (\lambda^2 - 4) = 0$ holds as an identity in x (i.e. is true for more than two distinct values of x). (A) $\lambda$ = 1 (B) $\lambda$ = 2 (C) $\lambda$ = 3 (D) $\lambda$ = -2 \hfill \textbf{((B))}
\item Q. The equations $5x^2 + 12x + 13 = 0$ and $ax^2 + bx + c = 0$ have a common root, where a, b, c $\in$ R are the side-lengths of $\triangle ABC$. Then $\angle C$ is (A) 45° (B) 60° (C) 90° (D) 30° \hfill \textbf{((C))}
\item Q. If $a^2 + 4b^2 + 9c^2 + 2ab + 6bc + 3ac \leq 0, a, b, c \in \mathbb{R}$, then the number of distinct real roots of the equation $ax^2 + bx + c = 0$ is/are \hfill \textbf{((D))}
\end{enumerate}
\section{Range Of Composite Functions}
\subsection{Medium}
\begin{enumerate}
\item Q. Let
$\quad f(x) = \begin{cases} [x], & -2 \le x \le -1, \\ |x| + 1, & -1 < x \le 2, \end{cases} \quad g(x) = \begin{cases} [x], & -\pi \le x < 0, \\ \sin x, & 0 \le x \le \pi. \end{cases}$
Define $\ h(x) = g(f(x)).$ The range of h is
(A) $\{-2,-1\} \cup [\sin 3, 1]$
(B) $\{-2, 1\} \cup [\sin 3, 1]$
(C) $\{-2,-1\} \cup [0, 1]$
(D) $[\sin 3,1]$ \hfill \textbf{((A))}
\end{enumerate}
\section{Range Of Functions}
\subsection{Medium}
\begin{enumerate}
\item Q. Determine the range of the function
$\quad f(x) = \log_2(\frac{\sin x - \cos x + 3\sqrt{2}}{\sqrt{2}}).$
(A) $[0,1]$
(B) $[1,2]$
(C) $[-1,1]$
(D) $[2,3]$ \hfill \textbf{((B))}
\item Q. Determine the range of the function
$\quad f(x) = \log_{\sqrt{2}}(2 - \log_2(16 \sin^2x + 1)).$
(A) $(-\infty,1]$
(B) $(-\infty,2]$
(C) $[0,2]$
(D) $[1,2]$ \hfill \textbf{((B))}
\end{enumerate}
\section{Rational Inequalities}
\subsection{Medium}
\begin{enumerate}
\item Number of non-negative integral values of $x$ satisfying the inequality $\frac{2}{x^2 - x + 1} - \frac{1}{x+1} - \frac{2x-1}{x^3+1} \geq 0$ is \hfill \textbf{(D)}
\end{enumerate}
\section{Relations And Functions}
\subsection{Medium}
\begin{enumerate}
\item Q. Let A = {1, 2, ..., 10} and define a relation R on A by
$(x, y) \in R \iff x + 2y = 10, \quad x, y \in A.$
Which of the following correctly describes R, its inverse R¯¹, and the domain and range of
R?
(A) R = {(2,4), (4, 3), (6, 2), (8, 1)}, R¯¹ = {(4, 2), (3, 4), (2, 6), (1, 8)},
Dom(R) = {2,4,6,8}, Range(R) = {4,3,2,1}.
(B) R = {(2, 4), (4, 3), (6, 2), (8,1)}, R-1 = {(4, 2), (3, 4), (2, 6), (1,8)},
Dom(R) = {1,2,3,4}, Range(R) = {2,4,6,8}.
(C) R = {(2,4), (4, 3), (6, 2), (8,1)}, R¯¹ = {(2, 4), (4, 3), (6, 2), (8, 1)},
Dom(R) = {2,4,6,8}, Range(R) = {4,3,2,1}.
(D) R = {(2,4), (4, 3), (6, 2), (8, 1)}, R¯¹ = {(4, 2), (3, 4), (2, 6), (1, 8)},
Dom(R) = {2,4,6}, Range(R) = {4,3,2}. \hfill \textbf{((A))}
\end{enumerate}
\section{Sequences And Series}
\subsection{Medium}
\begin{enumerate}
\item Q. If $a, x, y, z, b$ are in A.P., then $x + y + z = 15$. If $a, x, y, z, b$ are in H.P., then $\frac{1}{x} + \frac{1}{y} + \frac{1}{z} = \frac{5}{3}$. The numbers $a, b$ are:
(A) 8,2
(B) 9,1
(C) 8,1
(D) 11,3 \hfill \textbf{((B))}
\item Q. Sum of three numbers in a G.P. is 21. If 2 is subtracted from the first term and 1 is subtracted from the third term, the new numbers are in A.P. The largest of the original numbers is:
(A) 18
(B) 15
(C) 9
(D) 12 \hfill \textbf{((D))}
\item Q. Let $a$ be a fixed real number and $x, y, z$ satisfy
$\frac{a-x}{px} = \frac{a-y}{qy} = \frac{a-z}{rz}$
If $p, q, r$ are in arithmetic progression, then $x, y, z$ are in
(A) arithmetic progression (A.P.)
(B) geometric progression (G.P.)
(C) harmonic progression (H.P.)
(D) none of these \hfill \textbf{(C)}
\item Q. The positive integers are partitioned into consecutive “groups” as follows:
(1) , (2,3,4), (5,6,7,8,9), ...
1st group 2nd group 3rd group
where the $n$th group contains $2n - 1$ numbers. What is the sum of the numbers in the $n$th group?
(A) $n^3 + (n - 1)^3$
(B) $n^3 + (n + 1)^3$
(C) $2n^3 - n$
(D) $(2n - 1) n^2$ \hfill \textbf{(A)}
\end{enumerate}
\subsection{Hard}
\begin{enumerate}
\item Q. Let $a, b, c, d, p$ be distinct real numbers satisfying
$(a^2 + b^2 + c^2)p^2 - 2p (ab+bc+cd) + (b^2+c^2 + d^2) \leq 0$.
Then $a, b, c, d$ must be in
(A) Arithmetic progression (A.P.)
(B) Geometric progression (G.P.)
(C) Harmonic progression (H.P.)
(D) None of these \hfill \textbf{(B)}
\end{enumerate}
\section{Series}
\subsection{Medium}
\begin{enumerate}
\item Q. Find the sum of the first 16 terms of the series
$\frac{1^3}{1} + \frac{1^3 + 2^3}{1+3} + \frac{1^3 + 2^3 + 3^3}{1+3+5} + ...$
The value is
(A) 450
(B) 456
(C) 446
(D) none of these \hfill \textbf{(C)}
\end{enumerate}
\subsection{Hard}
\begin{enumerate}
\item Q. Find the sum of the infinite series
$4 - 9x + 16x^2 - 25x^3 + 36x^4 - 49x^5 + ...$
for $|x| < 1$. The sum $S$ equals
(A) $\frac{4 + 3x + x^2}{(1-x)^3}$
(B) $\frac{4 + 3x + x^2}{(1+x)^3}$
(C) $\frac{4- 3x + x^2}{(1+x)^3}$
(D) $\frac{4+ 3x - x^2}{(1+x)^3}$ \hfill \textbf{(B)}
\item Q. Let $T_r$ be a sequence whose partial sums satisfy
$\sum_{r=1}^{n} T_r = \frac{n(n + 1)(n+2)(n + 3)}{8}$.
Compute the infinite series $\sum_{r=1}^{\infty} \frac{1}{T_r}$.
(A) $\frac{1}{2}$
(B) $\frac{1}{3}$
(C) $\frac{2}{3}$
(D) $\frac{3}{4}$ \hfill \textbf{(A)}
\end{enumerate}
\section{Series And Summation}
\subsection{Medium}
\begin{enumerate}
\item Q. Let $x_1, x_2, \dots, x_{2019}$ be real numbers, none equal to 1, satisfying $\sum_{i=1}^{2019} \frac{x_i}{1 - x_i} = 1, \sum_{i=1}^{2019} \frac{x_i}{x_i - 1} = 1$. Find the value of $\sum_{i=1}^{2019} \frac{x_i^2}{1 - x_i}$. \hfill \textbf{((A))}
\end{enumerate}
\section{Set Theory And Counting}
\subsection{Medium}
\begin{enumerate}
\item Thirty-one candidates appeared for an examination. Fifteen candidates passed in English, fifteen in Hindi, and twenty in Sanskrit. Three candidates passed only in English, four passed only in Hindi, and seven passed only in Sanskrit. Two candidates passed in all three subjects. How many candidates passed in exactly two subjects? \hfill \textbf{(B)}
\end{enumerate}
\section{Set Theory And Logic}
\subsection{Medium}
\begin{enumerate}
\item Consider the following statements:\n1. $N \cup (B \cap Z) = (N \cup B) \cap Z$ for any subset $B$ of $\mathbb{R}$, where $N$ is the set of positive integers, $Z$ is the set of integers, and $\mathbb{R}$ is the set of real numbers.\n2. Let $A = \{n \in N : 1 \leq n \leq 24, n \text{ is a multiple of } 3\}$. There exists no subset $B$ of $N$ such that $|A| = |B|$.\nWhich of the above statements is/are correct? \hfill \textbf{(A)}
\end{enumerate}
\section{Set Theory And Number Theory}
\subsection{Medium}
\begin{enumerate}
\item Let $U$ be a set with number of elements in it equal to $2009$. $A$ is a subset of $U$ with $n(A) = 1681$ and out of these $1681$ elements, exactly $1075$ elements belong to a subset $B$ of $U$. If $n(A - B) = m^2 + p_1p_2p_3$ for some positive integer $m$ and distinct primes $p_1 < p_2 < p_3$, then for the least $m$, the value of $\frac{p_1p_3}{p_2}$ is \hfill \textbf{(A)}
\item Let $U$ be a set with number of elements in it equal to $2009$, and $A, B$ be subsets of $U$ with $n(A \cup B) = 280$. If $n(A' \cap B') = x_1^3 + x_2^3 = y_1^3 + y_2^3$ for some positive integers satisfying $x_1 < y_1 < y_2 < x_2$, then the value of $\frac{x_2 + y_2}{x_1 + y_1}$ is \hfill \textbf{(A)}
\end{enumerate}
\section{Statistics}
\subsection{Medium}
\begin{enumerate}
\item Q. Let $x_1,x_2,...,x_n$ be a data series with arithmetic mean $\bar{x}$. A new series is formed by adding $2i$ to the $i$th term, i.e. the terms are $x_i + 2i$ for $i = 1, 2, \dots, n$. The mean of this new series is $\bar{X}=?$ \hfill \textbf{(B)}
\item Q. A discrete random variable $X$ takes values $0, 1, 2, \dots, n$ with frequencies proportional to the binomial coefficients $\binom{n}{0}, \binom{n}{1}, \dots, \binom{n}{n}$. The arithmetic mean $\mu$ of this distribution is \hfill \textbf{(B)}
\item Q. Find the weighted mean of the first $n$ natural numbers $1, 2, \dots, n$, when each number $i$ is assigned weight $i^2$. That is, $\bar{X} = \frac{\sum_{i=1}^n i \cdot i^2}{\sum_{i=1}^n i^2}$. \hfill \textbf{(A)}
\item Q. A student scores 75%, 80%, and 85% in three subjects. A fourth subject is then added to his record. What is the minimum possible overall average percentage after including this fourth subject (assuming percentages range from 0% to 100%)? \hfill \textbf{(A)}
\item Q. The following frequency distribution shows the marks obtained by 100 students in a test. Find the median of this distribution. \hfill \textbf{(A)}
\item Q. The following data give the values $x_i$ and their frequencies $f_i$. Find the mean deviation about the mean. \hfill \textbf{(C)}
\item Q. Consider the arithmetic progression $a, a + d, a + 2d, \dots, a + 2nd$ which has $2n + 1$ terms. Its mean (arithmetic average) is $a + nd$. The mean deviation about the mean is defined as $\frac{1}{2n+1}\sum_{i=0}^{2n}|a+id - (a+nd)|$. This mean deviation equals \hfill \textbf{(B)}
\item Q. For a set of $n$ observations $x_1,x_2,...,x_n$, the mean square deviation about a point $c$ is defined by $\frac{1}{n}\sum_{i=1}^n(x_i - c)^2$. It is known that $\frac{1}{n}\sum_{i=1}^n(x_i + 2)^2 = 18$, $\frac{1}{n}\sum_{i=1}^n(x_i - 2)^2 = 10$. What is the standard deviation $\sigma$ of the observations? \hfill \textbf{(C)}
\item Q. The following frequency distribution is given (class-intervals of width 2): Find the variance $\sigma^2$ of this distribution (to two decimal places). \hfill \textbf{(A)}
\item Q. A data set $A$ has 5 observations with mean 8 and variance 24, and another data set $B$ has 3 observations also with mean 8 and variance 24. If the two sets are combined into one series of 8 observations, what is the variance of the combined series? \hfill \textbf{(B)}
\item Q. Find the variance of the first 20 natural numbers $1,2,...,20$ (treating them as the full population). \hfill \textbf{(A)}
\item Q. An experiment yields five observations whose arithmetic mean is 4 and population variance is 5.2. Three of the observations are 1, 2, and 6. What are the remaining two observations? \hfill \textbf{(C)}
\item Q. Let $\{x_i\}_{i=1}^{18}$ be a data set such that $\sum_{i=1}^{18}(x_i - 8) = 9$, $\sum_{i=1}^{18}(x_i - 8)^2 = 45$. What is the standard deviation $\sigma$ of the $x_i$'s? \hfill \textbf{(B)}
\item Q. The following seven data points are given: 340, 150, 210, 240, 300, 310, 320. Find the mean deviation about the median. \hfill \textbf{(C)}
\item Q. Let $a, b, c, d$ be real numbers satisfying $d > a > b > c$. If the mean of the four-point distribution $\{a, b, c, d\}$ is 5 and its median is 6, find the value of $-a + 3d + 3c - b$. \hfill \textbf{(C)}
\item Q. The standard deviation of the numbers 31, 32, 33, ..., 46, 47 is
(A) $2\sqrt{6}$
(B) $\sqrt{\frac{47^2-1}{12}}$
(C) $\sqrt{\frac{17}{12}}$
(D) $4\sqrt{3}$ \hfill \textbf{((A))}
\item Q. The mean of 5 observations is 4 and their variance is 9.2. If three of the observations are 1, 2 and 6, then the other two observations can be
(A) 2 and 9
(B) 3 and 8
(C) 4 and 7
(D) 5 and 6 \hfill \textbf{((A))}
\end{enumerate}
\section{Summation And Floor Function}
\subsection{Medium}
\begin{enumerate}
\item Q. Compute the sum
$\quad \sum_{k=0}^{2946} [\frac{1}{2} + \frac{k}{1000}],$
where $[\cdot]$ denotes the greatest integer function.
(A) 4332
(B) 4341
(C) 4342
(D) 4400 \hfill \textbf{((B))}
\end{enumerate}
\section{Trigonometry}
\subsection{Medium}
\begin{enumerate}
\item Q. Find the number of solutions $\theta \in [0, 2\pi]$ to the equation
$\sin(\pi \sin^2 \theta) + \sin(\pi \cos^2 \theta) = 2\cos(\frac{\pi}{2} \cos \theta)$. \hfill \textbf{((D))}
\item Q. If $\tan x = n\cdot \tan y$, $n \in R^+$ then maximum value of $\sec^2(x - y)$ is equal to
(A) $\frac{(n-1)^2}{2n}$
(B) $\frac{(n+1)^2}{2n}$
(C) $\frac{(n+1)^2}{2}$
(D) $\frac{(n+1)^2}{4n}$ \hfill \textbf{((D))}
\item Q. Given $a^2 + 2a + \operatorname{cosec}^2 \left(\frac{\pi}{2}(a+x)\right) = 0$, then which of the following holds good?
(A) $a = 1; \frac{x}{2} \in I$
(B) $a = -1; \frac{x}{2} \in I$
(C) $a \in R; x \in I$
(D) $a, x$ are finite but not possible to find \hfill \textbf{((B))}
\item Q. The number of positive integral solutions of the equation $\tan^{-1}x + \cos^{-1}\frac{y}{\sqrt{1+y^2}} = \sin^{-1}\frac{3}{\sqrt{10}}$ is
(A) 1
(B) 2
(C) 0
(D) none of these \hfill \textbf{((A))}
\item $\cot^{-1}(\sqrt{\cos \alpha}) – \tan^{-1}(\sqrt{\cos \alpha}) = x$ then $\sin x$ is equal to\\
(A) $\tan^2 (\frac{\alpha}{2})$\\
(B) $\cot^2 (\frac{\alpha}{2})$\\
(C) $\tan \alpha$\\
(D) $\cot (\frac{\alpha}{2})$ \hfill \textbf{((A))}
\item Q. For $x \in (0, \pi)$, the equation $\sin x + 2 \sin 2x - \sin 3x = 3$ has \hfill \textbf{((D))}
\item Q. If $\cos^{-1} \sqrt{p} + \cos^{-1} \sqrt{1 - p} + \cos^{-1} \sqrt{1 - q} = \frac{3\pi}{4}$, then the value of q is
(A) 1
(B) $\frac{1}{\sqrt{2}}$
(C) $\frac{1}{4}$
(D) $\frac{1}{2}$ \hfill \textbf{((D))}
\item Q. In any triangle ABC, if $\sin A - \cos B = \cos C$, then the angle B is \\ (A) $\frac{\pi}{2}$ \\ (B) $\frac{\pi}{3}$ \\ (C) $\frac{\pi}{4}$ \\ (D) $\frac{\pi}{6}$ \hfill \textbf{((a))}
\item Q. If $A + B + C = \frac{3\pi}{2}$, then $\cos 2A + \cos 2B + \cos 2C$ equals \\ (A) $1-4 \cos A \cos B \cos C$ \\ (B) $4 \sin A \sin B \sin C$ \\ (C) $1+2 \cos A \cos B \cos C$ \\ (D) $1 - 4 \sin A \sin B \sin C$ \hfill \textbf{((d))}
\item Q. Find the maximum value of $1 + \sin\left(\frac{\pi}{4}+\theta\right)+2\cos\left(\frac{\pi}{4}-\theta\right)$. \\ (A) $1$ \\ (B) $2$ \\ (C) $3$ \\ (D) $4$ \hfill \textbf{((d))}
\item Q. For all real $\theta$, the expression \\ $E(\theta) = 5 \cos\theta + 3\cos\left(\theta+\frac{\pi}{3}\right)+3$ \\ satisfies which of the following ranges? \\ (A) $-4 \le E(\theta) \le 10$ \\ (B) $-7 \le E(\theta) \le 7$ \\ (C) $-3 \le E(\theta) \le 9$ \\ (D) $0 \le E(\theta) \le 6$ \hfill \textbf{((a))}
\item Q. For any integer $n \ge 1$ and real $\theta$, evaluate the product \\ $P = (1 + \sec 2\theta) (1 + \sec 2^2\theta) (1 + \sec 2^3\theta) \cdots (1 + \sec 2^n\theta)$. \\ (A) $\tan 2^n\theta \cot \theta$ \\ (B) $\tan 2^n \tan \theta$ \\ (C) $\cot 2^n \tan \theta$ \\ (D) $\sec 2^n \cot \theta$ \hfill \textbf{((a))}
\item Q. For integer $n \ge 2$, evaluate the sum \\ $\sum_{r=1}^{n-1} \cos^2\left(\frac{r\pi}{n}\right)$. \\ (A) $\frac{n-2}{2}$ \\ (B) $\frac{n-1}{2}$ \\ (C) $\frac{n}{2}$ \\ (D) $\frac{n+1}{2}$ \hfill \textbf{((a))}
\item Q. Evaluate the expression \\ $\tan A + \tan(60^\circ + A) + \tan(120^\circ + A)$. \\ (A) $\tan 3A$ \\ (B) $2 \tan 3A$ \\ (C) $3 \tan 3A$ \\ (D) $4 \tan 3A$ \hfill \textbf{((c))}
\item Q. The product \\ $\tan(60^\circ + A) \tan(60^\circ - A)$ \\ is equal to \\ (A) $\frac{2 \cos 2A + 1}{2 \cos 2A-1}$ \\ (B) $\frac{2 \cos 2A-1}{2 \cos 2A + 1}$ \\ (C) $\frac{\cos 2A + 2}{\cos 2A - 2}$ \\ (D) $\frac{\cos 2A - 2}{\cos 2A + 2}$ \hfill \textbf{((a))}
\item Q. The value of \\ $\frac{\sin 5\theta+ \sin 2\theta - \sin \theta}{\cos 5\theta + 2 \cos 3\theta + 2 \cos^2 \theta + \cos \theta}$ \\ is \\ (A) $\tan\theta$ \\ (B) $\cos \theta$ \\ (C) $\cot \theta$ \\ (D) none of these \hfill \textbf{((a))}
\item Q. Evaluate the product \\ $\sin 12^\circ \sin 48^\circ \sin 54^\circ$. \\ (A) $\frac{1}{8}$ \\ (B) $\frac{1}{4}$ \\ (C) $\frac{1}{16}$ \\ (D) $\frac{1}{2}$ \hfill \textbf{((a))}
\item Q. If $\sin 2A = \lambda \sin 2B$, then the ratio \\ $\frac{\tan(A + B)}{\tan(A - B)}$ \\ is equal to \\ (A) $\frac{\lambda+1}{\lambda-1}$ \\ (B) $\frac{\lambda-1}{\lambda+1}$ \\ (C) $\frac{\lambda+1}{1-\lambda}$ \\ (D) $\frac{1-\lambda}{\lambda+1}$ \hfill \textbf{((a))}
\item Q. Solve for $x$: \\ $3 \cos x + 4 \sin x = 5$. \\ (A) $x = 2n\pi + 2 \arctan\frac{1}{2}$ \\ (B) $x = 2n\pi + \arctan\frac{1}{2}$ \\ (C) $x = n\pi + 2 \arctan\frac{1}{2}$ \\ (D) $x = 2n\pi - 2 \arctan\frac{1}{2}$ \hfill \textbf{((a))}
\item Q. Solve for $x$: \\ $\sin 2x + 3 \sin x = 1 + 3 \cos x$. \\ (A) $x = n\pi + \frac{\pi}{4}$, $n \in \mathbb{Z}$ \\ (B) $x = 2n\pi + \frac{\pi}{4}$, $n \in \mathbb{Z}$ \\ (C) $x = n\pi + \frac{3\pi}{4}$, $n \in \mathbb{Z}$ \\ (D) $x = n\pi + \frac{\pi}{3}$, $n \in \mathbb{Z}$ \hfill \textbf{((a))}
\item Q. Solve the equation \\ $\sin^2 x - \frac{\cos x}{4} = \frac{1}{4}$. \\ (A) $x = 2n\pi \pm \cos^{-1}\frac{3}{4}$, $n \in \mathbb{Z}$. \\ (B) $x = (2n + 1)\pi$, $n \in \mathbb{Z}$. \\ (C) $x = (2n + 1)\pi$ or $x = 2n\pi \pm \cos^{-1}\frac{3}{4}$, $n \in \mathbb{Z}$. \\ (D) $x = n\pi \pm \cos^{-1}\frac{3}{4}$, $n \in \mathbb{Z}$. \hfill \textbf{((c))}
\item Q. Solve for $x$: \\ $\frac{\sin^3\left(\frac{x}{2}\right) - \cos^3\left(\frac{x}{2}\right)}{2 + \sin x} = \frac{\cos x}{3}$. \\ (A) $x = 2n\pi + \frac{\pi}{2}$, $n \in \mathbb{Z}$. \\ (B) $x = 2n\pi + \pi$, $n \in \mathbb{Z}$. \\ (C) $x = 2n\pi + \frac{\pi}{4}$, $n \in \mathbb{Z}$. \\ (D) $x = n\pi$, $n \in \mathbb{Z}$. \hfill \textbf{((a))}
\item Q. For a real exponent x, define $f(\theta) = \frac{1}{1 + (\cot \theta)^x} = \frac{\sin^x \theta}{\sin^x \theta + \cos^x \theta}$, $\theta \in \{1^\circ, 2^\circ, \dots, 89^\circ\}$. Let $S = \sum_{\theta=1^\circ}^{89^\circ} f(\theta)$, and let [S] denote the greatest integer $\leq$ S. Then [S] equals \hfill \textbf{((B))}
\end{enumerate}
\subsection{Hard}
\begin{enumerate}
\item Q. Sum of all the solutions in $[0, 4\pi]$ of the equation $\tan x + \cot x + 1 = \cos(x + \frac{\pi}{4})$ is
(A) $3\pi$
(B) $\frac{\pi}{2}$
(C) $\frac{7\pi}{2}$
(D) $4\pi$ \hfill \textbf{((C))}
\end{enumerate}
\section{Unsorted}
\subsection{Medium}
\begin{enumerate}
\item A triangle has two vertices $A = (0, 1)$, $B = (2, 3)$, and its centroid is at $G = (\frac{5}{3},\frac{3}{3})$. Find the third vertex $C$, the circumcenter $P$, the circumradius $R$, and the orthocenter $H$. \hfill \textbf{((A))}
\item A rigid rod of length $l$ has its ends always sliding on the $x$- and $y$-axes, respectively. Let $P(x, y)$ be the point on the rod that divides it in the ratio $m_1 : m_2$ (measured from the $x$-axis end towards the $y$-axis end). The locus of $P$ is \hfill \textbf{((C))}
\item Let $A = (a,0)$ and $B = (-a, 0)$ be fixed points. A variable point $C = (x, y)$ moves so that $\cot \angle A + \cot \angle B = \lambda$, where $\lambda$ is a constant. The locus of $C$ is \hfill \textbf{((A))}
\item Two points A and B move along the positive x- and y-axes respectively so that OA+OB = K (a constant). Let M = (x, y) be the foot of the perpendicular from the origin O onto the line AB. The locus of M is \hfill \textbf{((C))}
\item A variable line through P(-2, -3) meets the pair of lines $x^2 + 3y^2 + 4xy – 8x – 6y – 9 =0$ at points Q and R. If PQ $\cdot$ PR = 20, then the equations of the line are \hfill \textbf{((A))}
\item Find the equations of the two legs of an isosceles right-angled triangle whose hypotenuse is $3x + 4y = 4$ and whose opposite vertex is $C = (2, 2)$. \hfill \textbf{((A))}
\item A circle of radius 3 has its center on the line $y = x - 1$ and passes through the point (7,3). Which of the following gives the equations of all such circles? \hfill \textbf{((A))}
\item A variable line $y - 2x + a = 0$ lies entirely between (and does not meet or touch) the two circles $C_1: x^2 + y^2 - 2x - 2y + 1 = 0$ and $C_2: x^2 + y^2 – 16x – 2y + 61 = 0$. Find the range of the parameter $a$. \hfill \textbf{((B))}
\item Find the equation of the circle whose center is (3, -1) and which cuts off a chord of length 6 on the line $2x – 5y + 18 = 0$. \hfill \textbf{((A))}
\item The straight line $ax + by = 2$, $a, b \neq 0$, is tangent to the circle $x^2 + y^2 - 2x = 3$ and is normal to the circle $x^2 + y^2 - 4y = 6$. Find $(a, b)$. \hfill \textbf{((C))}
\item From the point (14,7), three normals are drawn to the parabola $y^2 - 16x - 8y = 0$. The feet of these normals are \hfill \textbf{((A))}
\item Let $y^2 = 4ax$ be a parabola. A chord of this parabola subtends a right angle at the vertex. If $(h, k)$ is the midpoint of such a chord, then its locus is \hfill \textbf{((A))}
\item The parabolas $y^2 = 4ax$ and $y^2 = 4c (x - b)$, with $a \neq c$, admit a common normal (other than the x-axis). Which of the following is the correct necessary and sufficient condition? \hfill \textbf{((A))}
\item A line $lx + my = n$ is a normal to the ellipse $\frac{x^2}{a^2} + \frac{y^2}{b^2} = 1$ if and only if its coefficients satisfy which of the following conditions? \hfill \textbf{((A))}
\item A normal to the hyperbola $\frac{x^2}{a^2} - \frac{y^2}{b^2} = 1$ meets the x- and y-axes at M and N, respectively. From M draw MP $\perp$ x-axis, and from N draw NP $\perp$ y-axis, meeting at P(x, y). The locus of P is \hfill \textbf{((A))}
\item Q. f is a continuous function on R. Given that
$$x^2 + (f(x) – 2) x - \sqrt{3} f(x) + 2\sqrt{3}-3 = 0 \quad \forall x,$$then the value of $f(\sqrt{3})$ is
(A) $\frac{2(\sqrt{3}-2)}{\sqrt{3}}$
(B) $2(1 - \sqrt{3})$
(C) $0$
(D) cannot be determined \hfill \textbf{(B)}
\item Q. If
$$x^y + y^x = 2,$$then
$$\frac{dy}{dx} =?$$
(A) $\frac{y^x \ln y + x^y \frac{y}{x}}{x^y \ln x + y^x \frac{x}{y}}$
(B) $\frac{y^x \ln x + x^y \frac{y}{x}}{x^y \ln y + y^x \frac{x}{y}}$
(C) $\frac{x^y \ln y + y^x \frac{x}{y}}{y^x \ln x + x^y \frac{y}{x}}$
(D) $\frac{x^y \ln x + y^x \frac{y}{x}}{y^x \ln y + x^y \frac{x}{y}}$ \hfill \textbf{(A)}
\item Q. Suppose $y$ is defined by
$$y = \frac{\sin x}{1 + \frac{\cos x}{1 + \frac{\sin x}{1 + \frac{\cos x}{1 + \frac{\sin x}{1 + \frac{\cos x}{1+...}}}}}}$$Then
$$\frac{dy}{dx} =?$$
(A) $\frac{(1+y) \cos x + y \sin x}{1 + 2y + \cos x - \sin x}$
(B) $\frac{(1 + y) \sin x + y \cos x}{1 + 2y + \sin x - \cos x}$
(C) $\frac{(1 + y) \cos x - y \sin x}{1 + 2y - \cos x - \sin x}$
(D) $\frac{(1 + y) \sin x - y \cos x}{1 + 2y - \sin x + \cos x}$ \hfill \textbf{(A)}
\item Q. A curve is defined parametrically by
$$x = \frac{1+t}{t^3}, \quad y= \frac{3}{2t^2} + \frac{2}{t}$$If $y' = \frac{dy}{dx}$, which of the following differential relations does the curve satisfy?
(A) $x (y')^3 = 1 + y'$
(B) $x (y')^3 = 1 - y'$
(C) $x (y')^2 = 1 + y'$
(D) $x (y')^3 = y' – 1$ \hfill \textbf{(A)}
\item Q. Suppose $f$ is a function satisfying
$$f(x) = x^3 + x^2 f'(1) + x f''(2) + f'''(3) \quad \text{for all } x \in \mathbb{R}.$$Then $f(x)$ can be written in the form $x^3 + ax^2 + bx + c$ with constants $a, b, c$. Determine $f(x)$ explicitly (i.e. find $a, b, c$).
(A) $x^3 - 5x^2 + 2x + 6$
(B) $x^3 + 3x^2 - 4x + 6$
(C) $x^3 - 6x^2 + 3x + 1$
(D) $x^3 - x^2 + x - 1$ \hfill \textbf{(A)}
\item Q. If
$$x = a(t + \sin t), \quad y = a(1-\cos t),$$then
$$\frac{d^2y}{dx^2} =?$$
(A) $\frac{1}{2a}\sec^2\left(\frac{t}{2}\right)$
(B) $\frac{1}{2a}\sec^4\left(\frac{t}{2}\right)$
(C) $\frac{1}{4a}\sec^2\left(\frac{t}{2}\right)$
(D) $\frac{1}{4a}\sec^4\left(\frac{t}{2}\right)$ \hfill \textbf{(D)}
\item Q. Let $y = f(x)$ with $f$ continuously differentiable and let $x = g(y)$ be its inverse. Then $g'(y)$ and $g''(y)$ are given by
$$g'(y) = \dots, \quad g''(y) = \dots$$Which of the following pairs is correct?
(A) $g'(y) = \frac{1}{f'(x)}, \quad g''(y) = -\frac{f''(x)}{(f'(x))^3}$
(B) $g'(y) = \frac{1}{f'(x)}, \quad g''(y) = \frac{f''(x)}{(f'(x))^2}$
(C) $g'(y) = f'(x), \quad g''(y) = -\frac{f''(x)}{(f'(x))^3}$
(D) $g'(y) = f'(x), \quad g''(y) = \frac{f''(x)}{(f'(x))^2}$ \hfill \textbf{(A)}
\item Q. Evaluate
$$\lim_{n\to\infty} \left(\frac{e^n}{\pi}\right)^{1/n}.$$The value of this limit is
(A) $e$
(B) $\pi$
(C) $1$
(D) none of these \hfill \textbf{(A)}
\item Q. Evaluate
$$\lim_{x\to 0^+} \log_{\sin x} (\sin 2x).$$The value of this limit is
(A) $2$
(B) $1$
(C) $\frac{1}{2}$
(D) $0$ \hfill \textbf{(B)}
\item Q. Let
$$f(x) = \begin{vmatrix} x & x^2 & x^3 \\ 1 & 2x & 3x^2 \\ 0 & 2 & 6x \end{vmatrix}$$Then $f'(x) =$
(A) $0$
(B) $x^2$
(C) $6x^2$
(D) $6x^3$ \hfill \textbf{(C)}
\item Q. Determine the largest possible area of a rectangle with its base on the x-axis and the other two vertices on the curve $y = e^{-x^2}$. Writing the half-width of the rectangle as $a > 0$, its vertices are $(\pm a, 0)$ and $(\pm a, e^{-a^2})$, so its area is
$$A(a) = (\text{base}) \times (\text{height}).$$The maximum value of $A(a)$ is
(A) $\frac{1}{e}$
(B) $\frac{2}{e}$
(C) $\frac{\sqrt{2}}{\sqrt{e}}$
(D) $\frac{1}{\sqrt{2}e}$ \hfill \textbf{(C)}
\item Q. A conical vessel of maximum volume is to be made from a circular sheet of gold of radius $1$ by removing a sector and forming the remaining part into the cone. The area of the sector removed is
(A) $\pi$
(B) $\pi\left(\frac{2}{\sqrt{3}}-1\right)$
(C) $\pi\left(1 - \sqrt{3}\right)$
(D) $2\pi \left(1 - \frac{2}{\sqrt{3}}\right)$ \hfill \textbf{(C)}
\item Q. Find the coordinates of the point on the curve
$$y^2 = 8x$$that is at minimum distance from the line
$$x + y + 2 = 0.$$
(A) $(2, -4)$
(B) $(-2, -4)$
(C) $(2, 4)$
(D) $(-4, 2)$ \hfill \textbf{(A)}
\item Q. A right circular cylinder is inscribed in a cone of base radius $r$ and height $h$. Find the radius $x$ and height $y$ of the cylinder that maximize its volume.
(A) $x = \frac{r}{2}, \quad y = \frac{h}{2}$
(B) $x = \frac{2r}{3}, \quad y = \frac{h}{3}$
(C) $x = \frac{r}{3}, \quad y = \frac{2h}{3}$
(D) $x = \frac{3r}{4}, \quad y = \frac{h}{4}$ \hfill \textbf{(B)}
\item Q. Let
$$f(x) = x^2 - \ln |x|, \quad x \neq 0.$$On which of the following intervals is $f$ increasing?
(A) $\left(-\infty, -\frac{1}{\sqrt{2}}\right) \cup \left(0, \frac{1}{\sqrt{2}}\right)$
(B) $\left(-\frac{1}{\sqrt{2}}, 0\right) \cup \left(0, \frac{1}{\sqrt{2}}\right)$
(C) $\left(-\infty, -\frac{1}{\sqrt{2}}\right) \cup \left(\frac{1}{\sqrt{2}}, \infty\right)$
(D) $\left(-\frac{1}{\sqrt{2}}, 0\right) \cup \left(\frac{1}{\sqrt{2}}, \infty\right)$ \hfill \textbf{(D)}
\item Q. Evaluate I = \int \frac{x^2 - 1}{(x^4 + 3x^2 + 1) \tan^{-1}(x + \frac{1}{x})} dx.
The value of I is
(A) \ln|\tan^{-1}(x+\frac{1}{x})| + C
(B) \ln|x + \frac{1}{x}| + C
(C) \tan^{-1}(x + \frac{1}{x}) + C
(D) \frac{1}{x^2 + 1} + C \hfill \textbf{(A)}
\item Q. Evaluate I = \int \frac{1- \sqrt{x}}{1 + \sqrt{x}} dx.
The result is
(A) -2\ln|\sec\theta - \tan\theta| + 2\theta + C, \theta = \cos^{-1} \sqrt{x}.
(B) -2\ln|\frac{1+\sqrt{1-x}}{x}| + 2\cos^{-1} \sqrt{x} + C.
(C) 2\ln|\frac{1+\sqrt{x}}{1-\sqrt{x}}| - 2\sin^{-1} \sqrt{x} + C.
(D) -\ln|\frac{1-\sqrt{1-x}}{x}| + \theta + C, \theta = \sin^{-1} \sqrt{x}. \hfill \textbf{(B)}
\item Q. Evaluate I = \int \cos(\sqrt{x})dx.
The correct expression for I is
(A) 2[\sqrt{x} \sin\sqrt{x} - \cos\sqrt{x}] + C
(B) 2[\sqrt{x} \sin\sqrt{x} + \cos\sqrt{x}] + C
(C) \sqrt{x} \sin\sqrt{x} + \cos\sqrt{x} + C
(D) 2[\sqrt{x} \cos\sqrt{x} - \sin\sqrt{x}] + C \hfill \textbf{(B)}
\item Q. Evaluate the integral I = \int e^x \frac{x^4 + 2}{(1+x^2)^{5/2}} dx.
The value of I is
(A) \frac{e^x (x + 1)}{(1+x^2)^{3/2}} + C
(B) \frac{e^x (1 - x + x^2)}{(1+x^2)^{3/2}} + C
(C) \frac{e^x (1 - x)}{(1+x^2)^{3/2}} + C
(D) none of these \hfill \textbf{(D)}
\item Q. Evaluate I = \int \frac{4}{\sin^4 x + \cos^4 x} dx.
A convenient closed-form is
(A) 2\sqrt{2} \tan^{-1}(\frac{\tan x - \cot x}{\sqrt{2}}) + C
(B) 2\sqrt{2} \tan^{-1}(\frac{\tan x + \cot x}{\sqrt{2}}) + C
(C) 2\sqrt{2} \sinh^{-1}(\frac{\tan x - \cot x}{\sqrt{2}}) + C
(D) \frac{2}{\sqrt{2}} \ln|\frac{\tan x - \cot x + \sqrt{2}}{\tan x - \cot x - \sqrt{2}}| + C \hfill \textbf{(A)}
\item Q. Evaluate I = \frac{1}{2} \int \frac{2}{x^4 + 5x^2 + 1} dx.
The value of I can be written in closed form as
(A) \frac{1}{\sqrt{7}} [\tan^{-1}(\frac{x - 1/x}{\sqrt{7}}) - \frac{1}{\sqrt{3}} \tan^{-1}(\frac{x + 1/x}{\sqrt{3}})] + C
(B) \frac{1}{\sqrt{7}} [\tan^{-1}(\frac{x + 1/x}{\sqrt{7}}) - \frac{1}{\sqrt{3}} \tan^{-1}(\frac{x - 1/x}{\sqrt{3}})] + C
(C) \frac{1}{\sqrt{3}} [\tan^{-1}(\frac{x - 1/x}{\sqrt{3}}) - \frac{1}{\sqrt{7}} \tan^{-1}(\frac{x + 1/x}{\sqrt{7}})] + C
(D) none of these \hfill \textbf{(A)}
\item Q. Evaluate I = \int \frac{2+3 \cos \theta}{\sin \theta + 2 \cos \theta+3} d\theta.
The value of I is
(A) \frac{6\theta}{5} + \frac{3}{5} \ln|\sin \theta + 2 \cos \theta+3| - \frac{8}{5} \tan^{-1}(\frac{\tan(\theta/2) + 1}{\sqrt{2}}) + C
(B) \frac{6\theta}{5} - \frac{3}{5} \ln|\sin \theta + 2 \cos \theta+3| - \frac{8}{5} \tan^{-1}(\frac{\tan(\theta/2) + 1}{\sqrt{2}}) + C
(C) \frac{6\theta}{5} + \frac{3}{5} \ln|\sin \theta + 2 \cos \theta+3| + \frac{8}{5} \tan^{-1}(\frac{\tan(\theta/2) + 1}{\sqrt{2}}) + C
(D) none of these \hfill \textbf{(A)}
\item Q. Evaluate the integral \int \sin^2 x \cos^4 x dx.
The correct antiderivative is
(A) \frac{\sin 6x}{192} - \frac{\sin 4x}{64} + \frac{\sin 2x}{64} + \frac{x}{16} + C
(B) \frac{\sin 6x}{192} - \frac{\sin 4x}{32} + \frac{\sin 2x}{64} + \frac{x}{16} + C
(C) \frac{\sin 6x}{192} - \frac{\sin 4x}{64} + \frac{\sin 2x}{32} + \frac{x}{16} + C
(D) none of these \hfill \textbf{(A)}
\item Q. Evaluate the integral \int \frac{\sqrt{\sin x}}{\cos^{9/2} x} dx.
(A) \frac{2}{3}\tan^{3/2} x + \frac{2}{7}\tan^{7/2} x + C
(B) \frac{2}{3}\tan^{3/2} x + \frac{1}{7}\tan^{7/2} x + C
(C) \frac{3}{2}\tan^{3/2} x + \frac{2}{7}\tan^{7/2} x + C
(D) none of these \hfill \textbf{(A)}
\item Q. Evaluate the integral I = \int_1^2 (x[x^2] + [x^2]x) dx,
where [.] is the greatest-integer function. The value of I is
(A) \frac{5}{4} + \sqrt{3} + \frac{(2\sqrt{3}-2\sqrt{2})}{\ln 2} + \frac{1}{\ln 3} (9-3\sqrt{3})
(B) \frac{5}{4} + \sqrt{3} + \frac{(2\sqrt{3}-\sqrt{2})}{\ln 2} + \frac{1}{\ln 3} (9-3\sqrt{3})
(C) \frac{\sqrt{5}}{4} + \sqrt{3} + \frac{1}{\ln 2} (2\sqrt{3}-2\sqrt{2}) + \frac{1}{\ln 3} (9-3\sqrt{3})
(D) none of these \hfill \textbf{(B)}
\item Q. Let f(x) = \begin{vmatrix} \cos x & e^{x^2} & 2x \cos^2\frac{x}{2} \\ x^2 & \sec x & \sin x + x^3 \\ 1 & 2 & x + \tan x \end{vmatrix}
Evaluate \int_{-\pi/2}^{\pi/2} (x^2 + 1)(f(x) + f''(x)) dx.
The value is
(A) 1
(B) -1
(C) 2
(D) none of these \hfill \textbf{(D)}
\item Q. Evaluate I = \int_{-\pi}^{\pi} \frac{x \sin x}{e^x + 1} dx.
The value of I is
(A) 0
(B) \pi
(C) 2\pi
(D) -\pi \hfill \textbf{(B)}
\item Q. Evaluate the integral \int_0^1 \cot^{-1}(1 - x + x^2) dx.
The value of this integral is
(A) \frac{\pi}{2} + \ln 2
(B) \frac{\pi}{2} - \ln 2
(C) \pi - \ln 2
(D) none of these \hfill \textbf{(B)}
\item Q. Find the area enclosed between the curve y = (x - 1)(x - 2)(x - 3) and the x-axis, from x = 0 to x = 3. This area equals
(A) \frac{7}{4}
(B) \frac{9}{4}
(C) \frac{11}{4}
(D) \frac{13}{4} \hfill \textbf{(C)}
\item Q. Find the area of the region common to the circle x^2 + y^2 + 4x + 6y - 3 = 0 and the parabola x^2 + 4x = 6y + 14. This area equals
(A) \frac{4\sqrt{3} + 16\pi}{3}
(B) \frac{4\sqrt{3} + 8\pi}{3}
(C) \frac{8\sqrt{3} + 16\pi}{3}
(D) \frac{4\sqrt{3} + 16\pi}{6} \hfill \textbf{(A)}
\item Q. Evaluate the infinite sum
$\cot^{-1}(1^2+\frac{3}{4})+\cot^{-1}(2^2+\frac{3}{4})+\cot^{-1}(3^2+\frac{3}{4})+...$ \hfill \textbf{((B))}
\item Q. Evaluate the integral
$\int \frac{dx}{x^{2008}+x}$
in the form
$\frac{1}{p} \ln \left(\frac{x^q}{1+x^r}\right)+C$,
where $p, q, r \in \mathbb{N}$. Then find $p + q + r$. \hfill \textbf{((C))}
\item Q. Suppose real numbers $x_i, y_i$ satisfy
$\sum_{i=1}^6(\sin^{-1}(x_i)+\cos^{-1}(y_i))=9\pi$.
Evaluate the integral
$\int_{-6}^6 x \ln(1+x^2) \frac{e^x}{1+e^{2x}} dx$. \hfill \textbf{((A))}
\item Q. Let
$\vec{r} = (\vec{a} \times \vec{b}) \sin x + (\vec{b} \times \vec{c}) \cos y + 2(\vec{c} \times \vec{a}),$
where $\vec{a}, \vec{b}, \vec{c}$ are nonzero, noncoplanar vectors. If $\vec{r}$ is orthogonal to $\vec{a} + \vec{b} + \vec{c}$, then the
minimum value of
$\frac{20}{\pi^2}(x^2+y^2)$
is \hfill \textbf{((A))}
\item Q. Let $x, y, z$ be positive real numbers with $x+y+z=10$. The maximum value of
$x^2 y^3 z^5$
is $N$. If $k$ is the number of odd positive divisors of $N$, then $k/12$ equals \hfill \textbf{((D))}
\item Q. Define
$f(x) = \lim_{n\to\infty} \frac{\sin(\pi x^4)+(x+2)^n \frac{\tan(\pi x)}{x+1}}{1+(x+2)^n-x^4}$.
If $m = \lim_{x\to-1} f(x)$, then $[m/\pi]$ equals \hfill \textbf{((D))}
\item Q. Let $m$ and $n$ be the coefficients of $x^3$ in the expansions of
$(1+x+2x^2+3x^3)^3$ and $(1+x+2x^2+3x^3+4x^4)^3$,
respectively. Then $[m/n]$ equals \hfill \textbf{((B))}
\item Q. Let $\alpha, \beta \in \mathbb{R}$ and define
$f(x) = \alpha \sin x + \beta \sqrt{|x|} + 4, \quad x \in \mathbb{R}$.
If
$f(\log_{10}(\log_3 10)) = 5$,
then the value of
$f(\log_{10}(\log_{10} 3))$
is \hfill \textbf{((A))}
\item Q. Let $a, b, c \in \mathbb{R}$ satisfy
$a^2+4b^2+9c^2+2ab+6bc+3ac \le 0$.
How many distinct real roots does the quadratic equation
$ax^2+bx+c=0$
have? \hfill \textbf{((D))}
\item Q. Let $z$ be a complex number satisfying
$(z+\frac{1}{z})(z+\frac{1}{z}+1)=1$.
Compute the value of
$(3z^{100}+\frac{2}{z^{100}}+1)(z^{100}+\frac{2}{z^{100}}+3)$. \hfill \textbf{((A))}
\item Q. Let $A$ be a $2 \times 2$ real matrix satisfying
$\operatorname{tr} A = 5$ and $\det A = 6$.
Then
$\det(A+I_2)$
equals \hfill \textbf{((C))}
\item Q. Define
$f(\theta) = \frac{1}{1+(\cot \theta)^x}$
and let
$S = \sum_{\theta=1^\circ}^{89^\circ} f(\theta)$.
Then $[S]$ is \hfill \textbf{((A))}
\item Q. If
$(\frac{t^3}{t-1}, \frac{t^2-3}{t-1})$ and $(\frac{t}{c-1}, \frac{c^2-3}{c-1})$
are collinear and
$\alpha abc + \beta (a+b+c) = \gamma (ab+bc+ca),$
where $\alpha, \beta, \gamma \in \mathbb{N}$ are coprime, then the value of
$[\alpha + \beta + \gamma]$
(where $[.]$ denotes the greatest-integer function) is \hfill \textbf{((B))}
\item Q. A ray of light moving parallel to the $x$-axis is reflected from the parabolic mirror
$(y-2)^2 = 4(x+1)$.
After reflection, the ray must pass through a point on the axis of the parabola of the form
$(0, k)$, where $k \in \mathbb{N}$. The value of $k$ is \hfill \textbf{((B))}
\item Q. Let
$S_1: x^2+y^2+4x-6y-12=0, \quad S_2: x^2+y^2-3x-2y+1=0$
be two distinct circles with centres $C_1$ and $C_2$, respectively. If $A$ and $B$ are two points which
lie on their common chord as well as on their common tangents, and
$C_1A^2-C_2B^2=\frac{a}{b}$,
where $a, b \in \mathbb{N}$ are coprime, then the value of $a+b$ is \hfill \textbf{((C))}
\item Q. $f$ is a continuous function on $\mathbb{R}$. Given that
$x^2 + (f(x) – 2) x - \sqrt{3} f(x) + 2\sqrt{3}-3 = 0 \quad \forall x$,
then the value of $f(\sqrt{3})$ is
(A) $\frac{2(\sqrt{3}-2)}{\sqrt{3}}$
(B) $2(1 - \sqrt{3})$
(C) $0$
(D) cannot be determined \hfill \textbf{((B))}
\item Q. If
$x^y + y^x = 2$,
then $\frac{dy}{dx} =?$
(A) $\frac{y^x \ln y + x y^{\frac{y}{x}-1}}{x^y \ln x + y x^{\frac{x}{y}-1}}$
(B) $\frac{y^x \ln x + x y^{\frac{y}{x}-1}}{x^y \ln y + y x^{\frac{x}{y}-1}}$
(C) $\frac{x^y \ln y + y x^{\frac{x}{y}-1}}{y^x \ln x + x y^{\frac{y}{x}-1}}$
(D) $\frac{x^y \ln x + y x^{\frac{x}{y}-1}}{y^x \ln y + x y^{\frac{x}{y}-1}}$ \hfill \textbf{((A))}
\item Q. Suppose $y$ is defined by
$y = \frac{\sin x}{1+\frac{\cos x}{1+\frac{\sin x}{1+\frac{\cos x}{1+\frac{\sin x}{1+\ldots}}}}}$
Then $\frac{dy}{dx} =$
(A) $\frac{(1+y) \cos x + y \sin x}{1 + 2y + \cos x - \sin x}$
(B) $\frac{(1 + y) \sin x + y \cos x}{1 + 2y + \sin x - \cos x}$
(C) $\frac{(1 + y) \cos x - y \sin x}{1+2y - \cos x - \sin x}$
(D) $\frac{(1 + y) \sin x - y \cos x}{1+2y- \sin x + \cos x}$ \hfill \textbf{((A))}
\item Q. A curve is defined parametrically by
$x = \frac{1+t}{t^3}$, $y=\frac{3}{2t^2} + \frac{2}{t}$
If $y' = \frac{dy}{dx}$, which of the following differential relations does the curve satisfy?
(A) $x (y')^3 = 1 + y'$
(B) $x (y')^3 = 1 - y'$
(C) $x (y')^2 = 1 + y'$
(D) $x (y')^3 = y' – 1$ \hfill \textbf{((A))}
\item Q. Suppose $f$ is a function satisfying
$f(x) = x^3 + x^2 f'(1) + x f''(2) + f'''(3)$ for all $x \in \mathbb{R}$.
Then $f(x)$ can be written in the form $x^3 + ax^2 + bx + c$ with constants $a, b, c$. Determine
$f(x)$ explicitly (i.e. find $a, b, c$).
(A) $x^3 - 5x^2 + 2x + 6$
(B) $x^3 + 3x^2 - 4x + 6$
(C) $x^3 - 6x^2 + 3x + 1$
(D) $x^3 - x^2 + x - 1$ \hfill \textbf{((A))}
\item Q. If
$x = a(t + \sin t)$, $y = a(1-\cos t)$,
then $\frac{d^2y}{dx^2} =?$
(A) $\frac{1}{2a} \sec^2(\frac{t}{2})$
(B) $\frac{1}{2a} \sec^4(\frac{t}{2})$
(C) $\frac{1}{4a} \sec^2(\frac{t}{2})$
(D) $\frac{1}{4a} \sec^4(\frac{t}{2})$ \hfill \textbf{((D))}
\item Q. Let $y = f(x)$ with $f$ continuously differentiable and let $x = g(y)$ be its inverse. Then
$g'(y)$ and $g''(y)$ are given by
$g'(y) = \ldots, \quad g''(y) = \ldots$
Which of the following pairs is correct?
(A) $g'(y) = \frac{1}{f'(x)}, \quad g''(y) = -\frac{f''(x)}{(f'(x))^3}$
(B) $g'(y) = \frac{1}{f'(x)}, \quad g''(y) = \frac{f''(x)}{(f'(x))^2}$
(C) $g'(y) = f'(x), \quad g''(y) = -\frac{f''(x)}{(f'(x))^3}$
(D) $g'(y) = f'(x), \quad g''(y) = \frac{f''(x)}{(f'(x))^2}$ \hfill \textbf{((A))}
\item Q. Evaluate
$\lim_{n\to\infty} \left(\frac{e^n}{\pi}\right)^{1/n}$.
The value of this limit is
(A) $e$
(B) $\pi$
(C) $1$
(D) none of these \hfill \textbf{((A))}
\item Q. Evaluate
$\lim_{x\to0^+} \log_{\sin x} (\sin 2x)$.
The value of this limit is
(A) $2$
(B) $1$
(C) $\frac{1}{2}$
(D) $0$ \hfill \textbf{((B))}
\item Q. Let
$f(x) = \begin{vmatrix} x & x^2 & x^3 \\ 1 & 2x & 3x^2 \\ 0 & 2 & 6x \end{vmatrix}$
Then $f'(x) =$
(A) $0$
(B) $x^2$
(C) $6x^2$
(D) $6x^3$ \hfill \textbf{((C))}
\item Q. Determine the largest possible area of a rectangle with its base on the $x$-axis and the
other two vertices on the curve $y = e^{-x^2}$. Writing the half-width of the rectangle as $a > 0$,
its vertices are $(\pm a, 0)$ and $(\pm a, e^{-a^2})$, so its area is
$A(a) = (\text{base}) \times (\text{height})$.
The maximum value of $A(a)$ is
(A) $\frac{1}{e}$
(B) $\frac{2}{e}$
(C) $\frac{\sqrt{2}}{\sqrt{e}}$
(D) $\frac{1}{\sqrt{2}e}$ \hfill \textbf{((C))}
\item Q. A conical vessel of maximum volume is to be made from a circular sheet of gold of radius
$1$ by removing a sector and forming the remaining part into the cone. The area of the sector
removed is
(A) $\pi \sqrt{\frac{2}{3}}$
(B) $\pi (\frac{2}{3}-1)$
(C) $\pi(1 - \sqrt{\frac{2}{3}})$
(D) $2\pi (1 - \sqrt{\frac{2}{3}})$ \hfill \textbf{((C))}
\item Q. Find the coordinates of the point on the curve
$y^2 = 8x$
that is at minimum distance from the line
$x + y + 2 = 0$.
(A) $(2, -4)$
(B) $(-2, -4)$
(C) $(2, 4)$
(D) $(-4, 2)$ \hfill \textbf{((A))}
\item Q. A right circular cylinder is inscribed in a cone of base radius $r$ and height $h$. Find the
radius $x$ and height $y$ of the cylinder that maximize its volume.
(A) $x = \frac{r}{2}, y = \frac{h}{2}$
(B) $x = \frac{2r}{3}, y = \frac{h}{3}$
(C) $x = \frac{r}{3}, y = \frac{2h}{3}$
(D) $x = \frac{3r}{4}, y = \frac{h}{4}$ \hfill \textbf{((B))}
\item Q. Let
$f(x) = x^2 - \ln |x|, \quad x \neq 0$.
On which of the following intervals is $f$ increasing?
(A) $(-\infty, -\frac{1}{\sqrt{2}}) \cup (0, \frac{1}{\sqrt{2}})$
(B) $(-\frac{1}{\sqrt{2}}, 0) \cup (0, \frac{1}{\sqrt{2}})$
(C) $(-\infty, -\frac{1}{\sqrt{2}}) \cup (\frac{1}{\sqrt{2}}, \infty)$
(D) $(-\frac{1}{\sqrt{2}}, 0) \cup (\frac{1}{\sqrt{2}}, \infty)$ \hfill \textbf{((D))}
\item Q. A fair coin is tossed three times. Define the events:
A: no head appears, B: exactly one head appears,
C: at least two heads appear.
Which of the following is true?
(A) A, B, C are mutually exclusive but not exhaustive.
(B) A, B, C are exhaustive but not mutually exclusive.
(C) A, B, C are neither mutually exclusive nor exhaustive.
(D) A, B, C are both mutually exclusive and exhaustive. \hfill \textbf{((D))}
\item Q. n positive integers are chosen at random (each last digit equally likely 0, 1, 2, ..., 9). Let
P5 be the probability that the last digit of their product is 5, and Po the probability that it
is 0. Which of the following gives the correct expressions?
(A) P5 = \frac{5^n - 4^n}{10^n}, Po = \frac{10^n - 8^n - 5^n + 4^n}{10^n}
(B) P5 = \frac{4^n - 3^n}{10^n}, Po = \frac{10^n - 7^n - 4^n + 3^n}{10^n}
(C) P5 = \frac{6^n - 5^n}{10^n}, Po = \frac{10^n - 9^n - 6^n + 5^n}{10^n}
(D) none of these \hfill \textbf{((A))}
\item Q. A bag contains n white and n red balls. Pairs of balls are drawn at random without
replacement until the bag is empty. What is the probability that each of the n drawn pairs
consists of one white and one red ball?
(A) \frac{2^n}{\binom{2n}{n}}
(B) \frac{\binom{2n}{n}}{2^n}
(C) \frac{(n!)^2}{(2n)!}
(D) \frac{1}{\binom{2n}{n}} \hfill \textbf{((A))}
\item Q. Let A and B be two events. The odds against A are 2 : 1, and the odds in favor of A\cup B
are 3 : 1. Then P(B) satisfies
what is the range of P(B)?
(A) $0 \le P(B) \le \frac{1}{4}$
(B) $\frac{1}{12} \le P(B) \le \frac{1}{4}$
(C) $\frac{5}{12} \le P(B) \le \frac{3}{4}$
(D) $\frac{5}{12} \le P(B) \le 1$ \hfill \textbf{((C))}
\item Q. Three numbers are chosen at random without replacement from the set \{1, 2, 3, ..., 10\}.
What is the probability that either the minimum of the three chosen numbers is 4 or the
maximum is 8?
(A) \frac{11}{40}
(B) \frac{3}{10}
(C) \frac{1}{40}
(D) none of these \hfill \textbf{((A))}
\item Q. Two dice are thrown. Find the probability that the numbers appeared have a sum of 8
if it is known that the second die always exhibits 4.
(A) \frac{1}{4}
(B) \frac{1}{6}
(C) \frac{1}{5}
(D) \frac{1}{3} \hfill \textbf{((B))}
\item Q. A lot contains 50 defective and 50 non-defective bulbs. Two bulbs are drawn at random,
one at a time, with replacement. The events are defined as:
A = \{first bulb is defective\}, B = \{second bulb is non-defective\},
C = \{both bulbs are either both defective or both non-defective\}.
Determine which of the following is true:
(A) A, B, C are pairwise independent and also mutually independent.
(B) A, B, C are pairwise independent but not mutually independent.
(C) A, B, C are not pairwise independent but are mutually independent.
(D) A, B, C are neither pairwise independent nor mutually independent. \hfill \textbf{((B))}
\item Q. The probability that an anti-aircraft gun can hit an enemy plane at the first, second and
third shot are 0.6, 0.7 and 0.1 respectively. The probability that the gun hits the plane (i.e.,
at least one shot hits) is
(A) 0.108
(B) 0.892
(C) 0.14
(D) none of these \hfill \textbf{((B))}
\item Q. A pair of dice is rolled together until a sum of either 5 or 7 is obtained. Find the
probability that 5 comes before 7.
(A) \frac{1}{4}
(B) \frac{2}{5}
(C) \frac{3}{7}
(D) \frac{5}{18} \hfill \textbf{((B))}
\item Q. Let A, B, C be three events such that exactly one of A and B occurs with probability
$1 - a$, exactly one of B and C occurs with probability $1 - 2a$, exactly one of C and A occurs
with probability $1 - a$, and
$P(A \cap B \cap C) = a^2$.
Then the probability that at least one of A, B, C occurs, i.e. $P(A \cup B \cup C)$, is
(A) $(a-1)^2 + \frac{1}{2}$
(B) $\frac{3}{2} - a + a^2$
(C) $\frac{1}{2} + a - a^2$
(D) $1 - a + a^2$ \hfill \textbf{((A))}
\item Q. If a fair coin is tossed 10 times, the probabilities of getting (i) exactly six heads, (ii) at
least six heads, (iii) at most six heads are respectively:
(A) $\left(\frac{105}{512}, \frac{193}{512}, \frac{53}{64}\right)$
(B) $\left(\frac{105}{512}, \frac{193}{1024}, \frac{53}{64}\right)$
(C) $\left(\frac{105}{1024}, \frac{193}{512}, \frac{53}{64}\right)$
(D) $\left(\frac{105}{512}, \frac{193}{1024}, \frac{53}{64}\right)$ \hfill \textbf{((A))}
\item Q. A coin is tossed 7 times. Each time, a man calls “head.” The probability that he wins
the toss on more than three occasions is:
(A) \frac{1}{4}
(B) \frac{5}{8}
(C) \frac{1}{2}
(D) none of these \hfill \textbf{((C))}
\item Q. In a test an examinee either guesses, copies, or knows the answer to a multiple-choice
question with four options. The probability that he guesses is $\frac{1}{3}$, and the probability that
he copies is $\frac{1}{6}$. Given that he copied, the probability his answer is correct is $\frac{1}{8}$. If he answers
correctly, what is the probability that he knew the answer?
(A) \frac{12}{25}
(B) \frac{24}{29}
(C) \frac{5}{7}
(D) \frac{3}{4} \hfill \textbf{((B))}
\item Q. Three persons A, B, C cut cards in turn (in that order) from a well-shuffled pack of 52
cards, replacing each card before the next cut. The first to draw a spade wins a prize. What
are their respective chances of winning?
(A) $\frac{16}{37}, \frac{12}{37}, \frac{9}{37}$
(B) $\frac{1}{3}, \frac{1}{3}, \frac{1}{3}$
(C) $\frac{16}{37}, \frac{9}{37}, \frac{12}{37}$
(D) none of these \hfill \textbf{((A))}
\item Q. Let A be a set with n elements. A subset P \subseteq A is chosen at random (each of the $2^n$
subsets equally likely). Then, after returning the elements of P to A, a subset Q \subseteq A is
again chosen at random. What is the probability that P and Q have no common elements?
(A) $\left(\frac{3}{4}\right)^n$
(B) $\left(\frac{1}{2}\right)^n$
(C) $\left(\frac{2}{3}\right)^n$
(D) $\left(\frac{1}{4}\right)^n$ \hfill \textbf{((A))}
\end{enumerate}
\section{Vector Algebra}
\subsection{Medium}
\begin{enumerate}
\item Q. Let $\vec{r} = 3\vec{i} + 2\vec{j} - 5\vec{k}$, $\vec{a} = 2\vec{i} - \vec{j} + \vec{k}$, $\vec{b} = \vec{i} + 3\vec{j} - 2\vec{k}$, $\vec{c} = -2\vec{i} + \vec{j} - 3\vec{k}$. If $\vec{r} = \lambda\vec{a} + \mu\vec{b} + \nu\vec{c}$, then $\lambda + \mu + \nu =$
(A) 5
(B) 6
(C) 7
(D) 8 \hfill \textbf{((B))}
\item Q. Let $\vec{a}$ and $\vec{b}$ be non-collinear vectors. For which value of $x$ are the vectors $\vec{\alpha} = (x - 2)\vec{a} + \vec{b}$ and $\vec{\beta} = (3 + 2x)\vec{a} - 2\vec{b}$ collinear?
(A) $x = \frac{1}{4}$
(B) $x = -\frac{1}{2}$
(C) $x= 2$
(D) $x = -\frac{3}{2}$ \hfill \textbf{((A))}
\item Q. Let $\vec{a}, \vec{b}, \vec{c}, \vec{d}$ be the position vectors of four coplanar points A, B, C, D. If $(\vec{a} - \vec{d}) \cdot (\vec{b} - \vec{c}) = 0$ and $(\vec{b} - \vec{d}) \cdot (\vec{c} - \vec{a}) = 0$, then, with respect to $\triangle ABC$, the point D is the
(A) in-centre
(B) orthocentre
(C) circumcentre
(D) centroid \hfill \textbf{((B))}
\item Q. Let $\vec{a} = 7\vec{i} - 4\vec{j} - 4\vec{k}$ and $\vec{b} = -2\vec{i} - \vec{j} + 2\vec{k}$. A vector $\vec{c}$ is directed along the internal bisector of the angle between $\vec{a}$ and $\vec{b}$, and $|\vec{c}| = 5\sqrt{6}$. Then $\vec{c}$ equals
(A) $\frac{5}{3}(\vec{i}-7\vec{j}+2\vec{k})$
(B) $\frac{5}{3}(5\vec{i}+5\vec{j}+2\vec{k})$
(C) $\frac{5}{3}(\vec{i}+7\vec{j}+2\vec{k})$
(D) none of these \hfill \textbf{((A))}
\item Q. If $[a \times b, b \times c, c \times a] = \lambda [a, b, c]^2$, where $[u, v, w]$ denotes the scalar triple product, then $\lambda$ equals
(A) 0
(B) 3
(C) 1
(D) 2 \hfill \textbf{((C))}
\item Q. Let $a, b, c \in \mathbb{R}$ satisfy $3a + 2b + c = 1$. The least value of $a^2 + b^2 + c^2$ is:
(A) $\frac{1}{14}$
(B) $\frac{1}{7}$
(C) $\frac{3}{14}$
(D) $\frac{2}{7}$ \hfill \textbf{((A))}
\item Q. If the angles between the vectors $\mathbf{v}_1$ and $\mathbf{v}_2$, $\mathbf{v}_2$ and $\mathbf{v}_3$, $\mathbf{v}_3$ and $\mathbf{v}_1$ are $\frac{\pi}{3}$, $\frac{\pi}{4}$ and $\frac{\pi}{6}$ respectively, then the angle $\phi$ that $\mathbf{v}_1$ makes with the plane containing $\mathbf{v}_2$ and $\mathbf{v}_3$ is:
(A) $\cos^{-1}\sqrt{1 - \frac{\sqrt{3}}{2}}$
(B) $\cos^{-1}\sqrt{\frac{\sqrt{3}}{2} - 1}$
(C) $\sin^{-1}\sqrt{1 - \frac{\sqrt{3}}{2}}$
(D) $\sin^{-1}\sqrt{\frac{\sqrt{3}}{2} - 1}$ \hfill \textbf{((D))}
\item Q. ABCD is a parallelogram and $A_1$ and $B_1$ are the midpoint of sides BC and CD, respectively. If $\vec{AA_1} + \vec{AB_1} = \lambda \vec{AC}$, then $\lambda$ is equal to:
(A) $\frac{1}{2}$
(B) 1
(C) $\frac{3}{2}$
(D) 2 \hfill \textbf{((C))}
\item Q. Let $\vec{u}, \vec{v}$ and $\vec{w}$ be such that $|\vec{u}| = 1$, $|\vec{v}| = 2\&|\vec{w}| = 3$ If the projection $\vec{v}$ along $\vec{u}$ is equal to that of $\vec{w}$ along $\vec{u}$ and $\vec{v}, \vec{w}$ are perpendicular to each other, then $|\vec{u} - \vec{v} + \vec{w}|$ equal to
(A) $\sqrt{7}$
(B) $2$
(C) $\sqrt{14}$
(D) $14$ \hfill \textbf{((C))}
\item The position vector of a point in which a line through the origin and perpendicular to the plane $2x - y - z = 4$ meets the plane $\vec{r} \cdot (3\vec{i} - 5\vec{j} + 2\vec{k}) = 6$, is\\
(A) $(1,-1,-1)$\\
(B) $(-1,-1,2)$\\
(C) $(4,2,2)$\\
(D) $(\frac{4}{3}, \frac{-2}{3}, \frac{-2}{3})$ \hfill \textbf{((D))}
\item If $\vec{a}, \vec{b}$, and $\vec{c}$ are three unit vectors equally inclined to each other at an angle $\alpha$, then the angle between $\vec{a}$ and the plane of $\vec{b}$ and $\vec{c}$ is\\
(A) $\theta = \cos^{-1} (\frac{\cos \alpha}{\cos \frac{\alpha}{2}})$\\
(B) $\theta = \sin^{-1} (\frac{\cos \alpha}{\cos \frac{\alpha}{2}})$\\
(C) $\theta = \cos^{-1} (\frac{\sin \frac{\alpha}{2}}{\sin \alpha})$\\
(D) $\theta = \sin^{-1} (\frac{\sin \frac{\alpha}{2}}{\sin \alpha})$ \hfill \textbf{((A))}
\item Q. Let $\vec{A}$, $\vec{B}$ and $\vec{C}$ be any three unit vectors. If $\vec{A} \cdot \vec{B} = \vec{A} \cdot \vec{C} = 0$ and the angle between $\vec{B}$ and $\vec{C}$ is $\frac{\pi}{4}$, then which of the following options are correct?
(A) $\vec{A} = \pm 2(\vec{B} \times \vec{C})$
(B) $\vec{A} = \pm \sqrt{2}(\vec{B} \times \vec{C})$
(C) $\vec{A} = \pm 3(\vec{B} \times \vec{C})$
(D) $\vec{A} = \pm \sqrt{3}(\vec{B} \times \vec{C})$ \hfill \textbf{((B))}
\end{enumerate}
\subsection{Hard}
\begin{enumerate}
\item Q. If the pth, qth and rth terms of a G.P. are the positive numbers a, b, c respectively, then the angle between the vectors $(\log a^2 \vec{i} + \log b^2 \vec{j} + \log c^2 \vec{k})$ and $((q-r)\vec{i} + (r-p)\vec{j} + (p-q) \vec{k})$ is
(A) $\frac{\pi}{4}$
(B) $\frac{\pi}{2}$
(C) $\sin^{-1}(\frac{1}{\sqrt{a^2+b^2+c^2}})$
(D) none of these \hfill \textbf{((B))}
\item Q. If $\vec{r}\cdot \vec{a} = 0$, $\vec{r}\cdot \vec{b} = 1$ and $[\vec{r}\vec{a}\vec{b}] = 1$, $\vec{a} \cdot \vec{b} \neq 0$, $(\vec{a} \cdot \vec{b})^2 - |\vec{a}|^2|\vec{b}|^2 = -1$. The value of $\vec{r}$ in terms of $\vec{a}$ and $\vec{b}$ is
(A) $\vec{a} \times (\vec{a} \times \vec{b}) + \vec{a} \times \vec{b}$
(B) $-\vec{a} \times (\vec{a} \times \vec{b}) + \frac{\vec{a}\times\vec{b}}{|\vec{a}\times\vec{b}|^2}$
(C) $\vec{a} \times (\vec{b}\times\vec{a}) + \frac{\vec{b}\times\vec{a}}{|\vec{b}\times\vec{a}|^2}$
(D) $\frac{\vec{a}\times(\vec{a}\times\vec{b})}{|\vec{a}\times(\vec{a}\times\vec{b})|} + \vec{a}\times\vec{b}$ \hfill \textbf{((B))}
\end{enumerate}
\section{Vector Algebra - Geometry}
\subsection{Medium}
\begin{enumerate}
\item Q. Let the position vectors of the vertices A, B, C of $\triangle ABC$ be
$\mathbf{A} = \mathbf{i} + \mathbf{j} + 2\mathbf{k}$, $\mathbf{B} = \mathbf{i} + 2\mathbf{j} + \mathbf{k}$, $\mathbf{C} = 2\mathbf{i} + \mathbf{j} + \mathbf{k}$.
Let $l_1, l_2, l_3$ be the lengths of the perpendiculars dropped from the orthocenter O onto the sides AB, BC, and CA, respectively. The value of $4(l_1 + l_2 + l_3)^2$ is:
(A) 4
(B) 5
(C) 6
(D) 7 \hfill \textbf{((C))}
\end{enumerate}
\section{Vector Algebra And 3D Geometry}
\subsection{Medium}
\begin{enumerate}
\item Q. Find the locus of the point P(x, y, z) such that the sum of its distances from A(1,0,0) and B(-1,0,0) is 10, i.e. PA + PB = 10. Which of the following equations represents this locus?
(A) $24x^2 + 25y^2 + 25z^2 = 600$
(B) $x^2 + y^2 + z^2 = 25$
(C) $4x^2 + 5y^2 + 5z^2 = 100$
(D) $16x^2 + 9y^2 + 9z^2 = 400$ \hfill \textbf{((A))}
\item Q. Find the projection of the line segment joining the points $(-1,0,3)$ and $(2,5,1)$ on the line whose direction ratios are $6:2:3$. The length of this projection is
(A) $\frac{22}{7}$
(B) $\frac{10}{7}$
(C) $\frac{11}{7}$
(D) $\frac{26}{7}$ \hfill \textbf{((A))}
\end{enumerate}
\section{Vectors}
\subsection{Medium}
\begin{enumerate}
\item Q. If $\vec{p}$ and $\vec{q}$ are two diagonals of a quadrilateral such that $|\vec{p} - \vec{q}| = \vec{p} \cdot \vec{q}$, $|\vec{p}| = 1$, $|\vec{q}| = \sqrt{2}$, then the area of quadrilateral (in sq. units) is: \hfill \textbf{((B))}
\end{enumerate}
\section{Vectors And Geometry}
\subsection{Medium}
\begin{enumerate}
\item Q. Let $\mathbf{a} = \mathbf{i} + \mathbf{j} + \mathbf{k}$, $\mathbf{b} = 2\mathbf{i} + 2\mathbf{j} + \mathbf{k}$, and $\mathbf{c} = 5\mathbf{i} + \mathbf{j} - \mathbf{k}$. Consider the locus of points $\mathbf{r}$ satisfying
$\mathbf{r} \cdot \mathbf{a} = 5$ and $|\mathbf{r} - \mathbf{b}| + |\mathbf{r} - \mathbf{c}| = 4$.
This locus is an ellipse in the plane $\mathbf{r} \cdot \mathbf{a} = 5$. The area of this ellipse, to the nearest integer, is \hfill \textbf{((A))}
\end{enumerate}
\end{document}